% Archived verbatim material removed from the rebuilt article. This file is not input by main.tex.
% Original section: sec01_introduction.tex
\section{Introduction}\label{sec:intro}

Universality statements for a single binary primitive split into different
assertions.  We use three levels.  Level \(L0\), \emph{expression universality},
asserts existence of terms.  Level \(L1\), \emph{protocol universality}, asserts
computable normal forms, compiled operations, and semantic readout identities.
Level \(L2\), \emph{certificate universality}, adds finite certificates over the
fibres of a normalizing map.

The motivating \(L0\) primitive is Odrzywo{\l}ek's operation
\[
  \eml(x,y)=e^x-\operatorname{Log}y
\]
with the constant \(1\), cited from the public arXiv v2 record
\cite{Odrzywolek2026SingleBinaryOperator}.  We use that record only as a pinned
representative-existence theorem for its displayed scientific-calculator basis,
not as a normal form or equality algorithm for EML terms.  The positive \(L1\)
protocol is internal: Zeckendorf normal forms, valuation transport through
\(\NN\), and Church iteration give compiled addition and multiplication on
\(\NN\).

The unconditional Richardson obstruction in this article is the verified-triple
replacement theorem, not a theorem about the printed strings
\((U_{47},U_{111},U_{126})\).  It applies to any supplied finite pure EML triple
whose \(\pi\), sine, and punctured absolute-value rows have been proved, and it
uses the self-contained positive-root term only through
\[
  \Abs(u)=G(M(u,u))=|u|\qquad(u\ne0).
\]
This is exactly Richardson's condition that the absolute-value symbol agree with
\(|u|\) away from zero.  Hence \Cref{thm:separation} gives, for every verified
finite replacement triple, a source-indexed pure EML Richardson fragment whose
identity-to-zero problem is undecidable on the transported partial real domains.
It follows that
there is no computable semantic normalizer into finite decidable objects that
reflects and preserves equality on those domains, and no equality-reflecting
compiler into finite decidable protocol objects.  The named specialization using
\((U_{47},U_{111})\) is a conditional branch under \((H_{\mathcal W})\) or under
separate proofs of the two printed real-domain identities, and the endpoint
square-root string \(U_{126}\) is not part of this punctured obstruction.

Some displayed theorem environments below carry additional labels because earlier
parts of the argument refer to the same result under its interface role, its
scope-cut role, or its submission-summary role.  These are cross-reference
aliases for a single numbered statement, not separate theorems or extra proof
obligations.  The first label displayed on such an environment is the canonical
source label; later labels are stable semantic aliases retained for audits and
backward references.  An alias whose name begins with a theorem, proposition,
lemma, or corollary prefix does not change the logical type of the numbered
environment to which it is attached.  The only multi-label environments used as
main-chain dependency cuts have the following canonical labels:
\[
\begin{array}{ll}
\text{source-gate classification} & \texttt{thm:source-gate-rigidity},\\
\text{core non-collapse package} & \texttt{thm:proc-ams-core-obstruction-spine},\\
\text{exact hierarchy scope closure} & \texttt{thm:main-hierarchy-exact-scope-closure},\\
\text{B36 import boundary} & \texttt{thm:B36-import-exact-scope},\\
\text{verified Richardson replacement spine} & \texttt{thm:proc-ams-relative-richardson-theorem},\\
\text{cover-relative finite-state obstruction} & \texttt{thm:finite-state-strictness}.
\end{array}
\]
All other labels on those same numbered statements are aliases for backward
references and semantic audit recognition.  Thus the proof graph is read through
the canonical labels above; the aliases do not introduce additional hypotheses,
additional conclusions, or additional theorem instances.

For submission scope, the paper proves the exact intrinsic fibre-height law and
its large-moment entropy endpoint, but deliberately does not claim a full
Zeckendorf fibre thermodynamic formalism, folded histogram law, large-deviation
theory, finite-observation escape rate, or digit-local normalization automaton.
The cover results below attach auxiliary certificate data to the fixed normal-form
sets used by the \(L1\) protocol and remain separate from the intrinsic height law.
Numbered theorem references in the source may drift during theoremization; audit
reports should identify the finite-state strictness theorem by the label
\texttt{thm:finite-state-strictness} and the value resolved in \texttt{main.aux},
not by stale prose numbering from earlier drafts.

\begin{assumption}[Conditional printed-string Richardson supplement]
\label{ass:HW}
The notation \((H_{\mathcal W})\) denotes the external analytic assertion that
the source-locked printed EML strings
\[
  W_{\pi}=U_{47},\qquad W_{\sin}=U_{111},
  \qquad W_{\sqrt{}}=U_{126}
\]
from Odrzywo{\l}ek's arXiv:2603.21852v2 ancillary file
\texttt{SupplementaryInformation.pdf}, Part~I, Sec.~1.3, Table~S2 and Part~II,
Secs.~2.1--2.5, satisfy
\[
  W_{\pi}=\pi,
  \qquad W_{\sin}(x)=\sin x\quad(x\in\RR),
  \qquad W_{\sqrt{}}(t)=\sqrt t\quad(t\in\RR_{\ge0})
\]
in the principal-branch semantics of \Cref{def:eml}.  Only the first two displayed
identities are used for the optional printed Richardson specialization in this
article; the endpoint square-root row is recorded solely as part of the external
source-lock and is not an input to the punctured Richardson obstruction.  The
public identifiers are
the versioned arXiv record \cite{Odrzywolek2026SingleBinaryOperator}, DOI
\texttt{10.48550/arXiv.2603.21852}, and the provenance snapshot
\cite{Odrzywolek2026ZenodoCode}, DOI \texttt{10.5281/zenodo.19762318}.  This
assumption is not a theorem of the article: the present manuscript proves neither
\(U_{47}=\pi\) nor \(U_{111}(x)=\sin x\) on \(\RR\), and it does not ask the reader to
infer either identity from a printed formula, parser certificate, numerical test,
or calculator-basis table.  It is used only for the optional printed
\(U_{47}\)/\(U_{111}\) specialization of the verified-triple theorem.  Unless
\((H_{\mathcal W})\) is invoked explicitly, or those two exact real-domain
identities are proved separately on the stated real domains, every Richardson
obstruction in the submission manuscript factors through the finite
verified-triple interface of \Cref{thm:separation}; in particular the main
hierarchy does not use the printed
names \(U_{47}\), \(U_{111}\), or the endpoint square-root row \(U_{126}\) as
unconditional analytic facts.  Thus there is no unconditional printed-string
Richardson instance in the article.  A citation to the Odrzywo{\l}ek \(L0\) import,
to the finite parser certificate, or to the printed catalogue is therefore never
an implicit invocation of \((H_{\mathcal W})\); the hypothesis must be named, or
the two displayed identities must be proved, before the concrete printed pair can
be used as Richardson witnesses.
\end{assumption}

\begin{proposition}[Obstruction spine]
\label{thm:obstruction-spine}
The strictness part of the paper has two primary obstructions and one subordinate
certificate addendum:
\begin{enumerate}
\item the verified-triple punctured Richardson obstruction of
  \Cref{thm:verified-replacement-richardson-instantiation,thm:replacement-triple-richardson-interface,cor:unconditional-richardson-replacement-interface,thm:richardson-witness-row-rigidity,thm:separation,thm:proc-ams-relative-richardson-theorem,thm:relative-richardson-main-spine-closure,thm:unconditional-submission-richardson-closure,cor:printed-string-free-richardson-spine};
\item the one-free-monogenic-orbit multiplication obstruction of
  \Cref{thm:level-zero-obstruction};
\item the cover-relative finite-fibre/finite-state separation of
  \Cref{thm:finite-state-strictness}, where finite-state transfer is the fixed
  matrix interface of \Cref{def:fixed-finite-state-transfer}.
\end{enumerate}
The endpoint-root string \(U_{126}\) is not part of the spine.  The first item is
an unconditional scheme over all verified finite triples: for every finite EML
triple whose \(\pi\), sine, and punctured absolute-value rows have been proved,
the punctured Richardson obstruction follows.  The printed \(U_{47}\) and
\(U_{111}\) rows are only a specialization of that scheme under
\((H_{\mathcal W})\), or under separate proofs of the same two real-domain
identities.  Thus the core theorem is the replacement-triple interface, not a
theorem about the particular printed names.
\end{proposition}

\begin{proof}
The first item is exactly \Cref{thm:separation}, whose proof factors through the
punctured triple-interface theorem
\Cref{thm:replacement-triple-richardson-interface} and its Proc. AMS formulation
\Cref{thm:proc-ams-relative-richardson-theorem}.  The row-rigidity theorem
\Cref{thm:richardson-witness-row-rigidity} proves that the three semantic rows
are not bookkeeping artifacts but exactly the necessary and sufficient data for
that structural reduction.  The submission-facing closure is
\Cref{thm:unconditional-submission-richardson-closure}: those theorems use
supplied proved rows for \(\pi\) and sine, and the internally proved punctured
absolute-value row when \(\mathsf A=\Abs\).
The second item is \Cref{thm:level-zero-obstruction}, refined by the exponent
skeleton and cyclic faithful-restriction classification.  The third item is
\Cref{thm:finite-state-strictness}; it is formulated for computable finite-fibre
covers over the fixed Zeckendorf normal-form sets and uses only the fixed
finite-state transfer notion and budget of
\Cref{def:fixed-finite-state-transfer,prop:fixed-transfer-interface-budget}.
The final sentence follows from
\Cref{thm:verified-replacement-richardson-instantiation,prop:printed-witness-certificate-boundary,prop:printed-pi-sine-instantiation-criterion,thm:unconditional-submission-richardson-closure}.
\end{proof}

\begin{theorem}[Minimal source-gate classification for the hierarchy]
\label{thm:source-gate-rigidity}
\label{thm:main-chain-external-anchor-rigidity}
\label{thm:submission-source-gate-classification}
\label{thm:minimal-source-gate-classification}
The proof of \Cref{thm:proc-ams-core-obstruction-spine,thm:main-hierarchy}
uses the following three source gates, and these gates are minimal for the
corresponding headline clauses.
\begin{enumerate}
\item The full \(L0\) calculator-basis headline is equivalent, inside this
  manuscript, to the pinned arXiv:2603.21852v2 representative-existence assertion
  for the main-text Table~1 \(\mathcal B_{36}\) basis, with the source's branch
  domains.  It is not replaceable by, nor does it imply, a row-by-row in-paper
  compiler, an equality algorithm for EML terms, a later table, a canonical
  normal form, or a proof of the printed \(U_{47}\)/\(U_{111}\) identities.
\item The Richardson obstruction in the non-collapse theorem is exactly the
  verified punctured-triple interface: finite pure EML witnesses must have proved
  rows for \(\pi\), sine on \(\RR\), and \(|x|\) on
  \(\RR\setminus\{0\}\).  The concrete printed pair \((U_{47},U_{111})\) enters
  this interface if and only if \((H_{\mathcal W})\) is assumed, or the two exact
  real-domain identities \(U_{47}=\pi\) and \(U_{111}(x)=\sin x\) on \(\RR\) are
  proved separately.  The endpoint-root string \(U_{126}\) is not a row of this
  punctured interface.
\item The intrinsic \(q=\infty\) fold-height formula uses the cited ordinary
  Fibonacci interval-maximum theorem only through the positive-index statement
  \((E)\) of \Cref{thm:intrinsic-height-source-dependency-cut}.  The new in-paper
  contribution is the conversion from \((E)\) to shifted half-interval maxima and
  then to \(\Fold_m\)-fibre height; no ordinary maximizing-integer classification,
  fibre histogram, or large-deviation theorem is imported or advertised as new.
\end{enumerate}
Consequently the submission-facing hierarchy is invariant under replacing the
first or third gate by any source or self-contained proof with the same displayed
interface, and under deleting the optional printed \((U_{47},U_{111})\) branch.
After those replacements or deletions, the surviving theorem is precisely the
printed-string-free verified-triple hierarchy with the same Zeckendorf \(L1\),
one-free-orbit, intrinsic-height, and cover-relative \(L2\) obstructions.  Any
unconditional printed \((U_{47},U_{111})\) Richardson instance, self-contained
thirty-six-row \(\mathcal B_{36}\) theorem, or new proof of the ordinary interval
maximum formula is therefore a strictly stronger input and not a consequence of
\Cref{thm:main-hierarchy}.
\end{theorem}

\begin{proof}
For the first gate, \Cref{thm:eml-universality} is explicitly a citation-only
representative-existence proposition.  The exact-scope and rigidity results
\Cref{thm:B36-import-exact-scope,thm:self-contained-L0-core-cited-B36-boundary,prop:B36-public-anchor-rigidity,thm:no-hidden-self-contained-B36-theorem,thm:B36-self-contained-replacement-core}
prove both directions of the classification.  If the manuscript asserts the full
calculator-basis headline, the only labelled source for that assertion is the
pinned public \(\mathcal B_{36}\) existence theorem with its branch-domain
conventions.  Conversely, that source gives only existential representatives, so
it cannot supply an in-paper row table, compiler, normalizer, equality algorithm,
canonical form, or Richardson witness row.  Thus every use of the full
calculator-basis clause factors through that single external representative-
existence interface, while the later strictness arguments can be rerun with only
the internally proved core witnesses after the headline is removed.

For the second gate, \Cref{thm:eml-punctured-richardson-core} proves internally
only the punctured absolute-value row \(G(M(u,u))=|u|\) for \(u\ne0\).  The first
two rows are exactly the missing \(\pi\)- and sine-row hypotheses of the abstract verified-triple and
replacement-triple interface in
\Cref{thm:verified-replacement-richardson-instantiation,thm:replacement-triple-richardson-interface,cor:unconditional-richardson-replacement-interface,thm:richardson-witness-row-rigidity,thm:proc-ams-relative-richardson-theorem,thm:unconditional-submission-richardson-closure}.
The row-rigidity theorem classifies the obstruction to using any proposed finite
triple: if the triple is not verified, the structural Richardson translation
fails already at the primitive row \(\pi\), \(\sin\), or punctured \(|-|\).  Applied
to the source-locked printed pair, this says that the names \(U_{47}\) and
\(U_{111}\) contribute only syntax until the two real-domain identities are
available.  This is recorded in
\Cref{prop:printed-pi-sine-instantiation-criterion,cor:no-unconditional-printed-richardson-triple,cor:printed-richardson-specialization-conditional,thm:exact-richardson-witness-conditionality,cor:printed-string-free-richardson-spine}.
Hence deleting \((H_{\mathcal W})\) deletes exactly the named printed-pair
specialization and leaves the verified-triple Richardson obstruction used by the
main hierarchy.  Since Richardson's hypothesis on the absolute-value primitive is
punctured, the endpoint square-root string \(U_{126}\) supplies no additional row
for this obstruction.

For the third gate, \Cref{thm:source-audited-ordinary-maximum-import} identifies
the Koc\'abov\'a--Mas\'akov\'a--Pelantov\'a input and converts it only to the
positive-index formula \((E)\).  Then
\Cref{cor:citation-sensitive-half-interval-conversion,thm:intrinsic-height-source-dependency-cut,thm:ordinary-maximum-supports-intrinsic-height,prop:kmp-indexing-lock-intrinsic-height,thm:audited-shifted-half-interval-height-conversion}
prove that the subsequent height, entropy, and zero-temperature endpoint results
use only \((E)\), the internal recursion for truncated representation counts, and
the fold-fibre pairing.  Those proofs take maxima after the half-interval and
fold conversions; they do not import maximizing residues, an intrinsic histogram,
or any pressure theorem from the ordinary interval source.  Therefore replacing
the citation by any proof of exactly \((E)\) leaves the fold-height chain
unchanged, while omitting \((E)\) removes the numerical height formulas and leaves
only the unevaluated internal maximum identities.

The final assertion is now formal.  The four load-bearing interfaces enumerated
in \Cref{thm:main-hierarchy-exact-scope-closure} are the cited \(L0\)
representative-existence gate, the verified-triple Richardson gate, the internal
Zeckendorf/one-free-orbit \(L1\) gate, and the intrinsic-height plus
cover-relative \(L2\) gate.  Replacing a source gate by the same mathematical
interface does not change any implication in that list.  Removing a stronger
optional block invokes the printed-string-free Richardson spine or the load-
bearing core projection
\Cref{cor:printed-string-free-richardson-spine,thm:verified-core-hierarchy-projection},
so the theorem chain proving \Cref{thm:main-hierarchy} is unchanged.  Conversely,
each stronger conclusion named in the last sentence supplies one of the inputs
classified above as absent, and therefore cannot be derived from the displayed
interfaces alone.
\end{proof}

\begin{theorem}[Proc. AMS hierarchy non-collapse package]
\label{thm:proc-ams-core-obstruction-spine}
\label{cor:compact-hierarchy-form}
\label{cor:richardson-submission-boundary}
\label{cor:main-L2-addendum}
The submission-facing content is the following non-collapse theorem for the
three levels of \Cref{def:expr-univ,def:prot-univ,def:cert-univ}.
\begin{enumerate}
\item The cited EML representative-existence theorem gives an \(L0\) anchor for
  the displayed \(\mathcal B_{36}\) calculator basis, but \(L0\) does not supply a
  computable semantic normalizer, nor an equality-reflecting compiler into any
  finite decidable protocol object, for the verified punctured Richardson
  fragments of \Cref{thm:proc-ams-relative-richardson-theorem}.
\item Zeckendorf normal forms give an internal \(L1\) protocol for
  \((\NN,+,\times)\), but this protocol is not reducible to a same-level ambient
  product readout on one free monogenic orbit: every faithful such readout is
  either trivial or a single cyclic multiplicative ray.
\item Finite fibres give the \(L2\) certificate layer over the fixed Zeckendorf
  normal-form sets.  For the unmodified fold maps, the intrinsic finite-degree
  moments have primitive fixed finite-state transfer and a fixed-degree pressure,
  but the full bivariate moment tower has no single rational finite-state transfer
  law.  The \(q=\infty\) height is exactly
  \(M_{2s-1}=F_{s+1}\), \(M_{2s}=2F_s\).  Separately, computable finite-fibre
  covers preserving all \(L1\) readouts can have second collision moments with no
  fixed finite-state transfer formula.  The obstruction is not merely an example:
  by the spike-height rigidity criterion of
  \Cref{thm:L2-one-spike-separation-criterion}, a cover violates the exponential
  finite-state budget exactly when its cover spike height does, and the extremal
  witness is the one-spike cover.
  Thus finite-fibre certificates do not force the cross-resolution finite-state
  transfer property.
\end{enumerate}
The Richardson item is the verified-triple theorem scheme for supplied finite
pure EML witnesses whose three rows prove \(\pi\), sine on \(\RR\), and \(|x|\) on
\(\RR\setminus\{0\}\); the paper proves the punctured absolute-value row
internally, while the printed
\(U_{47}\)/\(U_{111}\) pair becomes an instance only under
\((H_{\mathcal W})\) or separate proofs of those two real-domain identities.
The endpoint-root string \(U_{126}\), a term-by-term proof of the full
\(\mathcal B_{36}\) table, and full intrinsic fibre-histogram or large-deviation
classifications for the unmodified Zeckendorf fold fibres are not part of this
theorem.  The citation-sensitive interfaces behind these exclusions are exactly
the source gates classified in \Cref{thm:source-gate-rigidity}.
\end{theorem}

\begin{proof}
The \(L0\) anchor is \Cref{thm:eml-universality}.  The first non-collapse clause
is \Cref{thm:proc-ams-relative-richardson-theorem}, equivalently
\Cref{thm:separation}: its proof is the structural translation of
\Cref{lem:richardson-translation} followed by Richardson's punctured
identity-to-zero theorem \Cref{lem:richardson-fragment}.  If a computable
semantic normalizer into finite decidable objects reflecting and preserving
source-domain equality, or an equality-reflecting compiler into finite decidable
protocol objects, existed on one of these verified fragments, equality of the
resulting finite objects would decide identity to zero in the Richardson
fragment.  This contradiction is the decidability
argument isolated in \Cref{prop:L1-compiles-L0}.  Hence representative existence
at level \(L0\) does not collapse to this semantic-normalizer or compiler form of
\(L1\).

The positive \(L1\) protocol is
\Cref{prop:zeckendorf-protocol-data,thm:zeckendorf-protocol-univ}.  The
same-level obstruction is \Cref{thm:level-zero-obstruction}.  Under the
one-free-monogenic-orbit and ambient-product readout hypotheses every readout
factors through the prime-weight exponent skeleton, and every faithful
restriction is either the trivial monoid or a cyclic multiplicative ray.  Since
\((\NN_{\ge1},\times)\) is not one such ray, the Zeckendorf multiplication
protocol cannot be compressed to that restricted level-zero readout.  The
positive construction avoids the obstruction precisely by the Church
power-operator level shift in \Cref{thm:zeckendorf-protocol-univ}.

For \(L2\), \Cref{prop:zeckendorf-cert-univ} supplies finite-fibre certificates
over the Zeckendorf fibres.  The intrinsic transfer and height statements for
the unmodified fold are
\Cref{thm:intrinsic-fold-moment-transfer,thm:primitive-intrinsic-moment-transfer,thm:intrinsic-all-degree-transfer-obstruction,thm:intrinsic-second-moment-recurrence,cor:polynomial-intrinsic-fiber-statistics,thm:ordinary-maximum-supports-intrinsic-height,thm:intrinsic-fold-fibre-height,cor:linfty-fold-entropy,thm:moment-height-cover-trichotomy}.  In particular
\Cref{thm:ordinary-maximum-supports-intrinsic-height,prop:kmp-indexing-lock-intrinsic-height,thm:intrinsic-fold-fibre-height} give the dependency firewall, the index lock, and the two displayed height formulas,
and
\Cref{thm:intrinsic-fold-moment-transfer,thm:primitive-intrinsic-moment-transfer}
give primitive finite-state transfer and pressure for each fixed finite moment
degree.  \Cref{thm:intrinsic-all-degree-transfer-obstruction} proves the
intrinsic separation from a single rational transfer law for the entire moment
tower.  The cover-relative separation is
\Cref{thm:intrinsic-cover-transfer-dichotomy,thm:finite-state-strictness}: the
one-spike computable covers preserve the normal-form readouts by
\Cref{lem:cover-preserves-protocol} but choose second-moment spikes that violate
the exponential budget satisfied by every fixed finite-state transfer formula in
\Cref{prop:fixed-transfer-interface-budget}.  The strengthened one-spike
classification \Cref{thm:L2-one-spike-separation-criterion} proves the converse
budget lower bound as well: any computable cover whose second covered moment
defeats every exponential transfer budget must carry non-exponentially bounded
spike height, so distributed bounded-height cover excess is harmless for the
finite-state growth budget.  The strengthened budget dichotomy
\Cref{thm:moment-height-cover-trichotomy}
also shows that such bad covers must have maximal cover fibres exceeding the
intrinsic height scale by more than every exponential factor.  This proves that
the certificate layer does not collapse to cross-resolution finite-state transfer
for arbitrary computable finite-fibre covers.
The intrinsic fixed-degree pressure is the Perron root formula of
\Cref{thm:primitive-intrinsic-moment-transfer,thm:zero-temperature-fold-pressure};
the all-degree nonrationality theorem and the cover-relative degree-two
obstruction are distinct conclusions.

The submission boundary for printed names is
\Cref{cor:no-unconditional-printed-richardson-triple,thm:unconditional-submission-richardson-closure}, which leaves the
replacement-triple theorem intact and withholds only the unproved named-string
\(\pi\)/sine instance and the endpoint-root specialization.  The exclusions in
the final sentence are therefore exactly the assertions not supplied by the
cited interfaces.
\end{proof}

\begin{theorem}[Verified-core relative single-primitive hierarchy]\label{thm:main-hierarchy}
With the definitions of \(L0\), \(L1\), and \(L2\) in
\Cref{sec:L0,sec:L1,sec:L2}, the hierarchy does not collapse.  The cited
Odrzywo{\l}ek EML theorem supplies only an \(L0\) representative-existence anchor,
whereas the explicit Zeckendorf code, normalizer, decoding, compiled operations,
and Church-iteration readouts give an internal computable \(L1\) protocol for
\((\NN,+,\times)\).  The strictness is witnessed by two primary
obstructions: the verified finite replacement-triple theorem transports
Richardson's identity-to-zero theorem for any finite EML
\(\pi\)/sine/punctured-absolute-value triple whose three semantic rows have been
proved, and therefore forbids equality-reflecting
computable semantic normalizers into finite decidable objects; and multiplication
on \(\NN_{\ge1}\) cannot be faithfully read from the ambient product on one free
monogenic orbit except on the trivial monoid or a single cyclic multiplicative ray.
At level \(L2\), finite-fibre certificates over the fixed Zeckendorf normal forms
exist at every resolution, but in the computable finite-fibre cover category they
do not force fixed finite-state transfer.

The Richardson clause of the theorem is unconditional only in this
verified-triple form.  This article proves the punctured absolute-value row
internally, but the exact printed pair \((U_{47},U_{111})\) instantiates the
clause only under \((H_{\mathcal W})\) or separate proofs of
\(U_{47}=\pi\) and \(U_{111}(x)=\sin x\) on the stated real principal-branch
domains; this gate is the rigidity statement of
\Cref{thm:printed-richardson-gate-rigidity}.  The intrinsic height formula has an
analogous source boundary: the ordinary positive-index interval maximum is the
cited input, and
\Cref{thm:ordinary-maximum-supports-intrinsic-height,prop:kmp-indexing-lock-intrinsic-height,thm:audited-shifted-half-interval-height-conversion}
prove only the shifted half-interval and fold-fibre consequences needed here.
The endpoint-root string \(U_{126}\), a term-by-term proof of the full
\(\mathcal B_{36}\) table, a new proof of the ordinary interval maximum, and a
full intrinsic fibre-histogram or large-deviation theory for the unmodified
Zeckendorf fold are not consequences of the theorem.
\end{theorem}

\begin{proof}
The \(L0\) assertion is exactly the pinned external anchor
\Cref{thm:eml-universality}, with its domain and dependency boundary fixed by
\Cref{thm:L0-citation-domain-support,thm:L0-external-anchor-criterion,thm:B36-import-exact-scope}.
Thus the \(L0\) anchor imports representative existence only; it imports no
equality algorithm, semantic compiler, or Richardson witness package.  The
\(L1\) positive protocol is
\Cref{prop:zeckendorf-protocol-data,thm:zeckendorf-protocol-univ}.  The primary
strictness assertions are \Cref{thm:proc-ams-core-obstruction-spine}: its
Richardson part is \Cref{thm:separation}, proved from the punctured triple
interface
\Cref{thm:replacement-triple-richardson-interface,thm:relative-richardson-main-spine-closure},
with the printed name boundary recorded in
\Cref{ass:HW,cor:no-unconditional-printed-richardson-triple,thm:unconditional-submission-richardson-closure};
its one-orbit part is the strengthened multiplication obstruction
\Cref{thm:level-zero-obstruction,thm:restricted-level-zero-protocol-dichotomy}.

The \(L2\) layer is exactly the consolidated interface
\Cref{thm:L2-consolidated-interface}: fixed-resolution Hankel--Prony certificates,
positive intrinsic fixed-degree transfer and height, and the cover-relative
finite-state obstruction are three distinct assertions over the same normal-form
sets.  The quantitative clause of \Cref{thm:moment-height-cover-trichotomy}
separates the bad cover from the exact intrinsic height scale.  In particular,
\Cref{thm:L2-cover-intrinsic-nonimplication} excludes
interpreting the cover example as an intrinsic failure or nontransfer theorem for
the original fibres, and \Cref{thm:height-squeeze-pressure-nondetermination}
identifies the zero-temperature endpoint used here as only a maximal-height
squeeze, not a full pressure, histogram, or large-deviation theorem.  The last
item is the dependency-cut theorem
\Cref{thm:load-bearing-core-dependency-cut}; it proves that replacing the full
  calculator-basis import by the self-contained core preserves the stated
  internal theorem chain and removes only the advertised optional headlines,
  including the concrete printed \(U_{47}\)/\(U_{111}\) Richardson specialization,
  which is present only under \((H_{\mathcal W})\) or separate proofs of the two
  displayed real-domain identities.
\end{proof}

\begin{proposition}[Exact scope closure of the main hierarchy]
\label{thm:main-hierarchy-exact-scope-closure}
\label{thm:missing-results-scope-discharge}
\label{thm:main-spine-scope-closure}
\label{thm:main-theorem-scope-projection}
Every strictness conclusion used by \Cref{thm:main-hierarchy} factors through one
of the following four explicit interfaces, and through no stronger hidden
interface:
\begin{enumerate}
\item the cited \(L0\) representative-existence anchor
  \Cref{thm:eml-universality,thm:B36-import-exact-scope,thm:no-hidden-self-contained-B36-theorem}, which imports the
  public \(\mathcal B_{36}\) existence sentence but no table, normal form,
  equality algorithm, semantic compiler, or Richardson witness identity;
\item the punctured Richardson interface
  \Cref{thm:separation,thm:proc-ams-relative-richardson-theorem}, whose
  witnesses prove \(\pi\), sine on \(\RR\), and \(|x|\) on
  \(\RR\setminus\{0\}\);
\item the one-free-monogenic-orbit ambient-product interface
  \Cref{thm:level-zero-obstruction,thm:restricted-level-zero-protocol-dichotomy},
  not a classification of arbitrary single-primitive protocols realizing
  \(\{+,\times\}\);
\item the \(L2\) finite-fibre interface, consisting of the positive intrinsic
  fixed-\(q\) transfer package
  \Cref{thm:audited-fixed-degree-transfer-interface,thm:intrinsic-fold-moment-transfer,thm:primitive-intrinsic-moment-transfer,thm:intrinsic-second-moment-recurrence,cor:polynomial-intrinsic-fiber-statistics},
  the intrinsic all-degree obstruction
  \Cref{thm:intrinsic-all-degree-transfer-obstruction},
  the exact intrinsic height and large-moment endpoint
  \Cref{thm:intrinsic-height-source-dependency-cut,thm:ordinary-maximum-supports-intrinsic-height,prop:kmp-indexing-lock-intrinsic-height,thm:audited-shifted-half-interval-height-conversion,thm:intrinsic-fold-fibre-height,cor:linfty-fold-entropy,thm:zero-temperature-fold-pressure,thm:diagonal-zero-temperature-fold-endpoint,thm:moment-height-cover-trichotomy},
  and the computable-cover strictness interface
  \Cref{def:computable-cover,thm:finite-state-strictness}.  The latter is
  cover-height separated from the exact intrinsic maximum and is not an intrinsic
  nonrecurrence or finite-state-failure theorem for the unmodified fold.
\end{enumerate}
Consequently the submission theorem remains valid after simultaneously deleting
the data blocks that are outside the Proc. AMS scope: the unprinted term-by-term
\(\mathcal B_{36}\) table, the optional named-string Richardson hypothesis
\((H_{\mathcal W})\), a full online or bounded-window Zeckendorf arithmetic
theory, any global rigidity classification of arbitrary single-primitive
protocols, any proposed full intrinsic fibre histogram, all-\(q\) thermodynamic
formalism, or large-deviation theory for the original \(\Fold_m\) fibres, and any
\(L3\) capacity or Chebotarev-density layer for Hankel kernels.  After those
deletions the remaining theorem chain is exactly \Cref{thm:main-hierarchy}: the
cited \(L0\) representative-existence anchor, the internally proved core EML
witnesses, the replacement-triple punctured Richardson implication, the internal
Zeckendorf \(L1\) protocol of
\Cref{prop:zeckendorf-protocol-data,thm:zeckendorf-protocol-univ}, the one-orbit
multiplication obstruction, the intrinsic \(L2\) fixed-\(q\) transfer and height
package, and the cover-relative \(L2\) finite-state obstruction.  Conversely, an
unconditional printed-string
Richardson instance for \((U_{47},U_{111})\), a self-contained thirty-six-row
\(\mathcal B_{36}\) representative proof, a full stable Zeckendorf arithmetic
development, a global protocol-rigidity theorem, a full intrinsic fold histogram
or large-deviation classification, or an \(L3\) capacity theorem would require
additional hypotheses or additional proofs not present in this article and is
not a consequence of the displayed interfaces.
\end{proposition}

\begin{proof}
The first interface is isolated in \Cref{thm:eml-universality} and made exact by
\Cref{thm:B36-import-exact-scope}: the cited input is representative existence
for the public \(\mathcal B_{36}\) basis on the source's branch conventions.  The
proof of that exactness is part of the theorem chain, not a synonym for the
import: \Cref{thm:B36-import-exact-scope} shows that the manuscript imports only
one versioned existence statement and that no term-by-term row, endpoint value,
canonical form, equality algorithm, compiler, or Richardson witness package is
available unless separately proved.
\Cref{thm:no-hidden-self-contained-B36-theorem,thm:B36-self-contained-replacement-core,thm:load-bearing-core-dependency-cut}
then show that the later Richardson, Zeckendorf, one-orbit, and finite-cover
arguments use only the internal core witnesses once the global calculator-basis
headline is removed.  Hence no later strictness proof depends on a printed
thirty-six-row table or on a self-contained proof of that whole table.

The second interface is exactly the punctured hypothesis of \Cref{thm:separation}:
finite witnesses for \(\pi\), sine on \(\RR\), and \(|x|\) on
\(\RR\setminus\{0\}\).  \Cref{thm:unconditional-submission-richardson-closure} proves that the
endpoint-root string \(U_{126}\) is not part of this interface, and
\Cref{thm:unconditional-submission-richardson-closure} proves that deleting
\((H_{\mathcal W})\) removes the printed \(U_{47},U_{111}\) instantiation but
leaves exactly the verified-triple implication for any replacement pair whose
\(\pi\) and sine identities have been proved.  In particular, the internally
proved absolute-value row \(G(M(u,u))=|u|\) for \(u\ne0\) supplies only the third
punctured Richardson row.  The first two printed rows would require analytic
proofs of the exact equalities \(U_{47}=\pi\) and
\(U_{111}(x)=\sin x\) on their real principal-branch domain; absent
\((H_{\mathcal W})\) or such proofs, the main theorem is the replacement-triple
statement and cannot be read as an unconditional printed-string instance.

The third interface is the one-free-monogenic-orbit obstruction.  By
\Cref{thm:restricted-level-zero-protocol-dichotomy,thm:main-chain-protocol-rigidity-boundary},
the negative multiplication statement assumes a free monogenic orbit and asks
that multiplication be read by the ambient product on that same orbit.  Its proof
uses the prime-weight exponent skeleton
\Cref{lem:level-zero-exponent-skeleton}; without those hypotheses the skeleton is
not available, and the paper asserts only the positive Zeckendorf protocol of
\Cref{thm:zeckendorf-protocol-univ}.  Thus no global rigidity theorem is a hidden
input to \Cref{thm:main-hierarchy}.

The fourth interface is the \(L2\) finite-fibre interface.  On the unmodified
fold, \Cref{thm:intrinsic-fold-moment-transfer} gives fixed finite-state transfer
for every fixed intrinsic moment degree,
\Cref{thm:primitive-intrinsic-moment-transfer} gives its Perron asymptotic and
fixed-degree pressure, and
\Cref{thm:intrinsic-all-degree-transfer-obstruction} rules out a single rational
transfer law for the full moment tower.
\Cref{thm:intrinsic-second-moment-recurrence} gives the displayed sharp recurrence
for degree two, and \Cref{cor:polynomial-intrinsic-fiber-statistics} records the
resulting polynomial-statistic C-finiteness.  The same intrinsic interface also
contains the exact fibre-height theorem and large-moment endpoint.  The ordinary
Fibonacci maximum input for that height package is only the source-audited
support isolated in
\Cref{thm:intrinsic-height-source-dependency-cut,thm:ordinary-maximum-supports-intrinsic-height,prop:kmp-indexing-lock-intrinsic-height}; the intrinsic fold statement itself is
\Cref{thm:intrinsic-fold-fibre-height,cor:linfty-fold-entropy,thm:zero-temperature-fold-pressure,thm:diagonal-zero-temperature-fold-endpoint}.
Separately, \Cref{thm:finite-state-strictness} constructs covers over the fixed
normal-form sets and proves failure of the fixed finite-state transfer budget for
a cover moment.  \Cref{thm:moment-height-cover-trichotomy,thm:intrinsic-cover-transfer-dichotomy,thm:L2-cover-relative-closure,thm:L2-cover-intrinsic-nonimplication}
state the formal boundary: intrinsic fixed-\(q\) transfer and the exact intrinsic
height law coexist with cover moments that can have either finite-state or
non-finite-state behaviour, so the cover obstruction is auxiliary certificate data
and not an intrinsic failure theorem for the original maps \(\Fold_m\).  The
large-moment endpoint result used in this paper is the height/endpoint squeeze of
\Cref{thm:zero-temperature-fold-pressure,thm:diagonal-zero-temperature-fold-endpoint};
the all-degree theorem settles rational uniform transfer but does not classify
all fibre sizes or deviations away from maximal fibres.  Thus
deleting full histogram, all-\(q\) thermodynamic-formalism, or large-deviation
claims for \(\Fold_m\) changes no proved assertion in the hierarchy.  The same is
true of deleting online Zeckendorf arithmetic algorithms, global protocol
rigidity, \(L3\) capacity-density material, or the appendix-only finite Hankel
star-kernel compression, since none of these appears in the four load-bearing
interfaces above.  Combining the four paragraphs gives the claimed exact scope
closure and the announced nonconsequences.
\end{proof}

\begin{proposition}[Verified-core projection of the hierarchy]
\label{thm:verified-core-hierarchy-projection}
If one removes from the submission all optional data not proved internally or supplied as
explicit hypotheses--the \((H_{\mathcal W})\)-conditional printed
\((U_{47},U_{111})\) Richardson specialization,
the term-by-term \(\mathcal B_{36}\) table, endpoint-root claims, full stable
Zeckendorf arithmetic algorithms, global protocol rigidity, full intrinsic
fold-histogram or large-deviation assertions, and any \(L3\) capacity layer--the
remaining theorem chain still proves \Cref{thm:main-hierarchy}.  More precisely, the surviving core
is the following directed chain of load-bearing assertions:
\[
\begin{gathered}
\text{external }L0\text{ representative existence}
\longrightarrow
  \text{unconditional printed-string-free verified-triple Richardson obstruction},\\
\text{internal Zeckendorf }L1\text{ protocol}
\longrightarrow
\text{one-free-orbit multiplication obstruction},\\
\text{intrinsic fixed-}q\text{ transfer and height}
\longrightarrow
\text{cover-relative finite-state obstruction}.
\end{gathered}
\]
Every arrow in this chain is a theorem or proposition cited in
\Cref{thm:main-hierarchy-exact-scope-closure}.  Conversely, none of the four
deleted optional data blocks can be reconstructed from the surviving core: the
printed \((U_{47},U_{111})\) branch requires the two real-domain witness
identities, the full \(\mathcal B_{36}\) table requires term-by-term external or
internal verification, endpoint-root claims are ruled out by the punctured
Richardson interface, and full fold-histogram or large-deviation conclusions
require data beyond the intrinsic height and fixed-degree transfer package.
\end{proposition}

\begin{proof}
The first displayed arrow uses only the cited representative-existence anchor
\Cref{thm:eml-universality,thm:B36-import-exact-scope} and the verified-triple
Richardson theorem \Cref{thm:separation,thm:proc-ams-relative-richardson-theorem}.
By
\Cref{thm:unconditional-submission-richardson-closure,cor:printed-string-free-richardson-spine},
that Richardson theorem is formulated for arbitrary finite triples whose three
semantic rows have been proved, and deleting \((H_{\mathcal W})\) removes only the
named printed-pair specialization.  The second arrow is the positive protocol
\Cref{prop:zeckendorf-protocol-data,thm:zeckendorf-protocol-univ} together with
the restricted obstruction
\Cref{thm:level-zero-obstruction,thm:restricted-level-zero-protocol-dichotomy};
the hypotheses include the one-free-monogenic-orbit and ambient-product readout
conditions, so no global single-primitive rigidity theorem is being used.  The
third arrow is the consolidated \(L2\) interface
\Cref{thm:L2-consolidated-interface}: intrinsic fixed-degree transfer and the
height package are positive facts for the original folds, while
\Cref{thm:finite-state-strictness,thm:L2-cover-intrinsic-nonimplication} give only
the cover-relative finite-state obstruction.

These are exactly the four interfaces listed and proved sufficient in
\Cref{thm:main-hierarchy-exact-scope-closure}, with the external gates classified
in \Cref{thm:source-gate-rigidity}; hence the surviving core proves the
main hierarchy theorem.  The converse assertions are the boundary clauses of the
same proposition, together with
\Cref{thm:richardson-witness-row-rigidity,cor:no-unconditional-printed-richardson-triple}
for the printed Richardson pair,
\Cref{thm:B36-import-exact-scope,thm:no-hidden-self-contained-B36-theorem} for the
calculator-basis table, \Cref{thm:separation} for the punctured rather than
endpoint root interface, and
\Cref{thm:moment-height-cover-trichotomy,thm:L2-cover-intrinsic-nonimplication}
for the finite-fibre boundary.  Thus the optional blocks are not hidden premises
or hidden consequences of the load-bearing theorem chain.

\end{proof}

\begin{theorem}[Verified core and optional-source decomposition]
\label{thm:verified-core-printed-specialization-decomposition}
\label{thm:printed-specialization-biconditional}
\label{thm:optional-source-decomposition}
Let \(\mathcal C\) be the theorem chain obtained from the cited \(L0\)
representative-existence anchor, the internally proved punctured absolute-value
row, the verified finite EML replacement-triple Richardson interface, the
Zeckendorf \(L1\) protocol, the one-free-monogenic-orbit multiplication
obstruction, and the intrinsic/cover-relative \(L2\) interface.  Let
\((H_U)\) denote the conjunction of the two exact printed real-domain identities
\[
  U_{47}=\pi,
  \qquad
  U_{111}(x)=\sin x\quad(x\in\RR)
\]
for the printed pure EML strings on their stated principal-branch domains.  Let
\((E_{\mathrm{KMP}})\) be the positive-index ordinary interval maximum statement
isolated in \Cref{prop:kmp-indexing-lock-intrinsic-height}.  Then the
submission theorem decomposes as follows.
\begin{enumerate}
\item The chain \(\mathcal C\), together with \((E_{\mathrm{KMP}})\) for the
  intrinsic height clause, proves \Cref{thm:main-hierarchy} with its Richardson
  obstruction quantified over verified finite triples and without using the
  printed names \(U_{47}\) and \(U_{111}\) as analytic identities.
\item Relative to the structural Richardson translation used in this paper and to
  the internally proved punctured row \(\Abs(u)=|u|\) for \(u\ne0\), the printed
  \((U_{47},U_{111})\) Richardson branch is equivalent to \((H_U)\): assuming
  \((H_U)\), or \((H_{\mathcal W})\) which contains it, gives the printed branch,
  and any source-locked printed branch for these two names yields \((H_U)\).
\item The numerical intrinsic height law is equivalent, over the internal
  fold-aware conversion lemmas of this paper, to the imported statement
  \((E_{\mathrm{KMP}})\).  Replacing \((E_{\mathrm{KMP}})\) by any cited source or
  self-contained proof with the same positive-index formula leaves every
  \(\Fold_m\)-height, entropy, large-moment endpoint, and cover-height budget
  conclusion unchanged; deleting it deletes only the numerical ordinary-maximum
  input, not the Richardson, \(L1\), one-orbit, or finite-cover obstruction
  chain.
\end{enumerate}
Consequently deleting \((H_U)\) deletes exactly the printed-pair specialization
and leaves the verified-core hierarchy unchanged, while changing the source of
\((E_{\mathrm{KMP}})\) changes only the external support for the numerical height
law and not the internal fold-aware conversion.
\end{theorem}

\begin{proof}
The core chain \(\mathcal C\) uses the \(L0\) citation only through
\Cref{thm:eml-universality,thm:B36-import-exact-scope}: a representative-existence
anchor and not a table of proved row identities.  Its Richardson step is
\Cref{thm:separation,thm:proc-ams-relative-richardson-theorem}, whose hypotheses
are a finite pure EML triple with proved semantic rows for \(\pi\), sine on
\(\RR\), and \(|x|\) on \(\RR\setminus\{0\}\).  The absolute-value row is supplied
internally by \Cref{thm:eml-punctured-richardson-core}, namely
\(G(M(u,u))=|u|\) for \(u\ne0\).  The proof of the Richardson obstruction then
proceeds by the effective translation of \Cref{lem:richardson-translation} and
Richardson undecidability, and nowhere requires the syntactic names \(U_{47}\) or
\(U_{111}\).  The remaining arrows of \(\mathcal C\) are independent of those
names: the positive protocol is
\Cref{prop:zeckendorf-protocol-data,thm:zeckendorf-protocol-univ}, the restricted
multiplication obstruction is
\Cref{thm:level-zero-obstruction,thm:restricted-level-zero-protocol-dichotomy}, and
the \(L2\) assertion is
\Cref{thm:L2-consolidated-interface,thm:finite-state-strictness}.  The numerical
height part of the same \(L2\) interface uses the cited ordinary maximum only
through \((E_{\mathrm{KMP}})\), by
\Cref{prop:kmp-indexing-lock-intrinsic-height,thm:intrinsic-height-source-dependency-cut}.
Therefore the cited arrows, with that single external ordinary-maximum input for
the height formula, are exactly the proof of \Cref{thm:main-hierarchy} in
verified-core form.

If \((H_U)\) is available, then with \(\mathsf P=U_{47}\),
\(\mathsf S=U_{111}\), and \(\mathsf A=\Abs\), the three rows required by
\Cref{thm:separation} hold.  Thus the verified-triple Richardson theorem
specializes to the printed pair.  Since \((H_{\mathcal W})\) includes these two
printed identities, it gives the same specialization.

Conversely suppose that these printed names supply a source-locked Richardson
branch through the same structural translation used by the paper.  By
\Cref{thm:richardson-witness-row-rigidity}, value preservation for that branch is
possible only when its \(\pi\)-row and sine row are correct on the required real
domains.  With the printed substitutions \(\mathsf P=U_{47}\) and
\(\mathsf S=U_{111}\), those two necessary rows are exactly \((H_U)\).  Hence the
printed branch is not created by finite syntax, by the \(\mathcal B_{36}\)
representative-existence citation, or by the internally proved punctured
absolute-value row; it is created precisely by the two printed real-domain
identities.  Removing \((H_U)\) therefore removes precisely this substitution
step, while
\Cref{cor:printed-string-free-richardson-spine,thm:unconditional-submission-richardson-closure}
leave every quantified verified-triple consequence of \(\mathcal C\) in place.

It remains to prove the height-source clause.  By
\Cref{prop:kmp-indexing-lock-intrinsic-height}, \((E_{\mathrm{KMP}})\) is precisely
the positive-index formula
\(C_k=\max\{R(N):F_k\le N<F_{k+1}\}\) imported from the cited ordinary theorem.
\Cref{thm:intrinsic-height-source-dependency-cut} proves that the later chain
\[
  (E_{\mathrm{KMP}})
  \Longrightarrow \Cref{thm:half-interval-fibonacci-representation-max}
  \Longrightarrow \Cref{thm:intrinsic-fold-fibre-height}
  \Longrightarrow
  \Cref{cor:linfty-fold-entropy,thm:zero-temperature-fold-pressure,thm:moment-height-cover-trichotomy}
\]
uses only this displayed ordinary maximum formula and internal conversions:
the recursion for truncated representation counts, the shifted half-interval
argument, and the fold-fibre pairing.  Hence any replacement source proving the
same positive-index formula supplies the same numerical height and endpoint
clauses.  Conversely, without such an ordinary-maximum input the manuscript still
has the internal pairing and conversion implications, but it has no proved values
for the ordinary maxima \(C_k\) and therefore no proved numerical values
\(M_{2s-1}=F_{s+1}\), \(M_{2s}=2F_s\).  Since none of the Richardson, \(L1\),
one-orbit, or cover-relative proofs cites \((E_{\mathrm{KMP}})\), deleting that
input affects only the numerical height branch just identified.
\end{proof}

The verified-core/optional-source decomposition of
\Cref{thm:verified-core-printed-specialization-decomposition} is a theorem-level
projection of \Cref{thm:main-hierarchy}; the scoped strictness gate is now part of
\Cref{thm:proc-ams-core-obstruction-spine,thm:main-hierarchy-exact-scope-closure,thm:verified-core-hierarchy-projection}:
the Richardson obstruction is universal over verified finite EML triples, the
multiplication obstruction is restricted to one free monogenic orbit, and the
\(L2\) statement is cover-relative.
The Richardson input to the main theorem is exactly the punctured interface of
\Cref{thm:proc-ams-relative-richardson-theorem}: finite pure-EML witnesses for
\(\pi\), sine on \(\RR\), and \(|x|\) on \(\RR\setminus\{0\}\).  Under this
interface, identity to zero for Richardson expressions reduces effectively to
equality with the translated zero EML term on the same source domain, so a
computable semantic normalizer that reflects and preserves equality on those
source-indexed partial real functions would decide Richardson's undecidable
problem.  The endpoint-root string \(U_{126}\) is not part of this input.  This
is \Cref{thm:proc-ams-relative-richardson-theorem} in interface form: the proof
there constructs the translation by structural recursion and applies
Richardson's undecidability theorem recorded in \Cref{lem:richardson-fragment}.

The positive parts of the hierarchy are constructive: \((\eml,1)\) gives the
imported expression representatives, Zeckendorf valuation transport gives the
compiled \(L1\) arithmetic operations, and finite fibre-size multisets give the
\(L2\) certificates, respectively by
\Cref{thm:eml-universality,prop:zeckendorf-protocol-data,thm:zeckendorf-protocol-univ,prop:zeckendorf-cert-univ}.

% End original section: sec01_introduction.tex
% Original section: sec02_expression_universality.tex
\section{\texorpdfstring{\(L0\)}{L0}: expression universality}\label{sec:L0}

This section fixes the single primitive and the precise boundary of the imported
expression theorem.  Throughout, \(\operatorname{Log}\) is the principal complex
logarithm on
\[
  \Omega_{\mathrm{Log}}:=\CC\setminus(-\infty,0],
  \qquad -\pi<\operatorname{Im}\operatorname{Log}z<\pi,
  \qquad e^{\operatorname{Log}z}=z .
\]

\begin{definition}[EML operator]\label{def:eml}
The Exp-Minus-Log operation is the partial binary operation
\[
  \eml\colon\CC\times\Omega_{\mathrm{Log}}\to\CC,
  \qquad \eml(x,y)=e^x-\operatorname{Log}y .
\]
An EML term is interpreted only at those inputs for which every second argument
of every occurrence of \(\eml\) lies in \(\Omega_{\mathrm{Log}}\).  In particular,
nonzero points on the chosen branch cut are outside the complex domain unless a
real-domain calculation separately keeps the relevant second arguments in the
positive real axis.
\end{definition}

\begin{lemma}[Real-domain safety for the internal EML core]
\label{lem:internal-core-branch-domain-safety}
Every occurrence of \(\operatorname{Log}\) in the internally proved real EML core
of \Cref{lem:eml-printed-core,lem:eml-real-witness-library} is evaluated at a
positive real number on the stated real domain of that macro.  Consequently the
internal core identities are identities in the exact partial semantics of
\Cref{def:eml}; they do not use any value of the principal logarithm on the
negative real branch cut.
\end{lemma}

\begin{proof}
For \(E(u)=\eml(u,1)\), the only logarithm input is \(1>0\).  For
\(L(u)=\eml(1,E(\eml(1,u)))\), the domain condition is \(u>0\), so the inner
second input \(u\) is positive and the outer second input
\(E(\eml(1,u))=e^{e-\log u}=e^e/u\) is positive.  In
\(D(a,b)=\eml(L(a),E(b))\), the displayed domain has \(a>0\), and the second
input of the outer \(\eml\) is \(E(b)=e^b>0\).  The macros
\(P,N,A,S\) in \Cref{lem:eml-printed-core} are compositions of these three
forms; their stated domains include exactly the positivity hypotheses just
listed, and \(P(u)=e^u-u\ge1>0\) supplies the additional positive input used by
\(N\) and \(A\).

The witnesses in \Cref{lem:eml-real-witness-library} are then built from
\(E,L,D,P,N,A,S\) together with the auxiliary positive inputs explicitly stated
there: \(E(v)=e^v>0\), positive rational constants, products of positive terms in
the positive-domain multiplier, nonzero denominators only through positive
squares, and the final square-root argument restricted to \((0,\infty)\).  Hence
every logarithm input appearing in those calculations is a positive real number,
which lies in \(\Omega_{\mathrm{Log}}\).  The internal core therefore respects the branch
domain of \Cref{def:eml} at every step.
\end{proof}

\begin{proposition}[External public EML representation anchor]
\label{thm:eml-universality}
Odrzywo{\l}ek's public arXiv v2 record
\cite{Odrzywolek2026SingleBinaryOperator} asserts that every primitive in the
scientific-calculator basis \(\mathcal B_{36}\) displayed there has a finite
representative over \((\eml,1)\) on its stated principal-branch domain.  This
proposition records the citation role of that source: it is a citation-only
representative-existence theorem, pinned to that exact public version and to the
source's own branch-domain conventions.  The
version and domain details are external source data rather than internally proved
data of this article, and no internal proof of the thirty-six-row representation
theorem is claimed here.  Its logical content inside this paper is exactly the
following finite-existence assertion:
for each primitive named in the cited \(\mathcal B_{36}\) table there exists at
least one finite EML representative with the source's branch and domain
conventions.  The thirty-six representatives and their domain checks are not
printed or verified inside this paper.  The anchor is used here only as an
\(L0\) representative-existence theorem; it supplies no normal form, no equality
algorithm for EML terms, no semantic compiler or normalizer, no procedure taking
source primitives to verified EML terms, and no Richardson witness identity
except through the separate verified-triple interface of \Cref{thm:separation}.
Thus every later use of the citation must factor through
\Cref{thm:B36-import-exact-scope}; in particular, replacing the public anchor by
an actual compiler, an equality procedure, or a verified Richardson witness row
would be a strictly stronger external hypothesis not made in this article.
Equivalently, \Cref{thm:eml-universality} is a named imported tool, not a new
universality theorem proved in this paper.
This citation-only anchor is also not a proof of any particular source-locked
printed identity such as \(U_{47}=\pi\) or \(U_{111}(x)=\sin x\); those exact
named strings enter the Richardson part only through the separate conditional
interface \((H_{\mathcal W})\), or through an independent proof of those same two
real-domain identities.
\end{proposition}

\begin{proof}
This is a citation step, not a term-by-term derivation.  The public anchor is the
versioned record arXiv:2603.21852v2, DOI
\texttt{10.48550/arXiv.2603.21852}; the arXiv metadata identifies v2 as updated
2026-04-04, with primary class \texttt{cs.SC}, and records the ancillary file
\texttt{SupplementaryInformation.pdf}.  The supplementary text identifies the
target class as the thirty-six primitives of main-text Table~1 and states its
``Completeness for the Table-1 class'' theorem in Part~II as support for that
representative-existence claim.  Thus every occurrence of the citation in this
manuscript means that fixed v2 source object, its main-text Table~1
\(\mathcal B_{36}\) list, and the ancillary support file cited by that public
record, not a later mutable arXiv version, a repository checkout, a
generated-code artifact, or a regenerated table.  The principal-branch domains are likewise imported only
as stated in that source; this manuscript does not rederive or enlarge them.  A
representative-existence statement has only
existential force: from
``there exists a representative'' one may infer neither a canonical
representative nor an algorithm choosing one, and one may not infer any equality
decision procedure among independently written EML terms.  The ancillary file
named in \Cref{ass:HW} is used only for the optional printed-string hypothesis;
the main theorem chain treats the full \(\mathcal B_{36}\) result as the single
external existence assertion stated here.  No later proof may appeal to an
unprinted row of the \(\mathcal B_{36}\) table; all self-contained EML
calculations used later are the core calculations isolated in
\Cref{prop:internal-core-eml-witnesses}.  In particular the Richardson
obstruction can use this proposition only to name an \(L0\) background class; its
actual reduction must pass through a triple whose \(\pi\), sine, and punctured
absolute-value rows have been proved as hypotheses or internally verified.
The final sentence is a formal consequence of the preceding exclusions: a cited
existence theorem for a finite source table has no in-paper derivation unless the
individual representatives and their branch-domain checks are printed and proved
here, while this article prints only the smaller core package named in
\Cref{prop:internal-core-eml-witnesses}.
\end{proof}


\begin{theorem}[Exact scope of the \texorpdfstring{\(L0\)}{L0} import]
\label{thm:L0-citation-domain-support}
\label{thm:L0-external-anchor-criterion}
\label{thm:B36-import-exact-scope}
\label{thm:B36-import-normal-form-separation}
\label{cor:external-compiler-isolation}
\label{prop:B36-principal-branch-domain-source}
\label{thm:L0-import-replacement-invariance}
\label{thm:B36-citation-stability-and-domain-cut}
The imported Odrzywo{\l}ek input used in this paper is exactly the pinned
arXiv:2603.21852v2 claim that the main-text Table~1 scientific-calculator basis
\(\mathcal B_{36}\) has finite representatives over \((\eml,1)\) on exactly the
source's stated principal-branch domains, with the versioned ancillary Part~II
serving only as external support cited by that public record.  It is an
external representative-existence hypothesis, not a row-by-row verification
performed in this article, and no statement below converts it into an internal
term-by-term proof.
The assertion is not internalized as a hidden table: neither this theorem nor any
later argument may choose or inspect thirty-six in-paper representatives unless
they are printed and proved separately in the manuscript.
The ancillary Part~I Table~S2 strings identified in \Cref{ass:HW} are separate:
they enter this paper only through the optional printed-string hypothesis
\((H_{\mathcal W})\).  This pinning is part of the mathematical hypothesis:
replacing the citation by an unversioned web page, by another branch convention,
or by an endpoint-extended semantics would be a different assumption.  It
supplies representative existence only: it is not a term-by-term
self-contained proof in this article, not a branch-free or endpoint theorem, not
a canonical-form theorem, not an equality algorithm for EML terms, not a
semantic compiler, and not an endpoint-root Richardson package for \(\pi\), sine
on \(\RR\), and nonnegative square root on \(\RR_{\ge0}\).  Moreover, the later
Richardson, Zeckendorf \(L1\), one-orbit, and finite-cover \(L2\) arguments
remain valid if the full \(\mathcal B_{36}\) existence sentence is replaced by
the self-contained core witness package of
\Cref{prop:internal-core-eml-witnesses}; under that replacement the paper simply
no longer asserts \(L0\) expression universality for the whole calculator basis.
The multiple labels on this theorem are dependency-cut aliases for these exact
source, domain, normal-form, and replacement boundaries; they are not separate
imports from the cited table.
\end{theorem}

\begin{proof}
The citation \Cref{thm:eml-universality} imports one versioned public statement:
finite representatives for the main-text Table~1 \(\mathcal B_{36}\) basis, with
the principal logarithm and the domains fixed in that source.  The bibliographic
record pins the canonical key \texttt{Odrzywolek2026SingleBinaryOperator}, the
arXiv v2 URL, DOI, ancillary file name, and the separate Zenodo v1.1 provenance
snapshot \url{10.5281/zenodo.19762318} with archive file
\path{VA00/SymbolicRegressionPackage-v1.1.zip}.  Hence the manuscript has a
stable object to cite, but the object cited is still only the stated source
assertion and the external support materials cited by that source.
The pinning has two parts: the main-text Table~1 list is the sole \(\mathcal B_{36}\)
object being cited, and the ancillary Part~II material is not rechecked here as an
in-paper table proof.  The manuscript does not import any additional row,
generated code artifact, mutable repository state, regenerated table, or
term-by-term proof of the thirty-six representatives.  In
particular, the imported \(\mathcal B_{36}\) existence statement is kept disjoint
from the internally verified core witnesses of
\Cref{prop:internal-core-eml-witnesses}.
Every stronger item listed in the theorem would require extra data
absent from such an existence statement.  A branch-free assertion would compare
logarithm branches; an endpoint assertion would check boundary values; a
canonical form or equality algorithm would prove uniqueness or decidability for
EML terms; and a Richardson triple would prove the three displayed real-domain
identities.  None of these follows from mere representative existence, and none
is asserted by the citation step.

It remains to check that the full calculator-basis import is not hidden in the
strictness proofs.  The Richardson translation uses the core arithmetic library,
\(\exp\), division on its nonzero-denominator domain, positive square root, and a
internally repaired punctured witness package; these dependencies are recorded in
\Cref{lem:richardson-translation-bookkeeping,thm:separation}.  The Zeckendorf
protocol is proved from normal forms and value transport in
\Cref{thm:zeckendorf-protocol-univ}.  The one-orbit obstruction is the exponent
skeleton of \Cref{thm:level-zero-obstruction}.  The finite-cover statement uses
finite fibre data and the one-spike cover construction of
\Cref{thm:finite-state-strictness}.  Thus replacing the global calculator-basis
existence assertion by the internal core package preserves all strictness
arguments and removes only the headline claim that every operation in
\(\mathcal B_{36}\) has a cited representative.
\end{proof}

\begin{proposition}[Internal core EML witnesses]
\label{prop:internal-core-eml-witnesses}
In the principal-branch EML semantics, pure EML terms over \((\eml,1)\) represent
the constant \(1\), the identity input, \(\exp\), positive logarithm, negation,
addition, subtraction, multiplication, division by a nonzero real, rational
constants, \(\log2\), and square root on the positive domain \((0,\infty)\).
No value at \(0\) is asserted by this internal core package.
\end{proposition}

\begin{proof}
The explicit macros and real-domain calculations are printed in
\Cref{lem:eml-printed-core,lem:eml-real-witness-library}.  Those proofs use only
the definition of \(\eml\), positivity of the relevant inputs, and elementary
real algebra.
\end{proof}

\begin{theorem}[Self-contained \texorpdfstring{\(L0\)}{L0} core versus cited \texorpdfstring{\(\mathcal B_{36}\)}{B36} background]
\label{thm:self-contained-L0-core-cited-B36-boundary}
The self-contained \(L0\) content proved in this paper is exactly the finite core
witness package of \Cref{prop:internal-core-eml-witnesses}, together with the
calculated real-domain identities in
\Cref{lem:eml-printed-core,lem:eml-real-witness-library}.  The stronger statement
that every operation in Odrzywo{\l}ek's displayed calculator basis
\(\mathcal B_{36}\) has a finite representative over \((\eml,1)\) is used only as
the cited background theorem \Cref{thm:eml-universality}.  In particular, no
argument in the main hierarchy may invoke the \(\mathcal B_{36}\) import as a
self-contained term-by-term proof, as a normal-form algorithm, or as a verified
Richardson triple for \(\pi\), sine on \(\RR\), and square root on
\(\RR_{\ge0}\).
\end{theorem}

\begin{proof}
The first sentence is precisely the content of
\Cref{prop:internal-core-eml-witnesses}: its proof points to the printed macros
and real-domain calculations in
\Cref{lem:eml-printed-core,lem:eml-real-witness-library}, and no other
term-by-term EML representative calculation is performed in this section.  The
full calculator-basis assertion is stated separately in \Cref{thm:eml-universality},
whose proof imports the versioned public record rather than rederiving its
representatives.  \Cref{thm:B36-import-exact-scope} then proves that this import
has only representative-existence scope, and
\Cref{thm:L0-import-replacement-invariance} proves that replacing it by any
external theorem with the same existence content leaves the later \(L1\), \(L2\),
one-orbit, and punctured Richardson arguments unchanged.  The forbidden
uses listed in the last sentence are exactly the additional data excluded by
\Cref{thm:B36-import-exact-scope,cor:external-compiler-isolation,prop:B36-principal-branch-domain-source}:
term-by-term self-contained verification, canonical or decidable normal forms,
and the three real-domain identities required for the Richardson interface are
not consequences of the imported \(L0\) statement inside this paper.  Hence the
boundary is exact.
\end{proof}

\begin{proposition}[Public-anchor rigidity of the \texorpdfstring{\(\mathcal B_{36}\)}{B36} claim]
\label{prop:B36-public-anchor-rigidity}
For the full calculator-basis headline, the cited arXiv v2 object is an
irreducible external anchor in this article: every assertion here that the whole
displayed basis \(\mathcal B_{36}\) has representatives over \((\eml,1)\) is
logically equivalent, within the manuscript, to invoking
\Cref{thm:eml-universality} together with the source and domain cuts of
\Cref{thm:B36-import-exact-scope}.  The internal EML calculations of
\Cref{prop:internal-core-eml-witnesses} are sufficient for the later Richardson,
\(L1\), and \(L2\) theorem chain, but they are not a substitute proof of the full
thirty-six-row \(\mathcal B_{36}\) representative table.  Consequently the full
\(L0\) positive clause is external support in the precise sense that changing,
withdrawing, or replacing the cited public version changes only that headline
representative-existence input; it does not change any internally proved
Richardson, Zeckendorf, one-orbit, or finite-cover theorem after the smaller
core witness package is substituted.
\end{proposition}

\begin{proof}
The proposition has two directions.  First, if the manuscript asserts the full
\(\mathcal B_{36}\) representative-existence headline, the only labelled result
from which that headline is derived is \Cref{thm:eml-universality};
\Cref{thm:B36-import-exact-scope} proves that the cited object is exactly the
pinned arXiv v2 record with its stated principal-branch domains and external
support material.  Thus the full headline factors through that external
anchor and through no unversioned table, repository state, generated artifact, or
endpoint-extended branch convention.

Conversely, within the logic of this manuscript, \Cref{thm:eml-universality} is
precisely the admitted external statement that the displayed source basis has
finite representatives over \((\eml,1)\) on those domains, so invoking it with the
cuts of \Cref{thm:B36-import-exact-scope} gives exactly the full calculator-basis
\(L0\) assertion and no stronger data.  The self-contained calculations in
\Cref{prop:internal-core-eml-witnesses} prove only the smaller package named
there.  The dependency check in
\Cref{thm:B36-self-contained-replacement-core} shows that this smaller package
preserves every later Richardson, \(L1\), one-orbit, and \(L2\) argument, while
the last sentence of \Cref{thm:B36-import-exact-scope} excludes deriving the
unprinted thirty-six-row table from those internal witnesses.  Hence the public
anchor is rigid exactly for the full \(\mathcal B_{36}\) headline and non-load
bearing for the later internal theorem chain.  If the external source data are
not admitted, the proof graph therefore has a unique deletion point: remove the
headline import \Cref{thm:eml-universality}; the proof of every later internal
theorem still factors through the smaller package already isolated in
\Cref{thm:B36-self-contained-replacement-core} and the separate Richardson triple
interface.
\end{proof}

\begin{theorem}[No hidden self-contained \texorpdfstring{\(\mathcal B_{36}\)}{B36} theorem]
\label{thm:no-hidden-self-contained-B36-theorem}
In the proof graph of this article, the full calculator-basis sentence
``every primitive of \(\mathcal B_{36}\) has a representative over
\((\eml,1)\)'' is available if and only if the external anchor
\Cref{thm:eml-universality} is admitted.  After removing that cited anchor, the
manuscript proves only the internal core witness package of
\Cref{prop:internal-core-eml-witnesses}; it does not prove, imply, or use a
term-by-term table for all thirty-six operations.  Nevertheless every later
strictness result in the main theorem chain remains valid with the smaller
internal core package substituted for the full \(\mathcal B_{36}\) headline,
provided the Richardson punctured triple is supplied through the separate
interface of \Cref{thm:separation}.
\end{theorem}

\begin{proof}
The first assertion is the dependency cut in
\Cref{thm:B36-import-exact-scope}.  That theorem identifies the full
\(\mathcal B_{36}\) input with one pinned Odrzywo{\l}ek v2 representative-existence
claim and explicitly excludes all stronger data: no printed thirty-six-row proof,
no endpoint theorem, no branch-free semantics, no canonical form, no equality
algorithm, no compiler, and no Richardson witness package are imported.  Hence,
if \Cref{thm:eml-universality} is removed, no labelled result remains whose
statement has the full \(\mathcal B_{36}\) conclusion.

What remains is precisely the self-contained core package.  The displayed EML
macros and real-domain identities in
\Cref{lem:eml-printed-core,lem:eml-real-witness-library} prove the representatives
listed in \Cref{prop:internal-core-eml-witnesses}; they mention only that finite
core list and contain no enumeration of Odrzywo{\l}ek's thirty-six source
primitives.  Therefore those internal calculations cannot be reassembled into a
self-contained proof of the full source table.

Finally, \Cref{thm:B36-self-contained-replacement-core} checks the later theorem
chain one interface at a time.  The Richardson step uses the separate punctured
triple hypotheses of \Cref{thm:separation}; the Zeckendorf protocol uses residue
normal forms and readouts; the one-orbit obstruction uses the exponent skeleton;
and the finite-cover strictness argument uses only the finite sets \(X_m\) and
their certificate fibres.  None of these arguments calls on an unprinted
calculator-basis row.  Thus the full \(\mathcal B_{36}\) statement is exactly an
external expression-universality anchor, while the internally proved core is the
self-contained load-bearing replacement for the submitted strictness chain.
\end{proof}

\begin{theorem}[Submission replacement for a term-by-term \texorpdfstring{\(\mathcal B_{36}\)}{B36} table]
\label{thm:B36-self-contained-replacement-core}
For the theorem chain proved in this article, the only self-contained EML
representative data needed from \(L0\) are the finite core witnesses of
\Cref{prop:internal-core-eml-witnesses}.  Replacing the full imported
\(\mathcal B_{36}\) sentence of \Cref{thm:eml-universality} by the following
smaller internal assertion leaves every later theorem unchanged:
\[
\text{the core witnesses of \Cref{prop:internal-core-eml-witnesses} exist and
have the stated real-domain identities.}
\]
Thus a referee demand for self-containment of the hierarchy chain is met inside
the present proof graph by shrinking the load-bearing \(L0\) library to this
internally verified core.  A demand for self-containment of the full
\(\mathcal B_{36}\) headline is not met, and is not claimed to be met, by an
unprinted term-by-term proof of all thirty-six representatives; that headline
remains precisely the external citation \Cref{thm:eml-universality}.
\end{theorem}

\begin{proof}
List the later dependencies on \(L0\).  The Richardson translation uses rational
constants, \(\log2\), arithmetic operations, \(\exp\), division on its nonzero
domain, and the positive square-root macro, together with a separately supplied
punctured triple for \(\pi\), sine, and absolute value off zero; these are exactly
the dependencies recorded in \Cref{lem:richardson-translation-bookkeeping} and
\Cref{thm:separation}.  The Zeckendorf protocol and semiring readouts in
\Cref{thm:zeckendorf-protocol-univ} are constructed from Fibonacci residues and
valuation transport, not from the calculator-basis representatives.  The
finite-fibre certificate and cover arguments use only finite sets, fibre counts,
Hankel matrices, Prony reconstruction, and the cover-size construction over
\(X_m\), as recorded in \Cref{prop:zeckendorf-cert-univ,thm:finite-state-strictness}.
Finally, the one-orbit obstruction is a statement about readouts from a free
monogenic orbit and prime exponent weights; it contains no EML representative
input.

The internal core representatives named in
\Cref{prop:internal-core-eml-witnesses} are proved by the printed calculations of
\Cref{lem:eml-printed-core,lem:eml-real-witness-library}.  Therefore substituting
that smaller assertion for the imported full \(\mathcal B_{36}\) existence
sentence preserves all hypotheses actually used by the Richardson interface,
the \(L1\) protocol, the \(L2\) cover strictness theorem, and the one-orbit
obstruction.  Conversely, \Cref{thm:B36-import-exact-scope} proves that the full
import supplies no normalizer, compiler, verified Richardson triple, or
term-by-term self-contained table inside this article.  Hence the replacement is
both sufficient for the submitted theorem chain and exact with respect to the
stronger \(\mathcal B_{36}\) self-containment claim.
\end{proof}

\begin{proposition}[No hidden \texorpdfstring{\(\mathcal B_{36}\)}{B36} compiler in the submission spine]
\label{prop:no-hidden-B36-compiler-spine}
Let \(\Gamma\) be any dependency path in the submitted proof graph beginning at
the imported \(L0\) statement \Cref{thm:eml-universality} and ending at one of
\Cref{thm:separation,thm:obstruction-spine,thm:main-hierarchy,thm:finite-state-strictness}.
Then the only information about \(\mathcal B_{36}\) transported along \(\Gamma\)
is representative existence on the source's principal-branch domains.  In
particular \(\Gamma\) never yields, as an intermediate conclusion, a computable
normal form for EML terms, an equality algorithm, a semantic compiler into
finite protocol objects, or a term-by-term in-paper proof of all thirty-six
source rows.  If \Cref{thm:eml-universality} is replaced by the internal core
package of \Cref{prop:internal-core-eml-witnesses}, the same dependency paths to
the Richardson, Zeckendorf, one-orbit, and finite-cover conclusions remain after
deleting only the global \(\mathcal B_{36}\) expression-universality headline.
\end{proposition}

\begin{proof}
The proof is a cut-elimination check using the labelled interfaces already
proved in this section.  The outgoing data from \Cref{thm:eml-universality} are
cut down by \Cref{thm:B36-import-exact-scope}: the citation imports exactly the
public representative-existence assertion for the source table and explicitly
excludes normal forms, equality algorithms, semantic compilers, endpoint
extensions, and a row-by-row in-paper derivation.  Therefore a dependency path
can acquire one of those stronger objects only if some later theorem proves it
from another hypothesis.

The later hypotheses are disjoint.  The Richardson path uses the punctured
triple identities isolated in \Cref{thm:separation} and
\Cref{thm:replacement-triple-richardson-interface}; those identities are
supplied either by verified replacement witnesses or, conditionally for the
printed pair, by \((H_{\mathcal W})\) or separate real-domain proofs.  The
Zeckendorf protocol path uses the normal-form sets and readout identities of
\Cref{sec:L1}.  The one-orbit obstruction uses only the free-monogenic exponent
skeleton, and the finite-cover obstruction uses only computable fibre-size data
over the fixed sets \(X_m\).  None of these interfaces has a premise naming the
thirty-six calculator-basis rows or a conclusion reconstructing them.  Hence no
hidden compiler or normalizer can appear along \(\Gamma\).

Finally, \Cref{thm:B36-self-contained-replacement-core} proves that the internal
core witnesses are sufficient for every later theorem except the global
calculator-basis existence headline itself.  Replacing the import by that core
therefore deletes exactly the external \(\mathcal B_{36}\) expression-universality
claim and preserves the stated strictness paths, which proves the final
assertion.
\end{proof}

\begin{theorem}[Interface rigidity for the load-bearing \texorpdfstring{\(L0\)}{L0} and Richardson inputs]
\label{thm:L0-richardson-exact-load-bearing-cut}
\label{thm:submission-interface-rigidity}
The Proc. AMS theorem chain uses the expression-universality input and the
Richardson input through disjoint interfaces, and these interfaces are rigid in
the following sense.  Its \(L0\) interface is exactly one of the following two
choices:
\begin{enumerate}
\item the cited Odrzywo{\l}ek representative-existence theorem for the displayed
  \(\mathcal B_{36}\) basis, with no term-by-term verification reproduced here;
\item the smaller self-contained core witness package of
  \Cref{prop:internal-core-eml-witnesses}, in which case the global
  \(\mathcal B_{36}\) expression-universality sentence is deleted but the later
  strictness theorems are unchanged.
\end{enumerate}
Its Richardson interface is exactly the punctured witness package of
\Cref{thm:richardson-witness-row-rigidity,thm:separation}: finite pure EML
witnesses for \(\pi\), sine on \(\RR\), and \(|x|\) on
\(\RR\setminus\{0\}\).  Consequently the cited
\(\mathcal B_{36}\) import, the finite syntax catalogue, and the endpoint-root
string \(U_{126}\) are not load-bearing substitutes for those identities.  In
particular, after deleting the optional printed-string hypothesis
\((H_{\mathcal W})\), every conclusion of \Cref{thm:main-hierarchy} that remains
in force factors through the replacement-triple Richardson interface; no
unconditional conclusion naming the printed pair \((U_{47},U_{111})\) remains
unless the two real-domain identities \(U_{47}=\pi\) and
\(U_{111}(x)=\sin x\) on \(\RR\) are explicitly assumed under
\((H_{\mathcal W})\) or separately proved.
\end{theorem}

\begin{proof}
The two possible \(L0\) inputs are precisely the alternatives isolated in
\Cref{thm:L0-external-anchor-criterion}.  The first alternative is the cited
representative-existence proposition \Cref{thm:eml-universality}; by
\Cref{thm:B36-import-exact-scope} and \Cref{thm:B36-import-normal-form-separation}
it imports only existence of representatives on the stated principal-branch
domains, not a normalizer, equality procedure, compiler, or Richardson witness.
The second alternative is the replacement of the global calculator-basis claim by
the finite internal package of \Cref{prop:internal-core-eml-witnesses};
\Cref{thm:B36-self-contained-replacement-core} proves that every later
Richardson, Zeckendorf, one-orbit, and finite-cover argument still has all of its
actual hypotheses under that replacement.

It remains to identify the Richardson input.  In
\Cref{thm:richardson-witness-row-rigidity,thm:separation} and
\Cref{thm:replacement-triple-richardson-interface} the hypotheses are not the
existence of a large calculator-basis representation theorem.  They are exactly
the finite punctured witnesses for \(\pi\), sine, and \(|x|\) off zero: row
rigidity proves these identities are necessary and sufficient for the structural
Richardson reduction.  The proof
of \Cref{thm:B36-import-exact-scope} excludes deriving endpoint or normal-form
data from the \(\mathcal B_{36}\) import alone, and
\Cref{prop:printed-pi-sine-instantiation-criterion,cor:no-unconditional-printed-richardson-triple}
shows that the finite printed syntax of \((U_{47},U_{111})\) becomes a
Richardson instance exactly when those two printed strings have their
real-domain identities proved.  The Richardson theorem is relative to any
supplied punctured triple with those identities proved, and it proves the
absolute-value row internally, so the endpoint-root string \(U_{126}\) is not
load-bearing for \Cref{thm:main-hierarchy}.

Finally, \Cref{thm:main-hierarchy-exact-scope-closure} enumerates the four
interfaces through which the main theorem factors: the pinned \(L0\) import or
its internal-core replacement, the replacement-triple Richardson theorem, the
one-orbit \(L1\) obstruction, and the cover-relative \(L2\) budget obstruction.
The deletion of \((H_{\mathcal W})\) affects only the named printed
\((U_{47},U_{111})\) branch of the second interface.  Since all remaining
arrows in that enumeration cite the replacement-triple theorem rather than the
printed names, every surviving conclusion factors through the relative
interface, as claimed.
\end{proof}

\begin{theorem}[Load-bearing core after deleting optional external data]
\label{thm:load-bearing-core-dependency-cut}
Let the \emph{core theorem package} consist of the internally proved EML core of
\Cref{prop:internal-core-eml-witnesses}, the punctured Richardson interface
of \Cref{thm:separation}, the Zeckendorf \(L1\) protocol of
\Cref{thm:zeckendorf-protocol-univ}, the one-free-monogenic-orbit obstruction of
\Cref{thm:level-zero-obstruction}, the intrinsic \(L2\) transfer and height package of
\Cref{thm:intrinsic-fold-moment-transfer,thm:intrinsic-fold-fibre-height,cor:linfty-fold-entropy,thm:zero-temperature-fold-pressure},
and the cover-relative \(L2\) finite-state obstruction of
\Cref{thm:finite-state-strictness}.  If one deletes
\begin{enumerate}
\item the unprinted term-by-term \(\mathcal B_{36}\) representative table, keeping
  the internal core witnesses and a supplied punctured Richardson triple; and
\item the optional printed-string hypothesis \((H_{\mathcal W})\), keeping the
  verified-triple Richardson theorem for all triples whose \(\pi\), sine, and
  punctured absolute-value identities have been proved,
\end{enumerate}
then every assertion in the core theorem package remains proved with the same
statements.  The only assertions removed by this deletion are the global cited
\(\mathcal B_{36}\) expression-universality headline and the named printed
\((U_{47},U_{111})\) Richardson specialization.  In particular, the cover-relative
\(L2\) strictness theorem remains a theorem about computable covers over the fixed
normal-form sets \(X_m\), while the positive intrinsic \(\Fold_m\) fixed-\(q\)
transfer and height endpoint package remains the one proved directly in
\Cref{sec:L2}; no full intrinsic histogram or large-deviation classification of
the original fibres is introduced by the deletion.
\end{theorem}

\begin{proof}
The internal EML core remains because it is proved directly by
\Cref{lem:eml-printed-core,lem:eml-real-witness-library} and summarized in
\Cref{prop:internal-core-eml-witnesses}; it does not use the full
\(\mathcal B_{36}\) table.  The Richardson theorem used in the main spine is the
triple-relative \Cref{thm:separation}: its \(\pi\) and sine rows must be supplied
with proofs, and its punctured absolute-value row can be taken from the internal
term \(\Abs\).  Deleting \((H_{\mathcal W})\) therefore deletes only the optional
claim that the particular printed names \((U_{47},U_{111})\) supply the first two
rows.  It does not delete the universal verified-triple theorem, nor any instance
whose first two rows have been proved independently.  This is exactly the
interface identified in \Cref{thm:L0-richardson-exact-load-bearing-cut} and
\Cref{thm:unconditional-submission-richardson-closure}.

The Zeckendorf protocol theorem \Cref{thm:zeckendorf-protocol-univ} is proved from
finite normal forms, Fibonacci valuation transport, and Church-iteration
identities, not from EML calculator representatives or from the printed
Richardson strings.  The one-orbit obstruction \Cref{thm:level-zero-obstruction}
is a statement about readouts from one free monogenic orbit and its prime-weight
exponent skeleton, so it has no dependency on either deleted datum.  Finally,
\Cref{thm:intrinsic-fold-moment-transfer} proves the positive intrinsic transfer
statement directly from the bounded carry automaton.  Separately,
\Cref{thm:finite-state-strictness} is proved by the computable-cover skeleton:
\Cref{thm:finite-fiber-cover-realization} turns a computable positive fibre-size
function over \(X_m\) into a cover, \Cref{lem:cover-preserves-protocol} preserves
all normal-form readouts, and \Cref{prop:finite-state-coefficient-budget} gives
the exponential budget contradicted by a superexponential one-spike choice such
as \(a_m=m!\).  None of these cover-obstruction steps uses the intrinsic
multiplicities of \(\Fold_m\); by \Cref{thm:L2-cover-intrinsic-nonimplication} a
cover-moment transfer statement or obstruction cannot be transported to the
original fold fibres without an additional equality or comparison theorem.
Therefore the stated deletion leaves
exactly the listed core package and removes exactly the two optional headlines.
\end{proof}

\begin{definition}[Expression universality]\label{def:expr-univ}
A system is \(L0\)-expression-universal for a signature if every primitive
operation in the signature has at least one term representative over the chosen
primitive operations on its stated domain.  No normal form, equality procedure,
or fibre certificate is part of the definition.
\end{definition}

\begin{remark}\label{rem:eml-variants}
All EML claims in this article use the principal-branch convention of
\Cref{def:eml}; no branch-free complex calculus is asserted.
\end{remark}

\begin{remark}\label{rem:no-canon}
Expression universality alone does not choose canonical representatives.  This
is the gap separated by the Richardson obstruction in \Cref{thm:separation}.
\end{remark}



% End original section: sec02_expression_universality.tex
% Original section: sec03_protocol_universality.tex
\section{\texorpdfstring{\(L1\)}{L1}: a Zeckendorf protocol}\label{sec:L1}

Let \(F_1=1,F_2=2,F_{k+2}=F_{k+1}+F_k\).  Let
\(X_\infty^{\mathrm{fin}}\) be the finite-support binary sequences with no
adjacent \(1\)'s, and set
\(\Val(c)=\sum_{k\ge1}c_kF_k\).

\subsection{Normal forms}\label{ssec:fold}

\begin{definition}[Finite Zeckendorf interface]\label{def:fold}
For \(m\ge1\), put
\[
  \Omega_m=\{0,1\}^{m+1},
  \qquad
  X_m=\{c\in X_\infty^{\mathrm{fin}}:0\le \Val(c)<F_{m+2}\}.
\]
The finite valuation is
\[
  \Val_m\colon X_m\to \{0,1,\ldots,F_{m+2}-1\},
  \qquad \Val_m(c)=\Val(c).
\]
For a finite binary word \(\omega\in\Omega_m\), \(\Fold_m(\omega)\in X_m\) is
the greedy Zeckendorf normal form of the residue
\(\sum_{k=1}^{m+1}\omega_kF_k\pmod {F_{m+2}}\).  The stable normalizer
\(\mathrm{NF}_\infty\) is defined analogously for finite words of arbitrary
length, without reducing modulo a finite level.
\end{definition}

\begin{equation}\label{eq:fold-residue}
  \Val_m(\Fold_m(\omega))\equiv\sum_{k=1}^{m+1}\omega_kF_k\pmod {F_{m+2}} .
\end{equation}

\begin{proposition}[Finite Zeckendorf residue interface]
\label{prop:finite-zeckendorf-residue-interface}
The map \(\Val_m\) is a computable bijection
\[
  X_m\longrightarrow \ZZ/F_{m+2}\ZZ
\]
after identifying \(\ZZ/F_{m+2}\ZZ\) with its representatives
\(0,\ldots,F_{m+2}-1\).  Hence \(|X_m|=F_{m+2}\), and
\(\Fold_m\colon\Omega_m\to X_m\) is a computable surjection.
\end{proposition}

\begin{proof}
For each \(0\le n<F_{m+2}\), \Cref{lem:greedy-zeckendorf-normalizer} gives a
unique Zeckendorf representative \(c_n\in X_\infty^{\mathrm{fin}}\) with
\(\Val(c_n)=n\).  By definition this representative belongs to \(X_m\), so
\(\Val_m\) is surjective onto the displayed representative set.  If
\(c,d\in X_m\) and \(\Val_m(c)=\Val_m(d)\), then \(\Val(c)=\Val(d)\) as ordinary
integers in the interval \([0,F_{m+2})\), and Zeckendorf uniqueness gives
\(c=d\).  The greedy algorithm computes the inverse representative
\(n\mapsto c_n\), and reduction of \(\sum_{k=1}^{m+1}\omega_kF_k\) modulo
\(F_{m+2}\), followed by that inverse representative, computes \(\Fold_m\).  To
see surjectivity of \(\Fold_m\), let \(c\in X_m\) and let
\(\Val(c)=n<F_{m+2}\).  The Zeckendorf representative of such an \(n\) uses no
summand larger than \(F_{m+1}\); otherwise its value would be at least
\(F_{m+2}\).  Thus the first \(m+1\) digits of \(c\), viewed as an element of
\(\Omega_m\), fold back to \(c\).
\end{proof}

\begin{lemma}[Greedy Zeckendorf normalizer]
\label{lem:greedy-zeckendorf-normalizer}
Every \(n\in\NN\) has a unique representative \(c_n\in
X_\infty^{\mathrm{fin}}\), and the map \(n\mapsto c_n\) is computable.
\end{lemma}

\begin{proof}
This is the classical Zeckendorf theorem \cite{Zeckendorf1972}: repeatedly
subtract the largest Fibonacci number not exceeding the remaining integer.  The
greedy choice terminates, produces no adjacent summands, and uniqueness follows
from the standard exchange inequality
\(F_{k+1}>F_k+
F_{k-2}+F_{k-4}+\cdots\).
\end{proof}

\begin{proposition}[Computable Zeckendorf arithmetic anchor]
\label{prop:zeckendorf-computable-arithmetic-anchor}
The stable Zeckendorf normal-form set \(X_\infty^{\mathrm{fin}}\) gives a
computable presentation of \(\NN\): the valuation \(\Val\) and its greedy inverse
are computable, equality is decidable, and the operations transported from
ordinary addition and multiplication are computable on normal forms.  The only
finite-level operation used below is the residue fold of \Cref{def:fold}.  No
online normalizer, finite-window arithmetic algorithm, automatic-structure
theorem, or Ostrowski-definability theorem is used as a hidden input to the
protocol theorem.
\end{proposition}

\begin{proof}
The map \(\Val\) is computed by the finite Fibonacci sum of the support, and
\Cref{lem:greedy-zeckendorf-normalizer}, the classical Zeckendorf theorem
\cite{Zeckendorf1972}, computes the inverse normal form.  Equality is decidable
by comparing finite binary strings after deleting trailing zeroes, or
equivalently by applying \(\Val\).  To compute transported addition or
multiplication of \(c,d\in X_\infty^{\mathrm{fin}}\), compute
\(a=\Val(c)\), \(b=\Val(d)\), form \(a+b\) or \(ab\) in ordinary integer
arithmetic, and apply the greedy inverse.  The finite fold is computable by
\Cref{prop:finite-zeckendorf-residue-interface}.  The broader literature gives
stronger finite-word algorithms and automatic/Ostrowski arithmetic interfaces
\cite{AhlbachEtAl2013,HieronymiTerry2018}; those results are cited as public
anchors for the numeration background, but the present argument needs only the
explicit computations just described.
\end{proof}

\begin{proposition}[Fold confluence]
\label{prop:fold-confluence}
Two finite words have the same stable fold if and only if they have the same
ordinary valuation.  At finite level \(m\), two words of \(\Omega_m\) have the
same finite fold if and only if their valuations have the same residue modulo
\(F_{m+2}\).  Hence equality of both stable and finite normal forms is decidable.
\end{proposition}

\begin{proof}
The stable assertion follows because both words fold to the unique Zeckendorf
representative of their ordinary valuation by
\Cref{lem:greedy-zeckendorf-normalizer}.  The finite assertion follows from the
bijection in \Cref{prop:finite-zeckendorf-residue-interface}.  Finite binary
words and residues are decidably comparable.
\end{proof}

\begin{definition}[Completions]\label{def:completions}
The arithmetic locus is
\begin{equation}\label{eq:X-infty}
  X_\infty^{\mathrm{fin}}
  =\{c\in\{0,1\}^{\NN}: c\text{ has finite support and no adjacent }1\text{'s}\}.
\end{equation}
The compact inverse-limit completion \(\widehat X\) is not used as an arithmetic
locus.
\end{definition}

\begin{remark}\label{rem:finite-vs-profinite}
The valuation is a bijection \(X_\infty^{\mathrm{fin}}\to\NN\), not a bijection
\(\widehat X\to\NN\).
\end{remark}

\subsection{Transported arithmetic}\label{ssec:arith}

\begin{definition}[Stable arithmetic]\label{def:stable-arith}
For \(c,d\in X_\infty^{\mathrm{fin}}\), define
\begin{align}
  c\oplus d&=\Val^{-1}(\Val(c)+\Val(d)),\label{eq:stable-add}\\
  c\otimes d&=\Val^{-1}(\Val(c)\Val(d)).\label{eq:stable-mul}
\end{align}
\end{definition}

\begin{proposition}[Finite residues are not stable arithmetic]
\label{prop:finite-residue-not-stable}
No fixed finite level \(X_m\) carries the stable arithmetic of
\((\NN,+,\times)\) through valuation.
\end{proposition}

\begin{proof}
The finite set \(X_m\) has only \(F_{m+2}\) elements, while addition or
multiplication on \(\NN\) leaves every finite interval.  Thus a fixed finite
residue level cannot be closed under the transported operations.
\end{proof}

\begin{proposition}[Finite residue-level transport]
\label{prop:finite-residue-level-transport}
For each \(m\), the bijection \(\Val_m\) transports the ring structure of
\(\ZZ/F_{m+2}\ZZ\) to \(X_m\).  Explicitly,
\[
\begin{aligned}
  c\oplus_m d&=\Val_m^{-1}(\Val_m(c)+\Val_m(d)\bmod F_{m+2}),\\
  c\otimes_m d&=\Val_m^{-1}(\Val_m(c)\Val_m(d)\bmod F_{m+2}).
\end{aligned}
\]
Then \(\Val_m\) is a finite ring isomorphism from
\((X_m,\oplus_m,\otimes_m)\) to \(\ZZ/F_{m+2}\ZZ\).  This is residue arithmetic
only and is not the stable arithmetic of \Cref{def:stable-arith}.
\end{proposition}

\begin{proof}
The displayed operations are the pullback of addition and multiplication in the
finite quotient ring through the bijection
\(\Val_m\colon X_m\to\ZZ/F_{m+2}\ZZ\) of
\Cref{prop:finite-zeckendorf-residue-interface}.  Therefore the ring axioms and
the isomorphism statement are immediate from the corresponding quotient-ring
facts.  The final sentence is \Cref{prop:finite-residue-not-stable}: a fixed
finite quotient records only residues modulo \(F_{m+2}\), while stable
arithmetic records ordinary natural-number values without quotienting.
\end{proof}

\begin{theorem}[Transported semiring]
\label{thm:semiring}
The map \(\Val\colon X_\infty^{\mathrm{fin}}\to\NN\) is an isomorphism from
\((X_\infty^{\mathrm{fin}},\oplus,\otimes)\) to \((\NN,+,\times)\).  At each
finite level, the only transported arithmetic asserted in this paper is the
residue ring interface of \Cref{prop:finite-residue-level-transport}.
\end{theorem}

\begin{proof}
This is immediate from the definitions of \(\oplus\) and \(\otimes\) and the
bijection in \Cref{lem:greedy-zeckendorf-normalizer}.  The finite-level clause
is \Cref{prop:finite-residue-level-transport}.
\end{proof}

Stable arithmetic in \Cref{thm:semiring} cannot be replaced by arithmetic on one
fixed finite residue set \(X_m\), by \Cref{prop:finite-residue-not-stable}.

\begin{lemma}[Value transport suffices]
\label{lem:value-transported-suffices}
Any computable bijection \(N\to\NN\) transports ordinary computable operations on
\(\NN\) to computable operations on \(N\).
\end{lemma}

\begin{proof}
Decode, apply the ordinary operation on \(\NN\), and encode back by the inverse
bijection.
\end{proof}

\begin{proposition}[Value-transport exact interface]
\label{prop:value-transport-exact-interface}
The Zeckendorf protocol uses only the computable bijection
\(\Val\colon X_\infty^{\mathrm{fin}}\to\NN\) and not any digit-local arithmetic
algorithm.
\end{proposition}

\begin{proof}
The operations are exactly those in \Cref{def:stable-arith}; both are obtained by
decode--compute--encode.
\end{proof}

\subsection{One primitive at two levels}\label{ssec:sole}

Let \(A\) be a nonempty set and let \(g\colon A\to A\) have infinite order.  Put
\[
  I_n(f)=f^{\circ n}\qquad(f\in\End(A)).
\]

\begin{theorem}[Two-level readout]
\label{thm:two-level}
With composition as the single semantic operation,
\begin{align}
  \Phi_0(c_n)&=g^{\circ n},\label{eq:level0}\\
  \Phi_1(c_n)&=I_n\label{eq:level1}
\end{align}
read addition and multiplication:
\begin{align}
  \Phi_0(c_{a+b})&=\Phi_0(c_a)\circ\Phi_0(c_b),\label{eq:add-from-comp}\\
  \Phi_1(c_{ab})&=\Phi_1(c_a)\circ\Phi_1(c_b).\label{eq:mul-from-comp}
\end{align}
\end{theorem}

\begin{proof}
The first identity is \(g^{a+b}=g^a\circ g^b\).  The second is
\(I_a\circ I_b(f)=(f^b)^a=f^{ab}=I_{ab}(f)\).
\end{proof}

Thus addition is read at the orbit level and multiplication at the power-operator
level; both readouts use composition as the only semantic operation.
\begin{lemma}[Level-zero exponent skeleton]
\label{lem:level-zero-exponent-skeleton}
Let \(C=\{g^n:n\in\NN\}\) be a free monogenic orbit.  Every multiplicative
readout \(\eta\colon\NN_{\ge1}\to C\) has the form
\[
  \eta(n)=g^{e(n)},
  \qquad e(ab)=e(a)+e(b),
\]
and hence
\[
  e(n)=\sum_{p}v_p(n)w(p)
\]
for uniquely determined weights \(w(p)=e(p)\in\NN\).
\end{lemma}

\begin{proof}
Since \(g
\) has infinite order, each element of \(C\) has a unique exponent.  The
homomorphism law for \(\eta\) becomes additivity of \(e\) on multiplication.
Unique prime factorization gives the displayed formula.
\end{proof}

\begin{theorem}[One-orbit multiplication obstruction and extremal skeleton]
\label{thm:level-zero-obstruction}
Fix one free monogenic orbit \(C=\{g^n:n\in\NN\}\) and require multiplication to
be read at level zero by the ambient product on \(C\).  Then every
multiplicative readout \(\eta\colon\NN_{\ge1}\to C\) is obtained from a unique
prime-weight function \(w\colon\mathcal P\to\NN\) by
\[
  \eta(n)=g^{e(n)},
  \qquad e(n)=\sum_p v_p(n)w(p),
\]
and no such readout is injective on all of \(\NN_{\ge1}\).  More sharply, if
\(H\subseteq\NN_{\ge1}\) is any multiplicative submonoid, then
\(\eta|_H\) is injective if and only if either \(H=\{1\}\), or there is an
integer \(q>1\) such that
\[
  H\subseteq\{q^k:k\in\NN\}
  \qquad\hbox{and}\qquad e(q)>0.
\]
Thus the obstruction is not a failure of computability or of normal forms: it is
the rank-one exponent skeleton forced by asking ambient product on one free
monogenic orbit to carry integer multiplication.  The positive replacement used
in the main hierarchy is the level shift to the Church power-operator readout of
\Cref{thm:two-level}.
\end{theorem}

\begin{proof}
The factorization and uniqueness of the weights are exactly
\Cref{lem:level-zero-exponent-skeleton}: the infinite order of \(g\) makes
exponents unique, the law \(\eta(ab)=\eta(a)\eta(b)\) makes \(e\) additive on
multiplication, and unique factorization gives
\(e(n)=\sum_p v_p(n)w(p)\) with \(w(p)=e(p)\).

This already excludes injectivity on all of \(\NN_{\ge1}\).  Choose distinct
primes \(p,r\).  If either \(w(p)=0\) or \(w(r)=0\), then the corresponding prime
has the same image as \(1\).  If both weights are positive, then
\[
  e(p^{w(r)})=w(p)w(r)=e(r^{w(p)}),
\]
so the two distinct prime powers have the same image.

It remains to prove the extremal statement for an arbitrary multiplicative
submonoid \(H\).  Suppose first that \(\eta|_H\) is injective.  If
\(h\in H\setminus\{1\}\) and \(e(h)=0\), then \(\eta(h)=g^0=\eta(1)\), a
contradiction.  Hence every nonunit in \(H\) has positive exponent weight.  For
\(a,b\in H\setminus\{1\}\), additivity gives
\[
  \eta(a^{e(b)})=g^{e(a)e(b)}=\eta(b^{e(a)}).
\]
Injectivity on \(H\) implies \(a^{e(b)}=b^{e(a)}\).  Comparing prime
factorizations yields
\[
  e(b)v_p(a)=e(a)v_p(b)
  \qquad\hbox{for every prime }p,
\]
so the exponent vectors of any two nonunit elements of \(H\) lie on the same
rational ray.  If \(H\ne\{1\}\), choose \(h\in H\setminus\{1\}\), let
\(\alpha=(\alpha_p)_p\) be the primitive nonnegative integer vector spanning the
rational ray through \((v_p(h))_p\), and set \(q=\prod_p p^{\alpha_p}\).  The
support is finite and \(q>1\).  Every \(x\in H\) then has exponent vector
\(k\alpha\) for some \(k\in\NN\), so \(x=q^k\).  Also \(e(q)>0\), because the
chosen element \(h=q^k\) is a nonunit of positive exponent weight.

Conversely, if \(H\subseteq\{q^k:k\in\NN\}\) and \(e(q)>0\), then
\(\eta(q^a)=\eta(q^b)\) implies \(ae(q)=be(q)\), hence \(a=b\).  Thus
\(\eta|_H\) is injective.  The final sentence is the positive identity
\(\Phi_1(c_{ab})=\Phi_1(c_a)\circ\Phi_1(c_b)\) from \Cref{thm:two-level}, which
reads multiplication after shifting from the orbit to the power-operator level.
\end{proof}

The obstruction applies only to same-level ambient-product readouts on one free
monogenic orbit.  It does not apply to the two-level readout of
\Cref{thm:two-level}, which reads multiplication by composition of power operators
in \(\End(\End(A))\), not by the ambient product on \(\{g^n\}\).  Within the
one-orbit model, \Cref{thm:level-zero-obstruction} is sharp: only the trivial
monoid or a single cyclic multiplicative ray can remain faithful.  The finer
prime-coordinate classification is split off from this short article.  No
finite-rank product-readout classification is used in this paper.

\begin{corollary}[Minimal two-level split]
\label{cor:minimal}
In the setting of \Cref{thm:two-level}, the two displayed readout levels are
jointly sufficient for the usual arithmetic semiring on \(\NN\): the map
\(n\mapsto c_n\) is injective, \(\Phi_0\) reads addition by composition in
\(\End(A)\), and \(\Phi_1\) reads multiplication by composition in
\(\End(\End(A))\).  Moreover, under the one-free-monogenic-orbit restriction in
which the same orbit \(C=\{g^{\circ n}:n\in\NN\}\) is required to carry an
injective multiplication readout by its ambient product, no level-zero-only
replacement can be injective on all of \(\NN_{\ge1}\).  Hence within that
restricted model the passage from the orbit level to the power-operator level is
the minimal extra readout needed to keep one semantic primitive, composition,
while reading both \(+\) and \(\times\).
\end{corollary}

\begin{proof}
The sufficiency assertions are exactly the identities of \Cref{thm:two-level}:
\(g^{\circ(a+b)}=g^{\circ a}\circ g^{\circ b}\) at level zero and
\(I_a\circ I_b=I_{ab}\) at the power-operator level.  The symbols \(c_n\) are
indexed by \(n\in\NN\), so \(n\mapsto c_n\) is injective as a syntactic normal
form map; when read at level zero, the infinite order of \(g\) also makes
\(n\mapsto g^{\circ n}\) injective.

It remains only to justify the stated minimality, which is restricted to the
one-orbit model named in the statement.  If multiplication were read there by a
level-zero map \(\eta\colon\NN_{\ge1}\to C\) satisfying
\(\eta(ab)=\eta(a)\eta(b)\) for the ambient product on \(C\), then unique
exponents in the free monogenic orbit would give
\(\eta(n)=g^{\circ e(n)}\) and \(e(ab)=e(a)+e(b)\).  This is the exponent skeleton
proved in \Cref{lem:level-zero-exponent-skeleton}.  The prime-power collision in
\Cref{thm:level-zero-obstruction} then shows that such an \(\eta\) cannot be
injective on all of \(\NN_{\ge1}\).  Thus a level-zero-only replacement fails
precisely in the restricted model, while \Cref{thm:two-level} supplies the next
readout level and realizes multiplication there by the same semantic operation,
composition.
\end{proof}


\begin{definition}[Protocol universality]\label{def:prot-univ}
An \(L1\) protocol for operations on a countable set \(S\) consists of the
following data.
\begin{enumerate}
\item A decidable raw-code language \(C\subseteq\{0,1\}^{<\infty}\) and a
  decidable set \(N\) of normal forms.
\item A computable normalizer \(\mathrm{NF}\colon C\to N\) and a computable
  decoding map \(\mathrm{Dec}\colon N\to S\).  Equality of normal forms in \(N\)
  is decidable, and \(\mathrm{Dec}\) is injective on the represented normal
  forms.
\item For each target operation \(\star\) of arity \(r\), a computable compiled
  operation \(\widehat\star\colon N^r\to N\) satisfying
  \[
    \mathrm{Dec}\bigl(\widehat\star(n_1,\ldots,n_r)\bigr)
      =\mathrm{Dec}(n_1)\star\cdots\star\mathrm{Dec}(n_r).
  \]
\item Semantic readouts from the represented normal forms to objects built from
  the permitted primitive, with identities that read the target operations by
  that primitive after compilation.
\end{enumerate}
The protocol is universal for a class of operations when these data are supplied
for every operation in the class.
\end{definition}

In the Zeckendorf protocol, the compiled arithmetic operations on normal forms
and the semantic readouts \(\Phi_0,\Phi_1\) are separate pieces of data.  The
raw codes are finite binary words, the normal forms are the finite-support
Zeckendorf words \(X_\infty^{\mathrm{fin}}\), the normalizer is the greedy
Zeckendorf map, the decoding is \(\Val\), the compiled operations are
\(\oplus,\otimes\) from \Cref{def:stable-arith}, and the readouts are the maps
in \Cref{thm:two-level}.  Thus the code layer, normalizer, decidable normal
forms, decoding map, compiled operations, and semantic readouts are distinct
pieces of data.  All computability in the \(L1\) construction occurs on
\(X_\infty^{\mathrm{fin}}\), by
\Cref{lem:greedy-zeckendorf-normalizer,prop:value-transport-exact-interface}; no
bounded-window or digit-local arithmetic algorithm is part of the protocol.

\begin{proposition}[Zeckendorf protocol data]
\label{prop:zeckendorf-protocol-data}
The Zeckendorf construction supplies all data required by
\Cref{def:prot-univ} for \(S=\NN\) and target operations \(\{+,\times\}\):
\[
  C=\{0,1\}^{<\infty},\qquad
  N=X_\infty^{\mathrm{fin}},\qquad
  \mathrm{Dec}=\Val.
\]
Here \(\mathrm{NF}\) is the greedy Zeckendorf normalizer, the compiled operations
are \(\oplus,\otimes\), and the semantic readouts are \(\Phi_0,\Phi_1\).  These
data make the addition and multiplication diagrams of \Cref{def:prot-univ}
commute and read addition and multiplication using composition as the only
semantic primitive.
\end{proposition}

\begin{proof}
The language \(\{0,1\}^{<\infty}\) is decidable.  The set
\(X_\infty^{\mathrm{fin}}\) is decidable because a finite word can be checked for
finite support and for the absence of adjacent ones, and equality of such words
is literal equality after trailing zeros are ignored.  For a raw word
\(\omega=(\omega_1,\ldots,\omega_\ell)\), the integer
\(\sum_{k=1}^{\ell}\omega_kF_k\) is computable; applying
\Cref{lem:greedy-zeckendorf-normalizer} to that integer gives a computable
normalizer \(\mathrm{NF}\colon C\to X_\infty^{\mathrm{fin}}\).  The same lemma
also says that \(\Val\colon X_\infty^{\mathrm{fin}}\to\NN\) is a bijection, hence
an injective computable decoding map.

By \Cref{def:stable-arith,thm:semiring}, the operations \(\oplus\) and
\(\otimes\) are computable and satisfy
\[
  \Val(u\oplus v)=\Val(u)+\Val(v),
  \qquad
  \Val(u\otimes v)=\Val(u)\Val(v).
\]
Thus the compiled-operation diagrams in \Cref{def:prot-univ} commute for
addition and multiplication.  Finally, \Cref{thm:two-level} gives
\(\Phi_0(c_{a+b})=\Phi_0(c_a)\circ\Phi_0(c_b)\) and
\(\Phi_1(c_{ab})=\Phi_1(c_a)\circ\Phi_1(c_b)\).  These are exactly the two
semantic readout identities, both using composition and no second semantic
primitive.
\end{proof}

\begin{theorem}[Restricted level-zero protocol dichotomy]
\label{thm:restricted-level-zero-protocol-dichotomy}
\label{cor:level-zero-protocol-dichotomy}
Let a protocol for multiplication on \(\NN_{\ge1}\) use a single free
monogenic orbit \(C=\{g^n:n\in\NN\}\) as its semantic multiplication target, and
suppose the multiplication readout is a map
\(\eta\colon\NN_{\ge1}\to C\) satisfying
\[
  \eta(ab)=\eta(a)\eta(b)
  \qquad(a,b\ge1),
\]
where the product on the right is the ambient monoid product on \(C\).  Thus
the assertion is made under this level-zero monogenic-orbit restriction.  Then
\(\eta\) is not injective.  Consequently any protocol that realizes
\(\{+,\times\}\) while satisfying this same level-zero one-free-monogenic-orbit
restriction must, for its multiplication readout, leave the one free monogenic
orbit, abandon ambient-product readout, or move to a higher readout level such
as the power-operator level of \Cref{thm:two-level}.  The theorem makes no
global rigidity assertion about protocols outside this restricted model.
\end{theorem}

\begin{proof}
The hypotheses are exactly the hypotheses of \Cref{thm:level-zero-obstruction},
so the readout \(\eta\) is not injective.  Hence an injective protocol cannot
simultaneously keep multiplication in one free monogenic orbit and read it by
the ambient product.  The displayed alternatives are the logical negations of
those two restrictions, together with the positive realized alternative:
\Cref{thm:two-level} reads multiplication by composition of the operators
\(I_n\) in \(\End(\End(A))\), not by the ambient product on \(C\).
\end{proof}

\begin{theorem}[Main-chain protocol rigidity boundary]
\label{prop:no-global-protocol-rigidity-input}
\label{thm:main-chain-protocol-rigidity-boundary}
The hierarchy theorem uses only the restricted dichotomy of
\Cref{thm:restricted-level-zero-protocol-dichotomy}.  It does not assert a
classification or rigidity theorem for all single-primitive protocols realizing
\(\{+,\times\}\).  Equivalently, every negative protocol inference in
\Cref{thm:main-hierarchy} factors through the one-free-monogenic-orbit hypothesis
and the ambient-product multiplication readout; protocols outside that
restricted model are judged only by the positive requirements of
\Cref{def:prot-univ} and by the explicit Zeckendorf realization of
\Cref{thm:zeckendorf-protocol-univ}.
\end{theorem}

\begin{proof}
The positive protocol is the value-transported Zeckendorf protocol of
\Cref{thm:zeckendorf-protocol-univ}; its multiplication readout is the
higher-level identity \(\Phi_1(c_{ab})=\Phi_1(c_a)\circ\Phi_1(c_b)\) from
\Cref{thm:two-level}.  The negative protocol input cited by
\Cref{thm:main-hierarchy} is \Cref{thm:level-zero-obstruction}, equivalently the
restricted dichotomy just proved.  Neither statement ranges over arbitrary
protocol data in \Cref{def:prot-univ}: the negative statement assumes a free
monogenic orbit \(\{g^n:n\in\NN\}\) and asks that multiplication be read by the
ambient product on that same orbit.  If either hypothesis is removed, the proof
has no exponent map \(e(ab)=e(a)+e(b)\) and hence no prime-weight skeleton from
\Cref{lem:level-zero-exponent-skeleton}.  Thus the only proved obstruction is
the restricted one, while the constructive Zeckendorf protocol supplies the
separate positive realization.  No global rigidity classification is therefore
used, proved, or needed in the main theorem chain.
\end{proof}

\begin{theorem}[Value-transported Zeckendorf protocol universality]
\label{thm:zeckendorf-protocol-univ}
The raw input grammar is the decidable set \(\{0,1\}^{<\infty}\) of finite
binary words.  The stable normalizer
\[
  \mathrm{NF}_\infty\colon \{0,1\}^{<\infty}\to X_\infty^{\mathrm{fin}}
\]
sends a word \(\omega=(\omega_1,\ldots,\omega_\ell)\) to the greedy Zeckendorf
normal form of \(\sum_{k=1}^\ell \omega_kF_k\).  Together with decoding by
\(\Val\), compiled operations \(\oplus,\otimes\), and readouts
\(\Phi_0,\Phi_1\), these data satisfy \Cref{def:prot-univ} and form an \(L1\)
protocol universal for \(\{+,\times\}\) on \(\NN\).  It is a value-transported
protocol and asserts no bounded-window or digit-local arithmetic theorem.
\end{theorem}

\begin{proof}
The code language, normal forms, normalizer, decoding, compiled operations, and
readouts are exactly the verified protocol data of
\Cref{prop:zeckendorf-protocol-data}.  The compiled operations are computable by
\Cref{prop:zeckendorf-computable-arithmetic-anchor,lem:value-transported-suffices},
and the semantic readout identities are \Cref{thm:two-level}.  Therefore the
data are an \(L1\) protocol universal for \(\{+,\times\}\) on \(\NN\).  The final
sentence is \Cref{prop:value-transport-exact-interface}.
\end{proof}

The positive \(L1\) result uses no operation on \(\NN\) beyond ordinary addition
and multiplication transported through the Zeckendorf valuation; this is exactly
\Cref{def:stable-arith} and \Cref{thm:zeckendorf-protocol-univ}.


% End original section: sec03_protocol_universality.tex
% Original section: sec04_certificate_boundary.tex
\begin{corollary}[L2 load split]\label{cor:L2-load-split}
The \(L2\) positive clause is the finite-fibre certificate construction; the
strictness clause is the separate transfer obstruction of \Cref{thm:finite-state-strictness}.
\end{corollary}

\begin{proof}
The positive clause is \Cref{prop:zeckendorf-cert-univ}.  The obstruction is the
failure of the finite-state budget, which is not a fixed-resolution moment
claim by
\Cref{prop:L2-novelty-finite-state-separation,thm:L2-finite-window-transfer-obstruction}.
\end{proof}

\begin{proposition}[Terminology bridge for the \texorpdfstring{\(L2\)}{L2} clause]
\label{thm:L2-terminology-bridge}
In this paper, \(L2\)-certificate universality means finite-fibre certificate
universality in the sense of \Cref{def:cert-univ}.  A finite-state transfer
formula is a strictly additional hypothesis on that certificate data.  Hence the
\(L2\) strictness statements are exactly cover-relative failures of this extra
finite-state transfer hypothesis, not failures of finite-fibre certificate
universality itself.
\end{proposition}

\begin{proof}
The first sentence is the definition of \(L2\)-certificate universality in
\Cref{def:cert-univ}: it asks for effective finite-resolution data whose
fibre-size multisets supply the fixed finite certificates of
\Cref{prop:classical-finite-atomic-skeleton}.  The same definition then states
finite-state transfer separately, by requiring for each fixed \(q\) a single
matrix expression \(S_q^{(m)}=u_q^{\mathsf T}A_q^m v_q\) valid for all
resolutions \(m\).  \Cref{prop:L2-novelty-finite-state-separation} proves that
this uniform-in-\(m\) matrix condition is not contained in the fixed-resolution
certificate data, and \Cref{cor:L2-load-split} identifies the positive clause
with \Cref{prop:zeckendorf-cert-univ} while assigning strictness to the transfer
budget obstruction.  Finally,
\Cref{thm:L2-finite-fiber-no-transfer-content} proves the needed obstruction
inside the computable finite-fibre cover category over the same normal forms.
Thus the terminology has the asserted unique reading.
\end{proof}

\begin{corollary}[Subordinate role of L2 finite-fibre facts]
\label{cor:L2-subordinate-role}
The finite-fibre facts do not prove the Richardson obstruction, the one-orbit
multiplication obstruction, or the Zeckendorf \(L1\) protocol.  They only attach
finite certificates after a normalizer has already been chosen.  In particular,
no result labelled in \Cref{sec:consequences} is a premise for any theorem in
\Cref{sec:L0,sec:L1,sec:strictness}.
\end{corollary}

\begin{proof}
Every result in this section starts from finite fibres of a given normalizing
map.  The \(L1\), Richardson, and one-orbit arguments are proved in
\Cref{sec:L1,sec:strictness} and do not use these certificates.
\end{proof}

\begin{theorem}[Cover moments do not determine intrinsic fold transfer]
\label{thm:L2-cover-intrinsic-nonimplication}
Let
\[
  d_m(x)=|\Fold_m^{-1}(x)|,
  \qquad S_2(m)=\sum_{x\in X_m}d_m(x)^2
\]
be the intrinsic second collision moment of the unmodified Zeckendorf fold.  Let
\(d_m^\sharp(x)\) be the fibre-size function of a computable finite-fibre cover
over the same normal-form set \(X_m\), and let
\[
  S_2^\sharp(m)=\sum_{x\in X_m}d_m^\sharp(x)^2 .
\]
The assertion that \((S_2^\sharp(m))\) has, or fails to have, a finite-state
transfer formula is a statement about the cover-size data \(d_m^\sharp\).  It
does not imply the corresponding assertion for the intrinsic sequence
\((S_2(m))\) unless an additional theorem identifies the two sequences, or gives
a transfer-preserving comparison between them, for all \(m\).  No such
identification or comparison is part of the \(L2\) strictness proof in this
article.
\end{theorem}

\begin{proof}
Both moments are finite sums over the same finite set \(X_m\), but their summands
come from different data.  The intrinsic summand \(d_m(x)\) is defined by the
original map \(\Fold_m\colon\Omega_m\to X_m\).  A computable finite-fibre cover,
by contrast, is specified by a separate computable positive integer function
\(d_m^\sharp(x)\) over \(X_m\); by \Cref{def:computable-cover} and
\Cref{thm:finite-fiber-cover-realization}, this function may be chosen without
changing \(X_m\), \(\Fold_m\), or any readout on normal forms.  Thus the cover
moment can be altered while the intrinsic fold moment remains the moment of the
unchanged map.

Finite-state transfer is a property of a particular sequence in \(m\): by
\Cref{def:cert-univ}, it asks for a fixed representation
\(u^{\mathsf T}A^m v\) for that sequence.  A representation, recurrence, or
growth obstruction for \((S_2^\sharp(m))\) therefore concerns only the sequence
computed from \(d_m^\sharp\).  It can be transported to \((S_2(m))\) only after a
separate equality or transfer-preserving comparison between the two sequences has
been proved.  The one-spike construction of
\Cref{thm:L2-finite-fiber-no-transfer-content} proves the opposite structural
fact needed here: it deliberately changes one cover fibre over an otherwise
fixed normal-form system.  Hence the cover-relative finite-state obstruction has
no intrinsic Fold\(_m\) nontransfer consequence in the present theorem chain; the
positive intrinsic transfer assertion is the separate theorem
\Cref{thm:intrinsic-fold-moment-transfer}.
\end{proof}

\begin{proposition}[Cover-category invariance of the finite-state obstruction]
\label{prop:cover-category-invariance-finite-state-obstruction}
\label{thm:cover-relative-transfer-is-not-intrinsic}
Fix the unmodified Zeckendorf normal-form system \((X_m,\Fold_m)_{m\ge1}\).  The
assignment
\[
  d_m^\sharp\longmapsto S_2^\sharp(m)=\sum_{x\in X_m}d_m^\sharp(x)^2
\]
from computable finite-fibre cover-size functions to second covered moments is
not determined by \((X_m,\Fold_m)\): over this same intrinsic system there is a
singleton computable cover whose second moment has fixed finite-state transfer,
and there is a factorial one-spike computable cover whose second moment has no
fixed finite-state transfer.  Hence finite-state transfer or nontransfer of
covered moments is not an intrinsic invariant of the fold map; it is an
additional property of the chosen cover-size data.
\end{proposition}

\begin{proof}
For the singleton cover take \(d_m^\sharp(x)=1\) for every \(x\in X_m\).  Then
\(S_2^\sharp(m)=|X_m|=F_{m+2}\), and the Fibonacci recurrence gives a fixed
finite-state transfer formula.  For the second cover, choose one normal form in
each \(X_m\), give it cover size \(m!+1\), and give all other normal forms cover
size \(1\).  The construction is computable and changes no normal form or
readout, so it is a cover of the same intrinsic system.  Its second moment is
\((m!+1)^2+F_{m+2}-1\), which grows faster than every exponential.  By
\Cref{prop:finite-state-coefficient-budget}, every fixed finite-state transfer
sequence is exponentially bounded; hence this second covered moment has no fixed
finite-state transfer.

Both covers lie over the identical data \((X_m,\Fold_m)\).  Since they have
opposite transfer behaviour, that behaviour cannot be a function of the
intrinsic fold map alone.  It is therefore exactly a cover-relative property of
the auxiliary multiplicities \(d_m^\sharp\).
\end{proof}
\begin{theorem}[Cover finite-fibre certificates do not imply transfer]
\label{thm:L2-finite-fiber-no-transfer-content}
\label{thm:cover-relative-finite-fibre-no-transfer}
In the Zeckendorf computable finite-fibre cover category, and only in that
cover-relative category, finite-fibre certificate data alone impose no
recurrence, rational generating series, exponential growth bound, or fixed
finite-matrix transfer formula on the cover moments \((S_q^\sharp(m))_{m\ge1}\).
More precisely, for every computable sequence \(a_m\ge1\) there is a computable
finite-fibre cover preserving the normal forms and readouts whose second moment is
\[
  S_2^\sharp(m)=(a_m+1)^2+F_{m+2}-1 .
\]
Taking \(a_m=m!\) gives a computable finite-fibre cover certificate system with
no finite-state transfer formula for \(S_2^\sharp(m)\).  No intrinsic failure
theorem for the unmodified \(\Fold_m\) fibre multiplicities is asserted.
\end{theorem}

\begin{proof}
At resolution \(m\), choose one normal form and assign its cover fibre size
\(a_m+1\), while assigning cover fibre size \(1\) to every other normal form.
This finite cover does not change the normal-form set or any readout on normal
forms, and it is computable because \((a_m)\) is computable.  The displayed
formula follows by adding the square of the enlarged fibre and the squares of
the remaining singleton cover fibres.  If \(a_m=m!\), then
\((S_2^\sharp(m))\) grows faster than
\(C\rho^m\) for every \(C,
\rho\), contradicting the exponential budget in
\Cref{prop:finite-state-coefficient-budget}.  Hence no fixed finite-state
transfer formula exists.
\end{proof}

\begin{theorem}[Two-cover skeleton for the intrinsic boundary]
\label{thm:L2-two-cover-intrinsic-boundary}
Fix the unmodified Zeckendorf data \((X_m,\Fold_m)\).  There are two computable
finite-fibre covers of this same data, both preserving all normal forms and all
\(L1\) readouts, with opposite finite-state-transfer behaviour for their second
collision moments:
\begin{enumerate}
\item the singleton cover has
  \(S_{2,0}^\sharp(m)=\lvert X_m\rvert=F_{m+2}\), hence satisfies a fixed
  finite-state transfer formula; and
\item the factorial one-spike cover has
  \(S_{2,!}^\sharp(m)=(m!+1)^2+F_{m+2}-1\), hence has no fixed finite-state
  transfer formula.
\end{enumerate}
Consequently finite-state transfer or failure for cover moments is not an
intrinsic invariant of the original \(\Fold_m\) fibres.  Any theorem about
recurrence, nonrecurrence, rational generating series, or finite-state transfer
for the intrinsic \(\Fold_m\) multiplicity sequence requires an additional bridge
identifying that intrinsic sequence with the chosen cover sequence; no such
bridge is part of the present article.
\end{theorem}

\begin{proof}
For the first cover, attach one certificate point over every normal form.  The
second collision moment is then the number of normal forms, namely
\(\lvert X_m\rvert=F_{m+2}\) by \Cref{prop:finite-zeckendorf-residue-interface}.  The Fibonacci
sequence satisfies the fixed two-state recurrence
\[
  \binom{F_{m+2}}{F_{m+3}}
  =
  \begin{pmatrix}0&1\\ 1&1\end{pmatrix}^{m}
  \binom{F_2}{F_3},
\]
after the harmless shift of the initial vector, so this moment sequence has a
finite-state transfer formula.

For the second cover, apply \Cref{thm:L2-finite-fiber-no-transfer-content} with
\(a_m=m!\).  It changes only one auxiliary certificate fibre over a fixed normal
form, preserves all readouts by construction, and has displayed moment
\((m!+1)^2+F_{m+2}-1\).  This sequence grows faster than every exponential, while
\Cref{prop:finite-state-coefficient-budget} gives an exponential bound for every
fixed finite-state transfer sequence; hence the factorial one-spike moment has
no finite-state transfer formula.

Both covers project to the identical unmodified data \((X_m,\Fold_m)\).  Since a
property of the auxiliary cover moment takes different truth values on two
covers with the same intrinsic fold map, that property is not determined by the
intrinsic fold multiplicities.  A conclusion for the intrinsic sequence would
therefore require an added equality or transfer-preserving comparison between
that intrinsic sequence and the chosen cover moment, exactly the bridge excluded
in \Cref{thm:L2-cover-intrinsic-nonimplication}.
\end{proof}

\begin{proposition}[Finite-fibre certificate skeleton and transfer pairing]
\label{thm:L2-fiber-multiset-factorization}
\label{thm:L2-finite-fiber-transfer-pairing}
For every fixed resolution \(m\), the \(L2\) certificates used in this paper
factor through the fibre-size multiset \(\mathcal D_m\).  If the same certificate
system also has a fixed finite-state transfer formula for \(S_q^{(m)}\), then
\((S_q^{(m)})_m\) satisfies a constant-coefficient recurrence and a rational
generating series.  Conversely, there are computable one-spike covers preserving
the same normal forms and readouts whose fixed-resolution finite-fibre
certificates exist at every \(m\) but whose second collision moments admit no
fixed finite-state transfer formula.
\end{proposition}

\begin{proof}
The fixed-resolution factorization is exactly
\Cref{prop:finite-fiber-certificate-taxonomy}: the finite-fibre certificates used
here are functions of \(\mathcal D_m\), hence are unchanged by replacing a datum
with another datum having the same fibre-size multiset.  If a fixed finite-state transfer formula
exists, \Cref{def:fixed-finite-state-transfer} gives
\(S_q^{(m)}=u_q^{\mathsf T}A_q^m v_q\) with one finite matrix model independent of
\(m\).  Applying \Cref{prop:finite-state-coefficient-budget} yields the stated
constant-coefficient recurrence and rational generating series.

For the converse, use \Cref{thm:L2-finite-fiber-no-transfer-content} with
\(a_m=m!\).  The construction changes only one auxiliary cover fibre over a fixed
normal form, so it preserves the normal-form set and all readouts, while its
second moment is \((m!+1)^2+F_{m+2}-1\).  This sequence grows faster than every
exponential sequence, contradicting the exponential consequence of any fixed
finite-state transfer formula in \Cref{prop:finite-state-coefficient-budget}.
Thus finite-fibre certificates and fixed finite-state transfer are paired only
when the additional transfer structure is explicitly supplied.
\end{proof}

\begin{corollary}[Cover finite-fibre certificates are independent of transfer]
\label{cor:hankel-psd-independent}
In the Zeckendorf cover category, the fixed-resolution finite-fibre certificates
used here can hold at every resolution while finite-state transfer fails.
\end{corollary}

\begin{proof}
They hold for every finite-fibre system by
\Cref{prop:classical-finite-atomic-skeleton}; finite-state transfer fails for the
factorial one-spike cover in \Cref{thm:L2-finite-fiber-no-transfer-content}.
\end{proof}



% End original section: sec04_certificate_boundary.tex
% Original section: sec04_certificate_universality.tex
\section{\texorpdfstring{\(L2\)}{L2}: finite certificates}\label{sec:L2}

An \(L1\) protocol supplies normal forms.  Level \(L2\) records computable
finite invariants of the normalizing fibres.  For the unmodified Zeckendorf
system these fibres are the fibres of \(\Fold_m\colon\Omega_m\to X_m\).  The
intrinsic fibre moments of this particular fold have finite-state transfer
degree by degree.  The strictness theorem below is a complementary cover
statement: in the computable finite-fibre cover category over the same normal
forms, fixed-resolution finite fibres give the usual finite-atomic moment
certificates, but arbitrary auxiliary cover fibres need not preserve the
intrinsic transfer structure in the resolution variable.

For \(x\in X_m\), put
\[
  d_m(x)=|\Fold_m^{-1}(x)|,
  \qquad
  S_q(m)=\sum_{x\in X_m} d_m(x)^q\quad(q\ge0).
\]
Thus \(S_0(m)=|X_m|=F_{m+2}\) and \(S_1(m)=|\Omega_m|=2^{m+1}\).

To keep the intrinsic tower visibly distinct from auxiliary covers, we also
write
\[
  S_q^{\mathrm{int}}(m)=\sum_{x\in X_m}|\Fold_m^{-1}(x)|^q .
\]
Thus \(S_q^{\mathrm{int}}(m)=S_q(m)\) for the unmodified fold, while the
superscript \(\sharp\) below is reserved for cover moments.

\begin{definition}[Collision moment]\label{def:collision-moment}
For any finite normal-form set \(N_m\) with finite fibres, the \(q\)-th
collision moment is
\begin{equation}\label{eq:Sq}
  S_q^{(m)}=\sum_{n\in N_m}|\mathrm{NF}_m^{-1}(n)|^q .
\end{equation}
\end{definition}

\begin{definition}[Fixed finite-state transfer]\label{def:fixed-finite-state-transfer}
Let \((b_m)_{m\ge m_0}\) be a scalar sequence over a fixed characteristic-zero
field \(K\) containing the values under discussion, for instance
\(K=\mathbb Q\) or \(K=\RR\).  The sequence has a \emph{fixed finite-state
transfer formula} if
there are an integer \(r\ge1\), a single matrix \(A\in M_r(K)\), and vectors
\(u,v\in K^r\), all independent of \(m\), such that
\[
  b_m=u^{\mathsf T}A^m v\qquad(m\ge m_0).
\]
For an effective finite-resolution datum, a collision-moment family has fixed
finite-state transfer in degree \(q\) when the sequence
\((S_q^{(m)})_{m\ge m_0}\) has such a formula.  It has fixed finite-state
transfer if this holds degree by degree.  The dimension \(r\), matrix \(A\), and
vectors may depend on \(q\), but for each fixed \(q\) they are fixed once and for
all and do not vary with the resolution \(m\).
\end{definition}

\begin{lemma}[Bounded signed-Fibonacci carry automaton]
\label{lem:bounded-signed-fibonacci-carry}
\label{lem:bounded-signed-fibonacci-carry-automaton}
Fix \(q\ge1\).  For a column
\(\alpha=(\alpha_1,\ldots,\alpha_q)\in\{0,1\}^q\), put
\[
  \Delta(\alpha)=(\alpha_2-
  \alpha_1,\ldots,\alpha_q-\alpha_1)
  \in\{-1,0,1\}^{q-1}.
\]
For \(a\in\{-1,0,1\}^{q-1}\), let
\[
  w_q(a)=\#\{\alpha\in\{0,1\}^q:\Delta(\alpha)=a\}.
\]
Then \(S_q^{\mathrm{int}}(m)\) is the weighted number of words
\(a_1,\ldots,a_{m+1}\in\{-1,0,1\}^{q-1}\), with weight
\(\prod_k w_q(a_k)\), for which there is a target marker
\(\sigma\in\{-1,0,1\}^{q-1}\) satisfying
\begin{equation}\label{eq:q-column-fib-target}
  \sum_{k=1}^{m+1} a_{k,j}F_k=\sigma_jF_{m+2}
  \qquad(1\le j\le q-1).
\end{equation}
These weighted words are recognized completely by a finite automaton whose carry states are
\[
  \mathcal C_q=\{-1,0,1\}^{q-1}\times\{-1,0,1\}^{q-1}.
\]
The initial state is \((0,0)\).  Reading a column
\(a\in\{-1,0,1\}^{q-1}\) sends \((p,r)\) to
\[
  \tau_a(p,r)=(r,p+a-r)
\]
when the second coordinate again lies in \(\{-1,0,1\}^{q-1}\), and has no
transition otherwise; the transition weight is \(w_q(a)\).  After
\(m+1\) columns, the accepting states are
\(\{(\sigma,0):\sigma\in\{-1,0,1\}^{q-1}\}\).  The weighted path count is
exactly \(S_q^{\mathrm{int}}(m)\).  Equivalently, no solution of
\eqref{eq:q-column-fib-target} is discarded by the bounded state restriction,
and every accepting bounded path satisfies \eqref{eq:q-column-fib-target} for a
unique target marker \(\sigma\).
\end{lemma}

\begin{proof}[Proof of \Cref{lem:bounded-signed-fibonacci-carry}]
We prove the lemma in the adjacent block.  The first sublemma classifies every
possible scalar exit from the bounded square; the second applies that scalar
classification coordinate by coordinate and proves that deleting transitions
leaving \(\mathcal C_q\) loses no target solution.  After those two sublemmas are
proved, the closing part of this same proof rewrites equality of
\(\Fold_m\)-fibres as \eqref{eq:q-column-fib-target} and applies the deletion
lemma.  Thus the bounded state restriction is not an independent hypothesis: the
displayed automaton, the two exit sublemmas, and the closing argument below are
one closed proof package, and every later reference to
\Cref{lem:bounded-signed-fibonacci-carry} means this entire package.

\begin{lemma}[Scalar exit classification for signed Fibonacci carries]
\label{lem:scalar-signed-fibonacci-exit-classification}
For the scalar transition
\[
  \tau_a(p,r)=(r,p+a-r),\qquad a\in\{-1,0,1\},
\]
every path starting in \((0,0)\) which ever leaves the square
\(Q=\{-1,0,1\}^2\) is trapped forever outside \(Q\) and can never reach a state
\((\sigma,0)\), \(\sigma\in\{-1,0,1\}\).  More precisely, the only one-step exits
from \(Q\) are
\[
  (0,-2),(1,-2),(1,-3),(-1,2),(-1,3),(0,2),
\]
and they fall into the following two forward-invariant classes:
\[
\begin{array}{c|c|c}
\text{exit class} & \text{first region} & \text{alternate region}\\ \hline
\text{positive} & U_+=\{p\le0,\ r\ge2\} & V_+=\{p\ge2,\ r\le-1\}\\
\text{negative} & U_-=\{p\ge0,\ r\le-2\} & V_-=\{p\le-2,\ r\ge1\}
\end{array}
\]
with \(\tau_a(U_+)\subseteq V_+\), \(\tau_a(V_+)\subseteq U_+\),
\(\tau_a(U_-)\subseteq V_-\), and \(\tau_a(V_-)\subseteq U_-\) for every
\(a\in\{-1,0,1\}\).
\end{lemma}

\begin{proof}
If \((p,r)\in Q\), then the first coordinate of \(\tau_a(p,r)\) is again in
\(\{-1,0,1\}\).  Hence a one-step exit occurs exactly when
\(p+a-r\notin\{-1,0,1\}\).  Checking the twenty-seven triples
\((p,r,a)\in\{-1,0,1\}^3\) gives precisely the six displayed exit states.  The
three exits with second coordinate positive lie in \(U_+\), and the three exits
with second coordinate negative lie in \(U_-\).

For \((p,r)\in U_+\), the image has first coordinate \(r\ge2\) and second
coordinate \(p+a-r\le 0+1-2=-1\), so it lies in \(V_+\).  For
\((p,r)\in V_+\), the image has first coordinate \(r\le-1\) and second coordinate
\(p+a-r\ge2-1-(-1)=2\), so it lies in \(U_+\).  Thus the positive class is
forward invariant and disjoint from \(Q\).  Similarly, if \((p,r)\in U_-\), then
\(r\le-2\) and \(p+a-r\ge0-1-(-2)=1\), so \(\tau_a(p,r)\in V_-\); if
\((p,r)\in V_-\), then \(r\ge1\) and \(p+a-r\le-2+1-1=-2\), so
\(\tau_a(p,r)\in U_-\).  The negative class is also forward invariant and
disjoint from \(Q\).  Since every accepting state \((\sigma,0)\) lies in \(Q\), a
path that has exited cannot later accept.
\end{proof}

\begin{lemma}[Coordinatewise exit deletion for signed Fibonacci carries]
\label{lem:coordinatewise-signed-fibonacci-exit-deletion}
Fix \(q\ge1\) and a length \(n\ge1\).  For columns
\(a_1,\ldots,a_n\in\{-1,0,1\}^{q-1}\), let the unrestricted vector carry states be
defined by
\[
  (p_0,r_0)=(0,0),
  \qquad (p_k,r_k)=(r_{k-1},p_{k-1}+a_k-r_{k-1})
  \quad(1\le k\le n),
\]
coordinate by coordinate.  If some coordinate of some state \((p_k,r_k)\) leaves
the scalar square \(Q=\{-1,0,1\}^2\), then the vector path can never end in an
accepting state \((\sigma,0)\), \(\sigma\in\{-1,0,1\}^{q-1}\).  Consequently the
bounded automaton obtained by deleting all transitions leaving
\(\mathcal C_q=\{-1,0,1\}^{q-1}\times\{-1,0,1\}^{q-1}\) deletes no solution of the
target equations
\[
  \sum_{k=1}^n a_{k,j}F_k=\sigma_jF_{n+1}
  \qquad(1\le j\le q-1).
\]
\end{lemma}

\begin{proof}
For \(q=1\) there are no scalar coordinates and the assertion is vacuous.  Assume
\(q\ge2\).  Fix a coordinate \(j\).  Its scalar states satisfy exactly the scalar
recursion of \Cref{lem:scalar-signed-fibonacci-exit-classification}.  Hence, if
this coordinate is the first coordinate to leave \(Q\), or if it leaves after
some other coordinate has already left, the scalar classification applies from
that exit time onward: the coordinate is trapped in one of the two invariant exit
classes and never reaches any scalar accepting state \((\sigma_j,0)\).  A vector
accepting state requires every coordinate to be scalar accepting, so any vector
path with any coordinate exit is nonaccepting.

It remains to connect nonacceptance with the displayed target equations.  The
coordinatewise induction
\[
  \sum_{i=1}^k a_{i,j}F_i=p_{k,j}F_{k+1}+r_{k,j}F_k
\]
is the scalar carry identity applied to coordinate \(j\).  If the \(j\)-th target
equation holds at length \(n\), then
\(p_{n,j}F_{n+1}+r_{n,j}F_n=\sigma_jF_{n+1}\).  The elementary bounds
\(|p_{n,j}|\le F_n-1\) and \(|r_{n,j}|\le F_{n+1}-1\), obtained from the same
recursion by induction, force \(r_{n,j}=0\) because \(F_{n+1}\mid r_{n,j}\), and
then force \(p_{n,j}=\sigma_j\).  Thus a simultaneous target solution must end in
the vector accepting state \((\sigma,0)\).  Since every path that ever leaves
\(\mathcal C_q\) is nonaccepting, deleting the leaving transitions cannot delete
any target solution.
\end{proof}

We now close the proof of the bounded signed-Fibonacci carry automaton.  The only
extra input used after the statement is the scalar exit classification of
\Cref{lem:scalar-signed-fibonacci-exit-classification}, packaged coordinate by
coordinate in \Cref{lem:coordinatewise-signed-fibonacci-exit-deletion}.
A \(q\)-tuple \((\omega^{(1)},\ldots,\omega^{(q)})\in\Omega_m^q\) contributes
exactly when all \(q\) words fold to the same normal form.  By
\eqref{eq:fold-residue} and the bijection \(\Val_m\colon X_m\to
\mathbb Z/F_{m+2}\mathbb Z\), this is equivalent to the \(q-1\) congruences
\[
  \sum_{k=1}^{m+1}(\omega^{(j+1)}_k-
  \omega^{(1)}_k)F_k\equiv0\pmod {F_{m+2}}
  \qquad(1\le j\le q-1).
\]
For one coordinate the absolute value of the left side is at most
\(\sum_{k=1}^{m+1}F_k=F_{m+3}-2<2F_{m+2}\).  Hence the congruence is equivalent
to equality with \(\sigma_jF_{m+2}\) for a unique
\(\sigma_j\in\{-1,0,1\}\).  Grouping digit differences by columns gives
\eqref{eq:q-column-fib-target}.  For a fixed difference column \(a_k\), the
number of binary \(q\)-columns producing it is exactly \(w_q(a_k)\), which gives
the stated product weight.

For the carry automaton it suffices to prove the scalar case, since the
\((q-1)\)-coordinate automaton is the direct product and the vector state is
bounded if and only if each scalar coordinate is bounded.  Put
\((p_0,r_0)=(0,0)\) and
\((p_k,r_k)=(r_{k-1},p_{k-1}+a_k-r_{k-1})\).  A direct induction using
\(F_{k+1}=F_k+F_{k-1}\) gives
\[
  \sum_{i=1}^k a_iF_i=p_kF_{k+1}+r_kF_k .
\]
For a word of length \(n\), the target equation
\(\sum_{i=1}^n a_iF_i=\sigma F_{n+1}\) is therefore equivalent to the final
identity \(p_nF_{n+1}+r_nF_n=\sigma F_{n+1}\).  The unrestricted carries also
satisfy the elementary bounds
\[
  |p_k|\le F_k-1,
  \qquad |r_k|\le F_{k+1}-1
  \qquad(k\ge1),
\]
proved by induction from \(p_k=r_{k-1}\) and
\(|r_k|\le |p_{k-1}|+1+|r_{k-1}|\le (F_{k-1}-1)+1+(F_k-1)=F_{k+1}-1\).  Since
\(\gcd(F_n,F_{n+1})=1\), the target identity implies
\(F_{n+1}\mid r_n\); the displayed bound forces \(r_n=0\), and then
\(p_n=\sigma\).  Conversely \((p_n,r_n)=(\sigma,0)\) plainly gives the target
equation.  Thus the unrestricted scalar recursion satisfies the target condition
exactly when its final state is \((\sigma,0)\).

It remains to justify the bounded state restriction.  This is exactly
\Cref{lem:coordinatewise-signed-fibonacci-exit-deletion}: any path with a
coordinate exit is nonaccepting, while any simultaneous target solution must end
at the vector accepting state.  Hence deleting transitions whose image is outside
\(\mathcal C_q\) deletes no target solution.  Conversely every bounded accepting
path satisfies the target equations by the displayed carry identity.

The vector transition is the direct product of the scalar transitions, its
accepting state records the unique vector
\(\sigma\in\{-1,0,1\}^{q-1}\), and the column weight \(w_q(a_k)\) counts exactly
the binary \(q\)-columns with difference column \(a_k\).  Therefore the bounded
weighted automaton counts all and only the \(q\)-tuples in the same fibre of
\(\Fold_m\), with the product weight stated above.  This is
\(S_q^{\mathrm{int}}(m)\).
\end{proof}

\begin{lemma}[Closed bounded-carry proof interface]
\label{lem:closed-bounded-carry-proof-interface}
Every later use in this paper of \Cref{lem:bounded-signed-fibonacci-carry} uses
the full adjacent proof chain: first
\Cref{lem:scalar-signed-fibonacci-exit-classification}, then
\Cref{lem:coordinatewise-signed-fibonacci-exit-deletion}, and only then
\Cref{lem:bounded-signed-fibonacci-carry},
and no independent boundedness assertion for signed Fibonacci carries is assumed.
Equivalently, the finite state set
\(\mathcal C_q=\{-1,0,1\}^{q-1}\times\{-1,0,1\}^{q-1}\) is lossless only for
the target equations \eqref{eq:q-column-fib-target} and only because the scalar
exit classes above prevent an exited coordinate from returning to an accepting
state.
\end{lemma}

\begin{proof}
The proof of \Cref{lem:bounded-signed-fibonacci-carry} first rewrites equality of
fold fibres as the coordinate target equations \eqref{eq:q-column-fib-target}.
It then proves the scalar carry identity
\(\sum_{i=1}^k a_iF_i=p_kF_{k+1}+r_kF_k\), identifies the scalar accepting states
\((\sigma,0)\), and invokes
\Cref{lem:coordinatewise-signed-fibonacci-exit-deletion}.  That lemma is exactly
the coordinatewise product of
\Cref{lem:scalar-signed-fibonacci-exit-classification}: once any scalar coordinate
leaves \(\{-1,0,1\}^2\), it stays in one of the two forward-invariant exit classes
and cannot end at \((\sigma,0)\).  Hence deleting transitions leaving
\(\mathcal C_q\) removes no solution of the target equations, while every bounded
accepting path satisfies those equations by the displayed carry identity.  These
are precisely the hypotheses used in
\Cref{thm:audited-fixed-degree-transfer-interface,thm:intrinsic-fold-moment-transfer};
there is no separate lemma or convention in the manuscript that asserts bounded
carry losslessness outside this chain.
\end{proof}

\begin{theorem}[Audited fixed-degree transfer interface]
\label{thm:audited-fixed-degree-transfer-interface}
Fix \(q\ge1\).  The finite-state transfer formula of
\Cref{thm:intrinsic-fold-moment-transfer} is obtained from the unmodified maps
\(\Fold_m\colon\Omega_m\to X_m\) by the following lossless finite construction:
\begin{enumerate}
\item the only deletion made in passing from unrestricted signed-Fibonacci
  carries to the finite state set
  \(\mathcal C_q=\{-1,0,1\}^{q-1}\times\{-1,0,1\}^{q-1}\) is the deletion of
  coordinatewise scalar exits classified in
  \Cref{lem:scalar-signed-fibonacci-exit-classification};
\item every deleted path is nonaccepting for every target marker
  \(\sigma\in\{-1,0,1\}^{q-1}\), and every solution of the fold congruences ends
  in the unique accepting state \((\sigma,0)\);
\item the matrix \(T_q\) is exactly the weighted adjacency matrix of this bounded
  automaton, with entry from a state \(c\) to a state \(c'\) equal to
  \(\sum_{\tau_a(c)=c'}w_q(a)\), and the vectors \(e_0,b_q\) are the
  initial-state indicator and accepting-state indicator.
\end{enumerate}
Consequently no cover multiplicity, truncated histogram, or resolution-dependent
state set is used in the intrinsic transfer theorem.  For each fixed \(q\), the
same finite matrix computes all moments \(S_q^{\mathrm{int}}(m)\) as \(m\)
varies, and this positive intrinsic assertion is independent of the
cover-relative obstruction in \Cref{thm:finite-state-strictness}.
\end{theorem}

\begin{proof}
The word counted in \(S_q^{\mathrm{int}}(m)\) is a \(q\)-tuple of words in
\(\Omega_m\) having the same value under the original fold map.  By
\eqref{eq:fold-residue} and the bijection \(\Val_m\colon X_m\to
\mathbb Z/F_{m+2}\mathbb Z\), this is equivalent to the signed Fibonacci
congruences displayed in the proof of \Cref{lem:bounded-signed-fibonacci-carry};
because the absolute signed sum is less than \(2F_{m+2}\), each congruence has the
unique target marker \(\sigma_j\in\{-1,0,1\}\).  Thus the target equations
\eqref{eq:q-column-fib-target} are exactly the fold-collision equations, not an
auxiliary cover condition.

For a scalar coordinate, the unrestricted carry recursion satisfies
\[
  \sum_{i=1}^k a_iF_i=p_kF_{k+1}+r_kF_k .
\]
If a path exits \(\{-1,0,1\}^2\),
\Cref{lem:scalar-signed-fibonacci-exit-classification} places it forever in one
of the two invariant exit classes, both disjoint from every scalar accepting
state \((\sigma,0)\).  Applying this coordinate by coordinate gives
\Cref{lem:coordinatewise-signed-fibonacci-exit-deletion}: a vector path with any
coordinate exit is nonaccepting.  Conversely, if the target equations hold, the
final carry identity and \(\gcd(F_n,F_{n+1})=1\) force the final state to be
\((\sigma,0)\); hence deleting the exits removes no target solution and adds none.

After the lossless deletion, the state set is finite and independent of the
resolution.  Reading a column \(a\in\{-1,0,1\}^{q-1}\) changes state by the fixed
rule \(\tau_a\) and contributes exactly the number \(w_q(a)\) of binary
\(q\)-columns with that difference column.  Therefore the weighted adjacency
matrix described in the statement has
\(e_0^{\mathsf T}T_q^{m+1}b_q\) equal to the weighted number of bounded
accepting paths of length \(m+1\), which is exactly
\(S_q^{\mathrm{int}}(m)\) by
\Cref{lem:bounded-signed-fibonacci-carry,lem:closed-bounded-carry-proof-interface}.
The construction uses only the fixed fold data \((\Omega_m,X_m,\Fold_m)\); the
finite-fibre cover construction of \Cref{thm:finite-state-strictness} changes
auxiliary certificate fibres over these fixed normal forms and is therefore a
separate strictness mechanism, not a premise of the transfer formula.
\end{proof}

\begin{theorem}[Intrinsic finite-state transfer for Zeckendorf fold moments]
\label{thm:intrinsic-fold-moment-transfer}
For every fixed \(q\ge0\), there are an effectively constructible integer matrix
\(T_q\) and integer vectors \(u_q,v_q\), independent of \(m\), such that
\[
  S_q^{\mathrm{int}}(m)=u_q^{\mathsf T}T_q^m v_q
  \qquad(m\ge1).
\]
Consequently \(\sum_{m\ge1}S_q^{\mathrm{int}}(m)z^m\) is a rational function and
\((S_q^{\mathrm{int}}(m))_{m\ge1}\) satisfies a constant-coefficient linear
recurrence with integer coefficients.
\end{theorem}

\begin{proof}
For \(q=0\), \(S_0^{\mathrm{int}}(m)=|X_m|=F_{m+2}\), so the Fibonacci matrix
supplies the transfer formula.  For \(q\ge1\), enumerate the finite carry state
set \(\mathcal C_q\) of \Cref{lem:bounded-signed-fibonacci-carry-automaton}.  Let
\(M_q\) be the integer matrix indexed by \(\mathcal C_q\) whose \((c,c')\)-entry
is the sum of the weights \(w_q(a)\) over all columns \(a\) with
\(\tau_a(c)=c'\).  Let \(e_0\) be the basis vector at \((0,0)\), and let \(b_q\)
be the characteristic vector of the accepting set.  Since an intrinsic
\(m\)-fold moment uses \(m+1\) columns,
\[
  S_q^{\mathrm{int}}(m)=e_0^{\mathsf T}M_q^{m+1}b_q
  =(M_q^{\mathsf T}e_0)^{\mathsf T}M_q^m b_q .
\]
Thus take \(T_q=M_q\), \(u_q=M_q^{\mathsf T}e_0\), and \(v_q=b_q\).  The
construction is effective because \(q\) fixes finite lists of states and columns.
Finally,
\[
  \sum_{m\ge1}S_q^{\mathrm{int}}(m)z^m
  =z\,u_q^{\mathsf T}(I-zT_q)^{-1}v_q,
\]
which is rational, and Cayley--Hamilton gives the stated recurrence.
\end{proof}

The transfer matrix can be trimmed without losing an accepted path.  Its
recurrent structure is stronger than rationality alone and supplies the
fixed-degree pressure which the height bounds by themselves do not determine.
The classical Fibonacci representation-function analyses of
\cite{Carlitz1968,Carlitz1970,KocabovaMasakovaPelantova2005} concern individual
representation counts and their extrema, while the golden-number Bernoulli
convolution analysis of \cite{Hu1997} concerns pointwise local dimensions.  The
result below concerns instead the positive-support transfer of the intrinsic
collision moments.

\begin{theorem}[Primitive intrinsic transfer and fixed-degree pressure]
\label{thm:primitive-intrinsic-moment-transfer}
For every fixed \(q\ge1\), let \(\mathcal T_q\subseteq\mathcal C_q\) be the set
of carry states which lie on a directed path of positive total weight from the
initial state to an accepting state in the automaton of
\Cref{lem:bounded-signed-fibonacci-carry-automaton}.  Equivalently,
\(\mathcal T_q\) is the trim set in the support digraph of the weighted
adjacency matrix \(M_q\) from
\Cref{thm:intrinsic-fold-moment-transfer}.  Let \(P_q\) be the weighted
adjacency matrix induced on \(\mathcal T_q\).  Then \(P_q\) is
primitive.  If \(\lambda_q>0\) is its Perron root, there are constants
\(c_q>0\) and \(0\le\theta_q<\lambda_q\) such that
\[
  S_q^{\mathrm{int}}(m)
  =c_q\lambda_q^m+O(\theta_q^m).
\]
Consequently the intrinsic fixed-degree pressure exists and is given by
\[
  p(q):=\lim_{m\to\infty}\frac1m\log S_q^{\mathrm{int}}(m)
  =\log\lambda_q .
\]
For \(q=0\), the same limit exists and equals \(\log\varphi\).
\end{theorem}

\begin{proof}
The columns of positive weight are exactly
\[
  \mathcal A_q
  =\{0,1\}^{q-1}\cup\{-1,0\}^{q-1}.
\]
We first show that every accepting marker which is reachable from the initial
state belongs to \(\mathcal A_q\).  Consider a path of positive total weight and
length \(n\) from the initial state to \((\sigma,0)\).  By
\Cref{lem:bounded-signed-fibonacci-carry}, it comes from binary words
\(\omega^{(1)},\ldots,\omega^{(q)}\) and satisfies
\[
  V_j-V_1=\sigma_{j-1}F_{n+1}\qquad(2\le j\le q),
  \quad
  V_j=\sum_{k=1}^n\omega_k^{(j)}F_k .
\]
Every \(V_j\) lies in the interval
\([0,F_{n+2}-2]\), whose length is strictly smaller than
\(2F_{n+1}\).  If \(\sigma\) had both a coordinate equal to \(1\) and a
coordinate equal to \(-1\), two of the values \(V_j\) would differ by
\(2F_{n+1}\), a contradiction.  Thus the nonzero coordinates of \(\sigma\)
have one sign, which is precisely \(\sigma\in\mathcal A_q\).

Let \(c\in\mathcal T_q\).  A positive-weight path from the initial state reaches
\(c\), and a positive-weight path from \(c\) reaches some reachable accepting
state \((\sigma,0)\).  Since
\(-\sigma\in\mathcal A_q\), one further column gives
\[
  \tau_{-\sigma}(\sigma,0)=(0,0).
\]
Hence every state of \(\mathcal T_q\) reaches the initial state, while the
initial state reaches every state of \(\mathcal T_q\) by the definition of the
trim set.  The directed graph of \(P_q\) is therefore strongly connected.  The
zero column has weight \(w_q(0)=2\) and gives a loop at \((0,0)\), so this graph
has period one.  Thus \(P_q\) is primitive.

Every path counted by \(S_q^{\mathrm{int}}(m)\) stays in the trim set, and the
initial and accepting vectors have nonzero restrictions to that set.  The
Perron--Frobenius theorem applied to \(P_q\) therefore gives the displayed
asymptotic with a positive leading coefficient and a strictly smaller
subdominant exponential rate.  Taking logarithms proves the pressure formula.
For \(q=0\), use
\(S_0^{\mathrm{int}}(m)=F_{m+2}\) and Binet's formula.
\end{proof}

\begin{lemma}[Initial second-moment enumeration]
\label{lem:initial-second-moment-enumeration}
For the counting function \(C_2(n)\) of
\Cref{thm:intrinsic-second-moment-recurrence}, one has
\[
  C_2(1)=2,
  \qquad C_2(2)=6,
  \qquad C_2(3)=14 .
\]
\end{lemma}

\begin{proof}
For a binary word of length \(n\), write
\(V_n(\omega)=\sum_{k=1}^n\omega_kF_k\), read modulo \(F_{n+1}\).  Then
\(C_2(n)\) is the sum of the squares of the residue-class multiplicities of
\(V_n\).  For \(n=1\), the residues modulo \(F_2=2\) have multiplicities
\((1,1)\), so \(C_2(1)=1^2+1^2=2\).  For \(n=2\), the four subset sums of
\(\{F_1,F_2\}=\{1,2\}\), modulo \(F_3=3\), have multiplicities
\((2,1,1)\), hence \(C_2(2)=2^2+1^2+1^2=6\).  For \(n=3\), the eight subset sums of
\(\{F_1,F_2,F_3\}=\{1,2,3\}\), modulo \(F_4=5\), have multiplicities
\[
  (2,2,1,2,1),
\]
because the residues are \(0,1,2,3,3,4,0,1\).  Therefore
\(C_2(3)=2^2+2^2+1^2+2^2+1^2=14\).
\end{proof}

\begin{theorem}[Sharp intrinsic second-moment recurrence]
\label{thm:intrinsic-second-moment-recurrence}
Let \(C_2(n)\) be the number of ordered pairs of binary words of length \(n\)
whose signed Fibonacci difference is congruent to zero modulo \(F_{n+1}\).  Then
\[
  C_2(1)=2,
  \qquad C_2(2)=6,
  \qquad C_2(3)=14,
\]
and, for \(n\ge4\),
\[
  C_2(n)=2C_2(n-1)+2C_2(n-2)-2C_2(n-3).
\]
Equivalently, since \(S_2^{\mathrm{int}}(m)=C_2(m+1)\),
\[
  S_2^{\mathrm{int}}(1)=6,
  \quad S_2^{\mathrm{int}}(2)=14,
  \quad S_2^{\mathrm{int}}(3)=36,
\]
and the same recurrence holds for \(S_2^{\mathrm{int}}(m)\) for \(m\ge4\).  The
first five intrinsic second moments are therefore
\[
  \bigl(S_2^{\mathrm{int}}(1),S_2^{\mathrm{int}}(2),S_2^{\mathrm{int}}(3),
  S_2^{\mathrm{int}}(4),S_2^{\mathrm{int}}(5)\bigr)
  =(6,14,36,88,220).
\]
The
ordinary generating functions are
\[
  \sum_{n\ge1}C_2(n)z^n
  =\frac{2z+2z^2-2z^3}{1-2z-2z^2+2z^3}
\]
and
\[
  \sum_{m\ge1}S_2^{\mathrm{int}}(m)z^m
  =\frac{6z+2z^2-4z^3}{1-2z-2z^2+2z^3}.
\]
\end{theorem}

\begin{proof}
For \(q=2\), the signed column alphabet is \(\{-1,0,1\}\) with weights
\(w(-1)=1,w(0)=2,w(1)=1\).  In the bounded carry automaton, let
\[
  A_n=W_n(0,0),\qquad B_n=W_n(1,0)=W_n(-1,0),
\]
where \(W_n(p,r)\) is the total weight of length-\(n\) paths from \((0,0)\) to
\((p,r)\), and put \(D_n=W_n(0,1)=W_n(0,-1)\).  The accepting targets are
\((-1,0),(0,0),(1,0)\), so \(C_2(n)=A_n+2B_n\).  The transition rule gives
\[
  A_{n+1}=2A_n+2B_n,
  \qquad B_{n+1}=D_n,
  \qquad D_{n+1}=A_n+2B_n=C_2(n).
\]
Hence \(B_{n+1}=C_2(n-1)\) for \(n\ge2\), and
\[
\begin{aligned}
  C_2(n+1)
    &=A_{n+1}+2B_{n+1} \\
    &=2A_n+2B_n+2C_2(n-1) \\
    &=2C_2(n)+2C_2(n-1)-2C_2(n-2).
\end{aligned}
\]
The initial values are the explicit length-\(1,2,3\) residue enumerations of
\Cref{lem:initial-second-moment-enumeration}.
The recurrence then gives \(S_2^{\mathrm{int}}(4)=88\) and
\(S_2^{\mathrm{int}}(5)=220\).
The generating functions are the standard translations of the recurrence and of
\(S_2^{\mathrm{int}}(m)=C_2(m+1)\).
\end{proof}

\begin{corollary}[Polynomial intrinsic fibre statistics are C-finite]
\label{cor:polynomial-intrinsic-fiber-statistics}
\label{cor:intrinsic-polynomial-statistics-cfinite}
Let \(P(t)\in\mathbb Z[t]\) be fixed.  Then
\[
  m\longmapsto \sum_{x\in X_m}P(|\Fold_m^{-1}(x)|)
\]
has a rational ordinary generating function and satisfies a constant-coefficient
linear recurrence.
\end{corollary}

\begin{proof}
Write \(P(t)=\sum_{q=0}^Q c_qt^q\).  The displayed statistic is
\(\sum_{q=0}^Q c_qS_q^{\mathrm{int}}(m)\).  By
\Cref{thm:intrinsic-fold-moment-transfer}, each fixed moment sequence is
C-finite, and finite linear combinations of C-finite sequences are C-finite.
\end{proof}

\begin{theorem}[Intrinsic-versus-cover transfer dichotomy]
\label{thm:intrinsic-cover-transfer-dichotomy}
\label{thm:intrinsic-versus-cover-transfer-dichotomy}
The unmodified Zeckendorf fold has fixed finite-state transfer in every fixed
moment degree.  In contrast, fixed finite-fibre certificates over the same
normal-form sets do not force such transfer for arbitrary computable covers:
for every computable sequence \(a_m\ge1\) there is a computable finite-fibre cover
preserving \(X_m\) and all normal-form readouts whose second covered moment is
\[
  S_2^\sharp(m)=(a_m+1)^2+F_{m+2}-1 .
\]
Thus intrinsic Fold\(_m\) moments are finite-state degree by degree, whereas
finite-state transfer is not invariant under arbitrary computable finite-fibre
covers.
\end{theorem}

\begin{proof}
The positive intrinsic assertion is \Cref{thm:intrinsic-fold-moment-transfer}.
For the cover assertion, choose the zero normal form in each \(X_m\), attach a
cover fibre of size \(a_m+1\) over it, and attach singleton cover fibres over
every other element of \(X_m\).  This changes only auxiliary certificate
multiplicities over already fixed normal forms, so it preserves all readouts
defined on those sets.  Since \(|X_m|=F_{m+2}\), the second covered moment is the
displayed quantity.  Taking \(a_m=m!\) gives a superexponential covered moment,
while every finite-state transfer sequence is exponentially bounded by
\Cref{prop:finite-state-coefficient-budget}.  Hence arbitrary covers can destroy
transfer although the intrinsic fold moments remain finite-state.
\end{proof}


% End original section: sec04_certificate_universality.tex
% Original section: sec04_finite_certificates.tex
\begin{proposition}[Classical finite-atomic certificate skeleton]
\label{prop:classical-finite-atomic-skeleton}
Let \(D\) be a finite multiset of positive integers and set
\(M_q(D)=\sum_{d\in D}d^q\).  Then the Hankel matrices
\([M_{i+j}(D)]_{0\le i,j\le r}\) are positive semidefinite for every \(r\), and
the moment sequence is log-convex:
\[
  M_q(D)^2\le M_{q-1}(D)M_{q+1}(D)\qquad(q\ge1).
\]
Moreover,
the moment segment \(M_0(D),\ldots,M_{2s-1}(D)\), where \(s\) is the number of
distinct values in \(D\), determines \(D\) uniquely.  Consequently fixed-fibre
certificate data determine exactly a positive finite atomic measure at each
resolution, supported on the distinct positive atoms \(a_1,\ldots,a_s\) with
positive integer weights \(n_1,\ldots,n_s\), and impose no recurrence or
finite-state transfer law as the resolution varies.
\end{proposition}

\begin{proof}
Write the distinct values as \(a_1<\cdots<a_s\), with positive integer
multiplicities \(n_1,\ldots,n_s\).  Thus the measure
\(\mu_D=\sum_i n_i\delta_{a_i}\) has finite support, positive weights, and
distinct atoms, and \(M_q(D)=\sum_i n_i a_i^q\).  For
\(Q(t)=\sum_{j=0}^r c_jt^j\),
\[
  \sum_{i,j=0}^r c_i c_j M_{i+j}(D)=\sum_{a\in D} Q(a)^2\ge0,
\]
which proves Hankel positivity; this is the finite positive-measure instance of
the moment problem \cite{Hamburger1920,ShohatTamarkin1943,Akhiezer1965}.  For
log-convexity, apply Cauchy's inequality to the two vectors
\((\sqrt n_i a_i^{(q-1)/2})_i\) and \((\sqrt n_i a_i^{(q+1)/2})_i\); their inner
product is \(M_q(D)\), and their squared norms are \(M_{q-1}(D)\) and
\(M_{q+1}(D)\).  For
reconstruction, the distinct-atom assumption makes the Hankel block
\([M_{i+j}]_{0\le i,j<s}\) nonsingular: it is the product of a Vandermonde
matrix, the positive diagonal matrix \(\operatorname{diag}(n_i)\), and the
transpose Vandermonde matrix.  The sequence \((M_q)\) therefore satisfies the
order-\(s\) recurrence whose characteristic polynomial is
\(\prod_i(T-a_i)\), and the first \(2s\) moments determine that recurrence by
the usual finite Hankel/Prony linear system; its roots are the \(a_i\), and the
Vandermonde system
\(M_q=
\sum_i n_i a_i^q\) for \(0\le q<s\) gives the multiplicities
\cite{Prony1795,BatenkovYomdin2013}.  All arguments are internal to the fixed
finite multiset \(D\), so they do not relate different resolutions.
\end{proof}

\begin{corollary}[Fixed-fibre Hankel positivity]
\label{thm:hankel-psd}
Let \(\mathrm{NF}_m\colon\Omega_m\to N_m\) be a finite surjection, and let
\(D_m=\{|\mathrm{NF}_m^{-1}(x)|:x\in N_m\}\) be its fibre-size
multiset.  For every fixed \(m\), if
\[
  S_q^{(m)}=\sum_{d\in D_m}d^q \qquad(q\ge0),
\]
then each finite Hankel matrix
\([S_{i+j}^{(m)}]_{0\le i,j\le r}\) is positive semidefinite, and
\[
  (S_q^{(m)})^2\le S_{q-1}^{(m)}S_{q+1}^{(m)}\qquad(q\ge1).
\]
In particular, for the intrinsic Zeckendorf fold over \(X_m\), the same assertions
hold for the collision moments \(S_q^{\mathrm{int}}(m)\).
\end{corollary}

\begin{proof}
Apply \Cref{prop:classical-finite-atomic-skeleton} to the finite positive-integer
multiset \(D_m\).  In the intrinsic Zeckendorf case, \(D_m\) is the multiset of
cardinalities of the fibres of \(\Fold_m\), and
\(S_q^{\mathrm{int}}(m)=\sum_{d\in D_m}d^q\) by definition.  The argument is
entirely internal to the chosen resolution \(m\).
\end{proof}

\begin{proposition}[Genericity of fixed-fibre Hankel certificates]
\label{rem:hankel-psd-generic}
The Hankel positivity and log-convexity assertions in \Cref{thm:hankel-psd}
are exactly generic finite-fibre consequences.  More precisely, a sequence
\((s_q)_{q\ge0}\) occurs as the moment sequence of the fibre-size multiset of a
finite surjection if and only if there are distinct positive integers
\(a_1,\ldots,a_t\) and positive integers \(n_1,\ldots,n_t\) such that
\[
  s_q=\sum_{i=1}^t n_i a_i^q\qquad(q\ge0).
\]
For every such sequence the finite Hankel matrices are positive semidefinite and
the sequence is log-convex.  Conversely, this description imposes no relation
between different resolutions: every computable sequence of such finite atomic
data can be realized by a computable sequence of finite surjections.
\end{proposition}

\begin{proof}
Let \(\mathrm{NF}\colon\Omega\to N\) be a finite surjection and let
\(D=\{|\mathrm{NF}^{-1}(x)|:x\in N\}\).  Group equal fibre sizes in \(D\).  If
the distinct sizes are \(a_1,
\ldots,a_t\) and the number of fibres of size \(a_i\) is \(n_i\), then
\(n_i\) and \(a_i\) are positive integers and the collision moments satisfy
\(s_q=\sum_i n_i a_i^q\).  \Cref{prop:classical-finite-atomic-skeleton} then
gives Hankel positivity and log-convexity.

Conversely, suppose distinct positive integers \(a_i\) and positive integers
\(n_i\) are given.  Take
\[
  N=\{(i,j):1\le i\le t,
  1\le j\le n_i\}
\]
and, for each \((i,j)\in N\), attach a disjoint set
\(\Omega_{i,j}\) of cardinality \(a_i\).  Let
\(\Omega=\bigsqcup_{(i,j)\in N}\Omega_{i,j}\), and define
\(\mathrm{NF}\) to send every point of \(\Omega_{i,j}\) to \((i,j)\).  This is
a finite surjection whose fibre-size multiset contains \(a_i\) exactly \(n_i\)
times, hence has moments \(s_q\).  Applying the same construction independently
at each resolution realizes any computable sequence of finite atomic data; if
the lists \((a_i^{(m)},n_i^{(m)})_i\) are computable, the corresponding finite
sets and maps are computable.  No step of the construction identifies data at
resolution \(m\) with data at any other resolution.
\end{proof}

\begin{corollary}[Finite Prony reconstruction of fibre-size data]
\label{thm:prony}
With the notation of \Cref{thm:hankel-psd}, suppose that \(D_m\) has \(s_m\)
distinct fibre sizes.  Then the finite moment segment
\[
  S_0^{(m)},S_1^{(m)},\ldots,S_{2s_m-1}^{(m)}
\]
determines the fibre-size multiset \(D_m\) uniquely.  Consequently the finite
segment of length at most \(2|N_m|-1\) always reconstructs \(D_m\).  For the
Zeckendorf fold this gives the explicit reconstruction bound
\(2|X_m|-1=2F_{m+2}-1\).
\end{corollary}

\begin{proof}
The first assertion is exactly the Prony part of
\Cref{prop:classical-finite-atomic-skeleton}, applied to \(D_m\): the first
\(2s_m\) moments determine the distinct support values and their multiplicities.
Since \(s_m\le |D_m|=|N_m|\), the initial segment through index \(2|N_m|-1\)
contains the required segment.  In the Zeckendorf fold, \(N_m=X_m\), and the
standard count of Zeckendorf words of length at most \(m\) gives
\(|X_m|=F_{m+2}\).
\end{proof}

\begin{theorem}[Finite-fibre certificates are not finite-state transfer]
\label{thm:hankel-prony-transfer-independence}
\label{thm:finite-fibre-finite-state-separation}
The fixed-resolution Hankel--Prony certificates of
\Cref{prop:classical-finite-atomic-skeleton,thm:hankel-psd,thm:prony} are exactly
finite-atomic fibre facts and do not imply any fixed finite-state transfer law in
the resolution variable.  More precisely, for every computable sequence
\((a_m)_{m\ge1}\) of positive integers there is a computable sequence of finite
surjections with one normal form at each resolution whose first collision moment
is \(S_1^{(m)}=a_m\).  Hence choosing \(a_m=m!\) gives effective finite-fibre
Hankel--Prony certificates at every resolution but no fixed finite-state transfer
formula for \((S_1^{(m)})_m\).  The same construction may be regarded as a
computable finite-fibre cover over singleton normal-form sets, so the obstruction
is already present before any Zeckendorf-specific structure is added.
\end{theorem}

\begin{proof}
The first sentence is the content of
\Cref{prop:classical-finite-atomic-skeleton}: at a fixed resolution the fibre-size
multiset is a positive finite atomic measure, Hankel positivity is positivity of
squares against that measure, and Prony reconstruction recovers only that finite
multiset from enough moments.  None of these operations compares two different
resolutions.

For a prescribed computable sequence \((a_m)\), take
\(N_m=\{*_m\}\), take \(\Omega_m=\{1,\ldots,a_m\}\), and let
\(\mathrm{NF}_m\) be the unique map to \(N_m\).  The fibre-size multiset is
\(D_m=\{a_m\}\), so the fixed-resolution conclusions of
\Cref{thm:hankel-psd,thm:prony} hold for every \(m\), and
\(S_1^{(m)}=a_m\).  The construction is computable because the sets and the
unique maps are effectively enumerable from \(m\) and \(a_m\).  If
\((m!)_m\) had a fixed finite-state transfer formula, then
\Cref{prop:finite-state-coefficient-budget} would make it exponentially bounded:
\(m!\le C\rho^m\) for some constants \(C,\rho>0\).  This is impossible,
since \(m!/\rho^m\to\infty\).  Thus finite-fibre Hankel positivity and
finite Prony reconstruction are classical fixed-resolution support facts only;
the uniform matrix representation of \Cref{def:fixed-finite-state-transfer} is
an additional cross-resolution hypothesis.
\end{proof}

The preceding Hankel and Prony statements are fixed-resolution support for
\(L2\).  \Cref{rem:hankel-psd-generic} identifies them as generic finite-fibre
facts.  They are classical support rather than the \(L2\) novelty claim of this article and, by
\Cref{thm:hankel-prony-transfer-independence}, they impose no recurrence in the
resolution variable.  A kernel-compression restatement of the same finite-atomic
calculation is relegated to the appendix.

\begin{definition}[Effective resolution data]\label{def:effective-resolution-data}
An effective finite-resolution datum is a computable sequence of finite
surjections \(\mathrm{NF}_m\colon\Omega_m\to N_m\) with finite fibres and
computable fibre-size lists.
\end{definition}

\begin{definition}[Certificate universality]\label{def:cert-univ}
An \(L1\) protocol is \(L2\)-certificate-universal if it is equipped with
effective finite-resolution data whose fibre-size multisets provide the fixed
finite certificates of \Cref{thm:hankel-psd,thm:prony}.  Fixed
finite-state transfer, in the sense of \Cref{def:fixed-finite-state-transfer},
is a stronger extra condition: for each fixed \(q\), the sequence
\((S_q^{(m)})_m\) must have one finite matrix model independent of the resolution
\(m\).
\end{definition}

\begin{proposition}[Finite-state coefficient budget]
\label{prop:finite-state-coefficient-budget}
If a sequence \((b_m)_{m\ge1}\) has a fixed finite-state transfer formula in the
sense of \Cref{def:fixed-finite-state-transfer}, equivalently
\(b_m=u^{\mathsf T}A^m v\), then \((b_m)\) satisfies a linear recurrence with
constant coefficients over the same field, and consequently its generating
series is rational over that field.  If, in addition, the coefficient field is
equipped with a fixed embedding in \(\mathbb C\), then there are constants
\(C,\rho\) such that \(|b_m|\le C\rho^m\) for the induced complex absolute value.
In particular this exponential budget applies to all integer-valued and
real-valued collision-moment sequences used below.
\end{proposition}

\begin{proof}
The Cayley--Hamilton theorem gives a linear recurrence of order at most the
size of \(A\) for the powers \(A^m\), hence for \(u^{\mathsf T}A^m v\).  A
sequence satisfying a constant-coefficient recurrence has rational generating
series over the coefficient field.  After fixing an embedding in \(\mathbb C\),
choose any complex matrix norm.  The exponential bound follows from
submultiplicativity: choose \(\rho>\|A\|\) and absorb the scalar factors
\(\|u\|\|v\|\), together with the finitely many initial terms, into \(C\).
\end{proof}

\begin{proposition}[Fixed-transfer interface for the budget obstruction]
\label{prop:fixed-transfer-interface-budget}
Every use of the finite-state transfer budget in this article refers to
\Cref{def:fixed-finite-state-transfer}.  In particular, a degree-\(q\)
collision-moment sequence with fixed finite-state transfer is subject to the
recurrence and rational-generating-series constraints of
\Cref{prop:finite-state-coefficient-budget}.  For the integer-valued and
real-valued collision moments considered below, it is also subject to the
exponential-growth constraint there; therefore any computable cover whose
degree-\(q\) moment violates every exponential bound has no fixed finite-state
transfer formula in that degree.
\end{proposition}

\begin{proof}
The definition gives exactly one finite matrix model
\(S_q^{(m)}=u_q^{\mathsf T}A_q^m v_q\) for all resolutions in the tail under
consideration.  Applying \Cref{prop:finite-state-coefficient-budget} to the
sequence \(b_m=S_q^{(m)}\) gives a constant-coefficient recurrence and a rational
generating series.  In the collision-moment applications below the sequence is
integer-valued, hence viewed inside \(\mathbb C\), so the same proposition gives
constants \(C,\rho\) with \(|S_q^{(m)}|\le C\rho^m\).  Hence a sequence which
eventually exceeds \(C\rho^m\) for every pair of constants \(C,\rho\) cannot
admit the fixed matrix representation in \Cref{def:fixed-finite-state-transfer}.
This is precisely the budget obstruction used later for the one-spike covers.
\end{proof}

\begin{proposition}[Finite-fibre certificate taxonomy]
\label{prop:finite-fiber-certificate-taxonomy}
For any effective finite-resolution datum, the finite-fibre certificate at
resolution \(m\) factors through the multiset of fibre sizes.  Thus every
fixed-resolution moment certificate recorded here is invariant under replacing
the datum by any other datum with the same fibre-size multiset.
\end{proposition}

\begin{proof}
Each listed invariant is a function only of the moments \(\sum d^q\) of the
finite fibre-size multiset.  The conclusions are precisely the finite-atomic
certificate skeleton of \Cref{prop:classical-finite-atomic-skeleton}.
\end{proof}

\begin{proposition}[Zeckendorf certificate universality]
\label{prop:zeckendorf-cert-univ}
The Zeckendorf \(L1\) protocol of \Cref{thm:zeckendorf-protocol-univ} is
\(L2\)-certificate-universal: its finite data are the computable surjections
\(\Fold_m\colon\Omega_m\to X_m\), and their fibre-size multisets give the
certificates of \Cref{prop:classical-finite-atomic-skeleton}.  The finite Prony
reconstruction bound is computable and may be taken to be
\(2|X_m|-1=2F_{m+2}-1\) at resolution \(m\).
\end{proposition}

\begin{proof}
For fixed \(m\), both \(\Omega_m\) and \(X_m\) are finite and effectively
enumerable, and \Cref{lem:greedy-zeckendorf-normalizer} computes the fold map.
Counting preimages gives the fibre-size list; the certificate assertions then
follow from
\Cref{prop:classical-finite-atomic-skeleton,rem:hankel-psd-generic}.  The
explicit Prony bound follows from \(|X_m|=F_{m+2}\), since the number of distinct
fibre sizes is at most the number of normal forms.
\end{proof}

\begin{theorem}[Consolidated \texorpdfstring{\(L2\)}{L2} interface]
\label{thm:L2-consolidated-interface}
For the Zeckendorf \(L1\) protocol, the \(L2\) layer used in
\Cref{thm:main-hierarchy} is exactly the following four-part interface.
\begin{enumerate}
\item At each fixed resolution, the fibre-size multiset gives the finite-atomic
  Hankel positivity, log-convexity, and Prony reconstruction certificates of
  \Cref{prop:classical-finite-atomic-skeleton,thm:hankel-psd,rem:hankel-psd-generic,thm:prony}.
\item For the unmodified maps \(\Fold_m\), every fixed finite intrinsic moment
  degree has fixed finite-state transfer through the audited bounded-carry
  interface of \Cref{thm:audited-fixed-degree-transfer-interface}.  Its trimmed
  matrix is primitive, so the fixed-degree pressure exists; in particular the intrinsic second
  moment satisfies the recurrence of \Cref{thm:intrinsic-second-moment-recurrence},
  polynomial intrinsic fibre statistics are C-finite, and the audited shifted
  half-interval conversion of \Cref{thm:audited-shifted-half-interval-height-conversion} gives the intrinsic height
  \(M_{2s-1}=F_{s+1}\), \(M_{2s}=2F_s\).
\item The full intrinsic moment tower has no single rational finite-state
  transfer law, because its bivariate moment series is nonrational by
  \Cref{thm:intrinsic-all-degree-transfer-obstruction}.
\item In the computable finite-fibre cover category over the same normal-form
  sets \(X_m\), the certificates in the first item do not force fixed finite-state
  transfer: the extremal one-spike criterion of
  \Cref{thm:L2-one-spike-separation-criterion} gives a computable cover
  preserving all \(L1\) readouts whose second cover moment has no fixed
  finite-state transfer formula, and there are two covers over the same intrinsic
  data \((X_m,\Fold_m)\) with opposite finite-state-transfer behaviour.
\end{enumerate}
Thus \(L2\) has an intrinsic separation between degreewise and a
rational-function all-degree transfer law in the moment marker.  The stronger
assertion that fixed-resolution certificates do not force transfer even in
degree two is the separate cover-relative obstruction.
\end{theorem}

\begin{proof}
For the first item, fix a resolution \(m\) and let
\(D_m=\{|\Fold_m^{-1}(x)|:x\in X_m\}\).  The sets \(\Omega_m\) and \(X_m\) are
finite and effectively enumerable, and \Cref{lem:greedy-zeckendorf-normalizer}
computes \(\Fold_m\).  Hence \(D_m\) is a computable finite multiset of positive
integers.  Applying \Cref{prop:classical-finite-atomic-skeleton} to \(D_m\) gives
the Hankel positivity and log-convexity statements, and the Prony part of the
same proposition reconstructs the multiset from the first \(2s_m\) moments, where
\(s_m\) is the number of distinct fibre sizes.  Since \(s_m\le |X_m|=F_{m+2}\),
the bound in \Cref{thm:prony} follows.  This is precisely the finite-resolution
certificate universality of \Cref{prop:zeckendorf-cert-univ}.

For the second item, \Cref{thm:intrinsic-fold-moment-transfer} supplies a fixed
finite-state transfer formula for each fixed intrinsic moment degree.  Applying
\Cref{prop:finite-state-coefficient-budget} gives the associated
constant-coefficient recurrence, rational generating series, and exponential
budget.  The degree-two recurrence and polynomial-statistic consequences are
\Cref{thm:intrinsic-second-moment-recurrence,cor:polynomial-intrinsic-fiber-statistics}.
The primitive asymptotic and pressure formula are
\Cref{thm:primitive-intrinsic-moment-transfer}.
The height dependency firewall is
\Cref{thm:ordinary-maximum-supports-intrinsic-height}; the height formula is
\Cref{thm:intrinsic-fold-fibre-height}, with the entropy and large-moment
endpoints recorded in
\Cref{cor:linfty-fold-entropy,thm:zero-temperature-fold-pressure,thm:diagonal-zero-temperature-fold-endpoint}.

The third item is
\Cref{thm:intrinsic-all-degree-transfer-obstruction}.  Its diagonal
superexponential coefficient growth rules out a rational bivariate moment
series and hence a single rational transfer matrix for all moment degrees.

For the fourth item, \Cref{thm:L2-one-spike-separation-criterion} identifies the
one-spike cover as the extremal cover shape for a prescribed excess and realizes
the factorial excess \(m!\).  The cover changes only auxiliary fibre sizes over
the unchanged normal-form sets, and \Cref{lem:cover-preserves-protocol} shows
that all normal forms and \(L1\) readouts are preserved.  Its second cover moment
is \((m!+1)^2+F_{m+2}-1\), which grows faster than \(C\rho^m\) for every fixed
pair \(C,\rho\).  By \Cref{prop:fixed-transfer-interface-budget}, no fixed
finite-state transfer formula can produce such a sequence.  The same skeleton
also gives the two-cover boundary of \Cref{thm:L2-two-cover-intrinsic-boundary}:
the singleton cover has the Fibonacci transfer sequence \(F_{m+2}\), while the
factorial one-spike cover over the same \((X_m,\Fold_m)\) violates the
exponential budget.

The first three items concern the intrinsic fibres of the original maps
\(\Fold_m\); the fourth concerns auxiliary cover multiplicities over the already
fixed normal forms.  Since \Cref{thm:L2-cover-intrinsic-nonimplication} proves
that cover-moment behaviour is not determined by the intrinsic fold data without
an additional bridge, the degree-two cover statement remains exactly
cover-relative; it is distinct from the intrinsic all-degree obstruction.
\end{proof}

\begin{theorem}[Extremal one-spike classification for \texorpdfstring{\(L2\)}{L2} separation]
\label{prop:L2-novelty-finite-state-separation}
\label{thm:L2-one-spike-separation-criterion}
Fix the Zeckendorf normal-form sets \(X_m\).  Let a computable finite-fibre cover
have cover fibre sizes \(d_m^\sharp(x)\ge1\) over \(x\in X_m\), and put
\[
  E_m=\sum_{x\in X_m}\bigl(d_m^\sharp(x)-1\bigr),
  \qquad
  S_2^\sharp(m)=\sum_{x\in X_m}d_m^\sharp(x)^2 .
\]
Let
\[
  H_m=\max_{x\in X_m}\bigl(d_m^\sharp(x)-1\bigr)
\]
be the spike height of the cover at resolution \(m\).  Then, for each fixed
resolution \(m\),
\[
  |X_m|+2E_m\le S_2^\sharp(m)
  \le |X_m|+2E_m+E_m^2,
\]
and the sharper integer lower bound
\[
  S_2^\sharp(m)
  \ge |X_m|+3E_m
\]
is attained exactly by distributing the excess over distinct normal forms, so
necessarily \(E_m\le |X_m|\).  If \(E_m>0\), the upper bound is attained exactly
when all excess lies on one normal form, that is, by a one-spike cover.  More
generally, if the excess is supported
on \(t_m\) normal forms, then
\[
  S_2^\sharp(m)
  \le |X_m|+2E_m+(E_m-t_m+1)^2+t_m-1,
\]
with equality exactly when one support point carries excess \(E_m-t_m+1\) and the
remaining \(t_m-1\) support points carry excess \(1\).  Consequently every cover
whose second moment violates all exponential bounds must have spike height which
violates all exponential bounds.  Conversely, if the spike-height sequence
\((H_m)_m\) violates all exponential bounds, then the second covered moment does
too.  Thus the finite-state-transfer budget obstruction in the computable-cover
category is exactly a spike-height obstruction: distributed low-height cover
excess cannot witness it, while a computable one-spike cover realizes the worst
case.  In particular, a one-spike cover with \(E_m=m!\) has no fixed finite-state
transfer formula for its second covered moment, although all fixed-resolution
finite-atomic certificates of
\Cref{prop:classical-finite-atomic-skeleton} exist at every \(m\) and the
unmodified intrinsic moments still have the fixed transfer package of
\Cref{thm:intrinsic-fold-moment-transfer}.  Thus the genuine \(L2\) novelty is
the cross-resolution finite-state separation for covers, not the classical
Hankel--Prony certificate at one resolution.
\end{theorem}

\begin{proof}
Write \(e_x=d_m^\sharp(x)-1\).  Then \(e_x\ge0\) and \(
\sum_{x\in X_m}e_x=E_m\).  Expanding the second moment gives
\[
  S_2^\sharp(m)=\sum_{x\in X_m}(1+e_x)^2
  =|X_m|+2E_m+
  \sum_{x\in X_m}e_x^2 .
\]
Since \(e_x\ge0\), one has \(0\le\sum_x e_x^2\le(\sum_x e_x)^2=E_m^2\), giving the
first displayed two-sided estimate.  Since the \(e_x\) are integers, also
\(e_x^2\ge e_x\) for every \(x\), and therefore
\(\sum_x e_x^2\ge E_m\).  Equality in this sharper lower estimate is equivalent
to \(e_x^2=e_x\) for every \(x\), hence to \(e_x\in\{0,1\}\) for every \(x\); this
is precisely distribution of the excess over \(E_m\) distinct normal forms and
therefore requires \(E_m\le |X_m|\).  When \(E_m>0\), equality in the upper
estimate occurs exactly when exactly one \(e_x\) is nonzero, which is precisely
the one-spike cover shape.

Now suppose the excess support has size \(t_m\ge1\).  On that support write the
positive excesses in nonincreasing order
\(b_1\ge\cdots\ge b_{t_m}\ge1\), with \(\sum_i b_i=E_m\).  Then
\[
  \sum_{i=1}^{t_m} b_i^2
  =E_m^2-2\sum_{1\le i<j\le t_m}b_ib_j
  \le E_m^2-2\sum_{j=2}^{t_m}b_1b_j
\]
is not the sharp estimate because \(b_1\) is itself variable.  A direct
compression argument gives the sharp form: if two support values satisfy
\(b_i,b_j>1\), replacing them by \(b_i+b_j-1\) and \(1\) preserves the total
excess and increases the sum of squares by
\[
  (b_i+b_j-1)^2+1-b_i^2-b_j^2
  =2(b_i-1)(b_j-1)>0 .
\]
Iterating this operation shows that the maximum at fixed support size \(t_m\) is
attained exactly at the partition
\((E_m-t_m+1,1,\ldots,1)\), and its square sum is
\((E_m-t_m+1)^2+t_m-1\).  Substitution into the expansion of
\(S_2^\sharp(m)\) gives the displayed support-size bound and its equality
criterion.

Since \(e_x^2\le H_m e_x\), every cover satisfies
\[
  S_2^\sharp(m)\le |X_m|+2E_m+H_mE_m .
\]
Because each fibre point contributes to exactly one normal form, also
\(E_m\le H_m|X_m|\).  As \(|X_m|=F_{m+2}\) is exponentially bounded, an
exponentially bounded spike height \(H_m\) would make both \(E_m\) and
\(S_2^\sharp(m)\) exponentially bounded.  Explicitly, if
\(H_m\le C\rho^m\) and \(|X_m|\le D\sigma^m\), then
\[
  E_m\le CD(\rho\sigma)^m,
  \qquad
  H_mE_m\le C^2D(\rho^2\sigma)^m,
\]
and the preceding estimate bounds \(S_2^\sharp(m)\) exponentially.  This proves
that no low-height distributed cover can violate the finite-state exponential
budget.

Conversely, \(H_m^2\le\sum_x e_x^2\le S_2^\sharp(m)\).  Hence an exponential
bound \(S_2^\sharp(m)\le C\rho^m\) would imply
\(H_m\le C^{1/2}(\rho^{1/2})^m\).  Thus failure of every exponential bound for
\((H_m)_m\) forces failure of every exponential bound for the second covered
moment.  Combining the two directions gives the announced spike-height rigidity
classification of the budget obstruction.

For the one-spike cover with \(E_m=m!\), the moment is
\[
  S_2^\sharp(m)=(m!+1)^2+|X_m|-1=(m!+1)^2+F_{m+2}-1,
\]
using \(|X_m|=F_{m+2}\).  This sequence is not exponentially bounded.  Indeed,
for any fixed \(C,
\rho>0\), \((m!)^2/(C\rho^m)\to\infty\), so eventually
\(S_2^\sharp(m)>C\rho^m\).  By
\Cref{prop:finite-state-coefficient-budget}, every fixed finite-state transfer
sequence is exponentially bounded; hence this cover moment has no fixed
finite-state transfer formula.  At each resolution the fibre-size multiset is
finite, so the finite-atomic Hankel and Prony certificates hold by
\Cref{prop:classical-finite-atomic-skeleton}.  The construction alters only
auxiliary cover multiplicities over the already fixed normal forms; it does not
alter \(\Fold_m\), whose intrinsic fixed-degree moments have transfer by
\Cref{thm:intrinsic-fold-moment-transfer}.  This proves the asserted separation.
\end{proof}

\begin{theorem}[Finite certificate windows do not certify transfer]
\label{thm:L2-finite-window-transfer-obstruction}
Fix integers \(M,Q\ge0\).  In the computable finite-fibre cover category over the
Zeckendorf normal-form sets, there are two cover systems which have identical
collision moments
\[
  S_q^\sharp(m)
  \qquad(1\le m\le M,\;0\le q\le Q)
\]
and identical fixed-resolution finite-fibre certificate output on that window,
but only one of them has second collision moment compatible with a fixed
finite-state transfer formula.  More strongly, every predicate depending only on
the fibre-size multisets
\[
  \{d_m^\sharp(x):x\in X_m\}\qquad(1\le m\le M)
\]
assigns the same truth value to the two systems, while the second system has no
finite-state transfer formula for \(S_2^\sharp(m)\).  Hence no finite window of
fixed-resolution finite-fibre certificate data or finitely many collision
moments can certify finite-state transfer for all resolutions.
\end{theorem}

\begin{proof}
Let the first cover be the singleton cover over every normal form, so every
covered fibre has size \(1\).  Its second moment is \(|X_m|=F_{m+2}\).  This
sequence satisfies the Fibonacci recurrence, equivalently a transfer formula
with the companion matrix
\(\begin{pmatrix}1&1\\1&0\end{pmatrix}\) and suitable initial vectors, and
therefore satisfies the finite-state condition for \(q=2\).

For the second cover, use the one-spike construction over the same normal-form
sets.  Choose the all-zero normal form at each resolution.  Put spike size
\(1\) for every \(1\le m\le M\), and put spike size \(m!+1\) for every
\(m>M\); all other fibres have size \(1\).  This is a computable finite-fibre
cover and it preserves the normal forms and \(L1\) readouts, as in
\Cref{lem:cover-preserves-protocol}.  For every \(1\le m\le M\) its fibre-size
multiset is exactly the singleton multiset of the first cover.  Therefore every
predicate or numerical invariant depending only on those multisets in the window
has the same value for the two covers.  In particular, all displayed moments
agree for \(0\le q\le Q\) and \(1\le m\le M\).

For \(m>M\), the second collision moment of the spiked cover is
\[
  (m!+1)^2+|X_m|-1 .
\]
This sequence grows faster than every exponential sequence, whereas every fixed
finite-state transfer sequence is exponentially bounded by
\Cref{prop:finite-state-coefficient-budget}.  Thus the second cover has no fixed
finite-state transfer formula for its second moment, despite agreeing with the
finite-state singleton cover on the prescribed finite certificate window.

All fixed-resolution finite-fibre certificates used in this article factor
through the finite fibre-size multiset by
\Cref{prop:classical-finite-atomic-skeleton,rem:hankel-psd-generic,prop:finite-fiber-certificate-taxonomy}.
They are therefore included among the finite-window data just matched.  This
proves that the transfer condition is genuinely cross-resolution structure and
not a consequence of any fixed-window finite-fibre certificate facts.
\end{proof}


% End original section: sec04_finite_certificates.tex
% Original section: sec04_intrinsic_height.tex
\subsection{Intrinsic fibre height and the zero-temperature endpoint}
\label{ssec:intrinsic-height-zero-temperature}

For \(n\ge0\) and \(N\in\ZZ\), let
\[
  R_n(N)=\#\left\{A\subseteq\{1,\ldots,n\}:
  \sum_{i\in A}F_i=N\right\},
\]
with \(R_n(N)=0\) when no such subset exists.  Separating subsets according as
they contain \(n\) gives
\begin{equation}\label{eq:R-recursion}
  R_n(N)=R_{n-1}(N)+R_{n-1}(N-F_n)
  \qquad(n\ge1).
\end{equation}

\begin{lemma}[Fold fibre--representation pairing]
\label{lem:fold-fibre-representation-pairing}
For \(x\in X_m\), put \(r=\Val_m(x)\in\{0,\ldots,F_{m+2}-1\}\).  Then
\[
  d_m(x)=|\Fold_m^{-1}(x)|
  =R_{m+1}(r)+R_{m+1}(r+F_{m+2})
  =R_{m+2}(r+F_{m+2}).
\]
Thus intrinsic \(\Fold_m\)-fibres are exactly a shifted half-interval slice of
the truncated Fibonacci representation function.
\end{lemma}

\begin{proof}
By \eqref{eq:fold-residue}, a word \(\omega\in\Omega_m\) lies over \(x\) exactly
when
\[
  \sum_{k=1}^{m+1}\omega_kF_k\equiv r\pmod {F_{m+2}} .
\]
The possible sums lie in
\([0,\sum_{k=1}^{m+1}F_k]=[0,F_{m+3}-2]\), and
\(F_{m+3}-2<2F_{m+2}\).  Hence the only possible sums in this congruence class
are \(r\) and \(r+F_{m+2}\).  Counting by the subset
\(A=\{k:\omega_k=1\}\subseteq\{1,\ldots,m+1\}\) gives the first equality.
For the second equality, apply \eqref{eq:R-recursion} with
\(n=m+2\) and \(N=r+F_{m+2}\): representations using the top part
\(F_{m+2}\) are counted by \(R_{m+1}(r)\), and those omitting it are counted by
\(R_{m+1}(r+F_{m+2})\).
\end{proof}

\begin{lemma}[Positive-index conversion for the ordinary maximum theorem]
\label{lem:fibonacci-maximum-normalization-conversion}
\label{lem:positive-index-conversion-ordinary-maximum}
\label{thm:source-audited-ordinary-maximum-import}
Let \((\mathcal F_k)_{k\ge0}\) denote the Fibonacci sequence used in the 2005
public article of Koc\'abov\'a--Mas\'akov\'a--Pelantov\'a,
\emph{Integers with a maximal number of Fibonacci representations}, RAIRO
Theor. Inform. Appl. \textbf{39} (2005), 343--359, available at the stable DOI
\texttt{10.1051/ita:2005022} and the public Numdam landing page recorded in
\cite{KocabovaMasakovaPelantova2005}.  The Numdam record identifies the same
title, authors, journal, volume \(39\), issue \(2\), pages \(343\)--\(359\), year
\(2005\), DOI, and PDF URL as the bibliography entry.  The citation role is
limited to this one ordinary maximum theorem: the only external result imported
here is Theorem~4.7 of that article
\cite[Thm.~4.7]{KocabovaMasakovaPelantova2005}.
In that source
\(\mathcal F_0=\mathcal F_1=1\),
\(\mathcal F_{k+1}=\mathcal F_k+\mathcal F_{k-1}\), Fibonacci
representations are sums of distinct parts \(\mathcal F_i\) with \(i\ge1\), and
the displayed maximum is
\[
  \operatorname{Max}(k)=\max\{R(N):\mathcal F_k\,\le N<\mathcal F_{k+1}\}
\]
for positive integers in the half-open interval.  With this paper's convention
\(F_1=1,F_2=2,F_{k+2}=F_{k+1}+F_k\), one has
\(\mathcal F_k=F_k\) for every \(k\ge1\).  Consequently the source maximum is
exactly
\[
  C_k=\max\{R(N):F_k\le N<F_{k+1}\}
\]
used below, with no index shift for \(k\ge1\).
The bibliographic record for the present version fixes both the DOI and the
Numdam landing page of the RAIRO source.  The theorem number, DOI, URL, and
positive-index convention have been checked against that public source record;
this paragraph is citation support for the imported ordinary maximum formula,
not a reproof of that source theorem.  Theorem~4.7 of that source is used here
exactly as the ordinary half-open interval maximum theorem for the above index
convention; it is not used for any modular, folded, or cover-relative assertion.
Theorem~4.7 of the cited source gives, after this positive-index normalization
and interval conversion,
\[
  C_{2t+1}=F_{t+1}\quad(t\ge0),
  \qquad
  C_{2t+2}=2F_t\quad(t\ge1),
\]
with the remaining value \(C_2=1\) supplied by the direct computation
\(R(2)=1\).  No classification of maximizing arguments from the cited paper is
used.  This ordinary interval maximum formula is the only external maximum
formula imported for \Cref{thm:half-interval-fibonacci-representation-max}; in
particular it is cited support for the intrinsic height package, not an
independent maximum theorem proved in this manuscript.  This lemma records the
citation boundary, the source metadata, and the normalization needed later; its
only mathematical input not proved here is the ordinary interval maximum formula
from the cited source.
All later uses of \(C_k\) are therefore to be read as uses of this cited
background theorem, with the displayed positive-index convention already locked;
the conversion lemmas below prove only the passage from these ordinary maxima to
the shifted half-intervals and then to intrinsic fold fibres.
\end{lemma}

\begin{proof}
The source conventions just listed are the conventions printed in the
introduction of the 2005 article \cite{KocabovaMasakovaPelantova2005}: its representations have
the form \(N=\mathcal F_{m_r}+\cdots+\mathcal F_{m_1}\) with
\(m_r>\cdots>m_1\ge1\), its representation number is denoted by the same symbol
\(R(N)\), and its \(\operatorname{Max}(k)\) is the maximum of that count over
\(\mathcal F_k\le N<\mathcal F_{k+1}\).  The two Fibonacci normalizations agree at
indices \(1\) and \(2\): \(\mathcal F_1=1=F_1\) and
\(\mathcal F_2=\mathcal F_1+\mathcal F_0=2=F_2\).  The common second-order
recurrence gives \(\mathcal F_k=F_k\) for all positive indices by induction.
Thus the admissible parts, the representation-counting function, and the
half-open intervals coincide term-by-term with this paper's
\([F_k,F_{k+1})\), so \(\operatorname{Max}(k)=C_k\) for every \(k\ge1\).

Theorem~4.7 in the cited source states the two interval maxima as
\[
  \max\{R(N):\mathcal F_{2t+1}\le N<\mathcal F_{2t+2}\}=\mathcal F_{t+1}
  \quad(t\ge0)
\]
and
\[
  \max\{R(N):\mathcal F_{2t+2}\le N<\mathcal F_{2t+3}\}=2\mathcal F_t
  \quad(t\ge1).
\]
Replacing \(\mathcal F_j\) by \(F_j\) for positive indices gives exactly the two
displayed parity clauses for \(C_k\), including the source restrictions
\(t\ge0\) and \(t\ge1\).  The exceptional value \(C_2\) is not imported from a
parity clause:
the interval \([F_2,F_3)=[2,3)\) contains only \(2=F_2\), so \(C_2=R(2)=1\).
No quantity with index \(0\) in the source sequence is used as an admissible part
or as an endpoint in the later half-interval argument.  Thus the source year,
theorem number, and positive-indexing convention have all been fixed before the
formula is used in the intrinsic height argument.
Thus the cited theorem supplies exactly the two ordinary interval maxima after
the positive-index normalization above, while every subsequent passage from
ordinary intervals to shifted half-intervals and then to \(\Fold_m\)-fibre height
is proved below inside the paper.
\end{proof}

\begin{corollary}[Citation-sensitive half-interval conversion]
\label{cor:citation-sensitive-half-interval-conversion}
The exact values in \Cref{thm:half-interval-fibonacci-representation-max} are
conditional only on the source-audited ordinary interval formula of
\Cref{thm:source-audited-ordinary-maximum-import} plus the in-paper conversion
from ordinary intervals to the shifted half-intervals defining \(H_n\).  No
ordinary-interval endpoint convention, source index \(0\), or classification of
maximizing integers is imported into the intrinsic fold-height theorem.  The
fold-aware shifted-half-interval conversion and the passage from those
half-intervals to intrinsic \(\Fold_m\)-fibre height are the in-paper content;
the citation supplies only the ordinary interval maximum values.  Equivalently,
the complete external interface is the three-line positive-index statement
\[
  C_{2t+1}=F_{t+1}\quad(t\ge0),
  \qquad C_{2t+2}=2F_t\quad(t\ge1),
  \qquad C_2=1,
\]
for \(C_k=\max\{R(N):F_k\le N<F_{k+1}\}\), with distinct parts drawn from
\(F_1,F_2,\ldots\).  Replacing the cited theorem by a self-contained proof of
exactly this positive-index ordinary-maximum statement would leave the
half-interval, intrinsic height, and large-moment endpoint conclusions unchanged.
\end{corollary}

\begin{proof}
The source-audited theorem identifies the external input as the two parity
formulas for \(C_k\) on intervals \([F_k,F_{k+1})\) with positive-index parts,
together with the directly checked exceptional value \(C_2=1\).  The proof of
\Cref{thm:half-interval-fibonacci-representation-max} uses those formulas only
after the elementary pairing of \Cref{lem:fold-fibre-representation-pairing}: the
shifted half-interval count is first reduced to the ordinary count on
\([F_{n+1},F_{n+2})\), and the remaining tail is bounded internally by
\eqref{eq:R-recursion} and the inequality \(C_n\le C_{n+1}\) read from the
displayed parity formulas.  Thus the conversion is a theorem of this manuscript
once the ordinary interval maximum formula is accepted, and the only
citation-sensitive premise is the audited Theorem~4.7 normalization recorded
above.

The source check is against the DOI and Numdam copy named in the bibliography:
Theorem~4.7 there states the two displayed parity formulas for
\(\operatorname{Max}(k)=\max\{R(n):\mathcal F_k\le n<\mathcal F_{k+1}\}\), with
\(\mathcal F_0=\mathcal F_1=1\) and distinct positive-index Fibonacci parts.  The
proof of \Cref{thm:source-audited-ordinary-maximum-import} gives the complete
parity and positive-index conversion needed to replace \(\mathcal F_k\) by this
paper's \(F_k\) for every \(k\ge1\).  If a referee asks for closure without that
citation, the only missing ingredient is therefore a proof of the displayed
ordinary formula for \(C_k\); every subsequent implication is already the
in-paper argument just described.
\end{proof}

\begin{proposition}[Index-locked source interface for intrinsic height]
\label{prop:kmp-indexing-lock-intrinsic-height}
\label{thm:intrinsic-height-kmp-index-lock}
Let \((E_{\mathrm{KMP}})\) denote the following normalized external statement:
with admissible parts \(F_1,F_2,\ldots\) and
\[
  C_k=\max\{R(N):F_k\le N<F_{k+1}\},
\]
one has
\[
  C_{2t+1}=F_{t+1}\quad(t\ge0),
  \qquad C_{2t+2}=2F_t\quad(t\ge1),
  \qquad C_2=1.
\]
The intrinsic height chain
\[
  (E_{\mathrm{KMP}})
  \Longrightarrow \Cref{thm:half-interval-fibonacci-representation-max}
  \Longrightarrow \Cref{thm:intrinsic-fold-fibre-height}
\]
is rigid with respect to the source indexing: a cited ordinary maximum theorem
can be used in this article only after it has been converted to exactly
\((E_{\mathrm{KMP}})\).  Here \((E_{\mathrm{KMP}})\) is precisely Theorem~4.7 of
Koc\'abov\'a--Mas\'akov\'a--Pelantov\'a, after the positive-index identification
\(\mathcal F_k=F_k\) for \(k\ge1\), with the same half-open interval
\(F_k\le N<F_{k+1}\) and the same parity restrictions \(t\ge0\) and \(t\ge1\).
Once this conversion is made, all subsequent assertions
about \(H_n\), \(M_m\), the parity formulas, and the recurrence
\(M_m=M_{m-2}+M_{m-4}\) are internal consequences of
\Cref{thm:half-interval-fibonacci-representation-max,lem:fold-fibre-representation-pairing}.
In particular no theorem about folded residues, maximizing residues, fibre
histograms, or pressure is imported from
Koc\'abov\'a--Mas\'akov\'a--Pelantov\'a.
\end{proposition}

\begin{proof}
The positive-index conversion in
\Cref{thm:source-audited-ordinary-maximum-import} identifies the cited source
sequence \((\mathcal F_k)_{k\ge0}\) with this paper's sequence on all positive
indices: \(\mathcal F_1=F_1\), \(\mathcal F_2=F_2\), and the common Fibonacci
recurrence gives \(\mathcal F_k=F_k\) for every \(k\ge1\).  The admissible parts in
the cited representation problem therefore coincide with this paper's parts
\(F_1,F_2,\ldots\), and the cited half-open intervals
\([\mathcal F_k,\mathcal F_{k+1})\) become exactly \([F_k,F_{k+1})\).  The two
parity clauses of Theorem~4.7 of the cited article are therefore precisely the
two displayed parity clauses for \(C_{2t+1}\) and \(C_{2t+2}\), with the same
restrictions \(t\ge0\) and \(t\ge1\); the value \(C_2=1\) is the direct one-point
calculation on \([F_2,F_3)=[2,3)\).  This proves that every use of the cited
ordinary maximum theorem in the paper is exactly a use of
\((E_{\mathrm{KMP}})\).

Assume now \((E_{\mathrm{KMP}})\).  The proof of
\Cref{thm:half-interval-fibonacci-representation-max} uses only the recursion
\(R_{n+1}(N)=R_n(N)+R_n(N-F_{n+1})\), the largest-sum bound
\(\sum_{j=1}^nF_j=F_{n+2}-2\), and the monotonicity consequence
\(C_n\le C_{n+1}\) read from the displayed parity clauses.  It proves internally
that the shifted half-interval maximum is \(H_n=C_{n+1}\).  The proof of
\Cref{thm:intrinsic-fold-fibre-height} then uses only
\Cref{lem:fold-fibre-representation-pairing} and the bijection
\(\Val_m\colon X_m\to\{0,\ldots,F_{m+2}-1\}\) to identify \(M_m=H_{m+1}\).
Substituting \(n=m+1\) gives \(M_{2s-1}=F_{s+1}\), \(M_{2s}=2F_s\), the listed
initial values, and the separate parity recurrence.  None of these steps cites or
requires any source statement about folded fibres or about the set of maximizing
integers, so the dependency and indexing cut is exact.
\end{proof}

\begin{theorem}[Exact dependency cut for the intrinsic height package]
\label{prop:intrinsic-height-source-dependency-cut}
\label{thm:intrinsic-height-source-dependency-cut}
\label{thm:intrinsic-height-dependency-cut}
Let \((E)\) denote the ordinary interval maximum formula
\[
  C_{2t+1}=F_{t+1}\quad(t\ge0),
  \qquad
  C_{2t+2}=2F_t\quad(t\ge1),
  \qquad C_2=1,
\]
for positive-index Fibonacci representations on \([F_k,F_{k+1})\), with
\(C_k=\max\{R(N):F_k\le N<F_{k+1}\}\).  In this manuscript the whole intrinsic
height and zero-temperature lower/upper envelope chain from
\Cref{thm:half-interval-fibonacci-representation-max} through
\Cref{thm:intrinsic-fold-fibre-height,cor:linfty-fold-entropy,thm:zero-temperature-fold-pressure,thm:diagonal-zero-temperature-fold-endpoint,thm:moment-height-cover-trichotomy}
uses external Fibonacci-representation information only through \((E)\), as
source-audited in \Cref{thm:source-audited-ordinary-maximum-import}.  Conversely,
if \((E)\) is replaced by any self-contained proof or by any other cited theorem
with exactly the same statement and indexing, every displayed intrinsic height,
entropy, and large-moment lower/upper endpoint conclusion in that chain remains
unchanged.
No classification of the maximizing integers, fibre histogram, or
large-deviation statement is used or obtained by this dependency cut.  The
existence of the fixed-degree pressure is supplied separately by
\Cref{thm:primitive-intrinsic-moment-transfer}.
\end{theorem}

\begin{proof}
The first step after the source-audited import is purely internal.
\Cref{cor:citation-sensitive-half-interval-conversion} identifies \((E)\) as the
only cited premise entering \Cref{thm:half-interval-fibonacci-representation-max}:
the proof there reduces the shifted half-interval
\([F_{n+1},2F_{n+1})\) to the ordinary interval \([F_{n+1},F_{n+2})\), controls the
remaining tail by \eqref{eq:R-recursion}, and uses only the monotonicity
\(C_n\le C_{n+1}\) read from the displayed parity formula.  Hence any proof of
\((E)\) with the same positive-index convention gives the same half-interval
values \(H_{2s}=F_{s+1}\) and \(H_{2s+1}=2F_s\).

The next step is also internal.  \Cref{lem:fold-fibre-representation-pairing}
identifies every fibre size of \(\Fold_m\) with
\(R_{m+2}(F_{m+2}+r)\) for a unique residue \(0\le r<F_{m+2}\).  Therefore
\Cref{thm:intrinsic-fold-fibre-height} obtains
\(M_m=H_{m+1}\) by taking a maximum over the residue interval; substituting the
already proved values of \(H_n\) gives the stated formulas for \(M_{2s-1}\) and
\(M_{2s}\).  The entropy corollary uses only those two formulas and the standard
asymptotic \(F_s=\Theta(\varphi^s)\).  The two zero-temperature endpoint theorems
use only the elementary inequalities
\(M_m^q\le S_q^{\mathrm{int}}(m)\le F_{m+2}M_m^q\) and their diagonal analogue,
together with \(|X_m|=F_{m+2}\).  Finally,
\Cref{thm:moment-height-cover-trichotomy} cites the resulting height and endpoint
statements verbatim for its second item.

Thus replacing the audited source theorem by any theorem proving the same
ordinary interval formula changes no subsequent line of argument.  The proof has
never selected or classified a maximizing integer and has never counted the full
distribution of fibre sizes, so the excluded histogram, pressure, and
large-deviation conclusions do not follow from the cut.
\end{proof}

\begin{theorem}[Dependency firewall for the intrinsic height theorem]
\label{thm:ordinary-maximum-supports-intrinsic-height}
\label{thm:intrinsic-height-dependency-firewall}
Let \((E)\) be the ordinary positive-index interval maximum formula of
\Cref{thm:intrinsic-height-source-dependency-cut}.  The part of the \(L2\)
theorem chain concerning the intrinsic \(q=\infty\) height factors through the
following three internal implications and through no other external
Fibonacci-representation input:
\[
  \begin{aligned}
  (E)&\Longrightarrow H_n=C_{n+1}
      \Longrightarrow M_m=H_{m+1},\\
      &\Longrightarrow
      M_{2s-1}=F_{s+1},\quad M_{2s}=2F_s,\\
      &\Longrightarrow
      \lim_{m\to\infty}m^{-1}\log M_m=\frac12\log\varphi
      \quad\text{and the zero-temperature endpoint bounds.}
  \end{aligned}
\]
More explicitly, after \((E)\) is granted, every displayed assertion in
\Cref{thm:half-interval-fibonacci-representation-max,thm:intrinsic-fold-fibre-height,cor:linfty-fold-entropy,thm:zero-temperature-fold-pressure,thm:diagonal-zero-temperature-fold-endpoint} follows from in-paper arguments,
while none of those arguments imports a classification of maximizing integers, a
full fibre histogram, or a large-deviation principle.  The fixed-degree pressure
limit uses the primitive carry transfer in addition to this height input.  Thus
the ordinary maximum formula is supporting background for the
height package, whereas the intrinsic fold-height statement itself is the
internal fold-aware consequence of the representation pairing.  Consequently the
article proves a conditional conversion theorem with cited premise \((E)\), not a
new ordinary-interval maximum theorem: deleting \((E)\) deletes the numerical
height formulas, while leaving intact the fold-aware identity
\(M_m=\max_{0\le r<F_{m+2}}R_{m+2}(F_{m+2}+r)\).
This firewall is also a replacement principle.  Any internal proof or different
citation may replace Koc\'abov\'a--Mas\'akov\'a--Pelantov\'a at this point only if
it proves exactly \((E)\), with the same positive-index parts and half-open
intervals.  A source theorem with shifted indices, different endpoint
conventions, or extra maximizing-integer classifications supplies no numerical
intrinsic height conclusion until it has first been converted to \((E)\).
\end{theorem}

\begin{proof}
By \Cref{thm:source-audited-ordinary-maximum-import}, the cited source is used
only to obtain the positive-index formula \((E)\) for
\(C_k=\max\{R(N):F_k\le N<F_{k+1}\}\), after the normalization of conventions
recorded there.  No maximizing argument or distributional statement from the
source enters the proof.

Assume \((E)\).  The proof of
\Cref{thm:half-interval-fibonacci-representation-max} is an internal conversion:
on the ordinary subinterval \([F_{n+1},F_{n+2})\), the truncated count
\(R_{n+1}\) agrees with the ordinary representation count, so the maximum there
is \(C_{n+1}\).  On the remaining tail, write
\(N=F_{n+2}+u\) with \(0\le u<F_{n-1}\).  The recursion \eqref{eq:R-recursion}
gives
\[
  R_{n+1}(N)=R_n(F_{n+2}+u)+R_n(F_n+u)=R_n(F_n+u),
\]
because the largest sum formed from \(F_1,\ldots,F_n\) is \(F_{n+2}-2\).  Since
\(F_n\le F_n+u<F_{n+1}\), this tail contribution is at most \(C_n\), and the
parity formula \((E)\), including the direct values \(C_1=C_2=1\), gives
\(C_n\le C_{n+1}\).  Hence the shifted half-interval maximum is exactly
\(H_n=C_{n+1}\), and substituting the two parity clauses of \((E)\) gives the
displayed formulas for \(H_n\).

The passage from shifted half-interval maxima to intrinsic fold height is also
internal.  By \Cref{lem:fold-fibre-representation-pairing}, for
\(r=\Val_m(x)\),
\[
  |\Fold_m^{-1}(x)|=R_{m+2}(F_{m+2}+r),\qquad 0\le r<F_{m+2}.
\]
The map \(\Val_m\) is a bijection from \(X_m\) to the residue interval, so taking
maxima gives \(M_m=H_{m+1}\).  The formulas
\(M_{2s-1}=F_{s+1}\) and \(M_{2s}=2F_s\) are then the half-interval formulas with
\(n=m+1\).

If the cited premise \((E)\) is withheld, the argument up to this point still
proves the purely intrinsic reduction
\[
  M_m=\max_{0\le r<F_{m+2}}R_{m+2}(F_{m+2}+r),
\]
because that equality is only \Cref{lem:fold-fibre-representation-pairing} plus
the bijectivity of \(\Val_m\).  What disappears is exactly the evaluation of this
maximum by the ordinary interval values \(C_k\).  Hence the proof cannot be read
as an independent derivation of \((E)\), and it cannot export the ordinary maximum
formula back to the representation problem.

The entropy and zero-temperature endpoint statements use no further external
representation theory.  The entropy is the asymptotic consequence of
\(F_s=\Theta(\varphi^s)\).  For each fixed \(q\), the zero-temperature squeeze uses
only
\[
  M_m^q\le S_q^{\mathrm{int}}(m)
  \le |X_m|M_m^q=F_{m+2}M_m^q,
\]
and the diagonal variant uses the same elementary inequality with the chosen
moment degree.  These steps neither select maximizing integers nor count the
whole fibre-size multiset.  Therefore the dependency factorization displayed in
the statement is exact for the height package used later in the \(L2\) interface.

It remains only to justify the replacement clause.  Inspecting the preceding
paragraphs, the proof has used the external source through exactly two facts:
the values of \(C_k\) in \((E)\) and the resulting monotonicity
\(C_n\le C_{n+1}\).  The latter is an immediate consequence of the displayed
values in \((E)\), including \(C_1=C_2=1\), and is not an additional imported
theorem.  Hence any proof of precisely \((E)\), whether internal or cited, can be
substituted verbatim.  Conversely, a theorem stated with a different Fibonacci
normalization or interval convention does not identify the quantities
\(C_k=\max\{R(N):F_k\le N<F_{k+1}\}\) used in the tail estimate above; before the
positive-index and endpoint conversion is supplied, the proof still has only the
unevaluated fold-aware maximum identity and cannot derive any of the displayed
numerical formulas for \(H_n\) or \(M_m\).  This proves the claimed firewall and
prevents the ordinary interval maximum formula from being presented as a new
result of the manuscript.
\end{proof}

\begin{theorem}[Audited shifted-half-interval height conversion]
\label{thm:audited-shifted-half-interval-height-conversion}
Grant only the ordinary positive-index maximum formula \((E)\) of
\Cref{thm:intrinsic-height-source-dependency-cut}.  Then the conversion to the
intrinsic height theorem is completely determined by the following internal
chain:
\[
  C_{n+1}\ge C_n,
  \qquad
  H_n=C_{n+1},
  \qquad
  M_m=H_{m+1}.
\]
In particular the parity formulas
\[
  H_{2s}=F_{s+1},\quad H_{2s+1}=2F_s,
  \qquad
  M_{2s-1}=F_{s+1},\quad M_{2s}=2F_s
\]
follow without importing any theorem about folded fibres, maximizing residues,
fibre histograms, or pressure.  The only Fibonacci-index conversion used is
\(\mathcal F_k=F_k\) for \(k\ge1\), as proved in
\Cref{lem:positive-index-conversion-ordinary-maximum}.
\end{theorem}

\begin{proof}
The formula \((E)\) gives the ordinary maxima
\(C_k=\max\{R(N):F_k\le N<F_{k+1}\}\) in this paper's positive-index convention.
The displayed parity values immediately give \(C_n\le C_{n+1}\) for every
\(n\ge2\), and the initial value \(H_1=1\) is checked directly from
\(R_2(2)=R_2(3)=1\).  Thus it remains to treat \(n\ge2\).

On the subinterval \([F_{n+1},F_{n+2})\), no part larger than \(F_{n+1}\) can occur
in a representation of \(N\).  Hence the truncated count \(R_{n+1}(N)\) is the
ordinary count there, and the maximum on that subinterval is \(C_{n+1}\).  On the
remaining part of the shifted half-interval, write
\(N=F_{n+2}+u\) with \(0\le u<F_{n-1}\).  The recursion
\eqref{eq:R-recursion} gives
\[
  R_{n+1}(N)=R_n(F_{n+2}+u)+R_n(F_n+u)=R_n(F_n+u),
\]
because the largest sum using only \(F_1,\ldots,F_n\) is \(F_{n+2}-2\).  Since
\(F_n\le F_n+u<F_{n+1}\), the tail maximum is at most \(C_n\le C_{n+1}\).  Therefore
\(H_n=C_{n+1}\).  Substituting the two parity clauses of \((E)\), with index
\(k=n+1\), gives \(H_{2s}=F_{s+1}\) and \(H_{2s+1}=2F_s\).

Finally, \Cref{lem:fold-fibre-representation-pairing} identifies each intrinsic
fibre size with \(R_{m+2}(F_{m+2}+r)\) for a unique
\(0\le r<F_{m+2}\), because \(\Val_m\colon X_m\to
\mathbb Z/F_{m+2}\mathbb Z\) is bijective.  Taking maxima over residues gives
\(M_m=H_{m+1}\), and the stated formulas for \(M_{2s-1}\) and \(M_{2s}\) are the
substitution \(n=m+1\).  Every step after \((E)\) is the displayed recursion or
the fold-fibre pairing, so no additional external enumerative or thermodynamic
assertion is used.
\end{proof}

\begin{theorem}[Half-interval Fibonacci representation maximum]
\label{thm:half-interval-fibonacci-representation-max}
For \(n\ge1\), set
\[
  H_n=\max_{0\le r<F_{n+1}} R_{n+1}(F_{n+1}+r).
\]
Then \(H_1=1\) and, for \(s\ge1\),
\[
  H_{2s}=F_{s+1},\qquad H_{2s+1}=2F_s .
\]
Equivalently,
\[
  H_n=H_{n-2}+H_{n-4}\qquad(n\ge6),
\]
with initial values \(H_1=1,H_2=2,H_3=2,H_4=3,H_5=4\).
The ordinary interval maximum formula itself is not a new theorem of this
article: it is exactly the cited positive-index formula \((E)\), equivalently
\((E_{\mathrm{KMP}})\), isolated as external background in
\Cref{thm:source-audited-ordinary-maximum-import,prop:kmp-indexing-lock-intrinsic-height}, namely Theorem~4.7 of
Koc\'abov\'a--Mas\'akov\'a--Pelantov\'a
\cite[Thm.~4.7]{KocabovaMasakovaPelantova2005}.  The bibliography and Numdam
metadata used for this import are the source-audited metadata recorded in
\Cref{thm:source-audited-ordinary-maximum-import}; no later step treats that
external theorem as internally proved.  The imported interval convention is
exactly
\[
  F_k\le N<F_{k+1}
\]
for ordinary distinct positive-index Fibonacci representations.  The theorem
proved here converts that imported ordinary interval maximum to the shifted
truncated half-interval
\[
  F_{n+1}\le N<2F_{n+1},
  \qquad N=F_{n+1}+r,\qquad 0\le r<F_{n+1},
\]
with parts restricted to \(F_1,\ldots,F_{n+1}\).  No maximizing-integer
classification or folded-fibre theorem is imported from the citation.
Thus the theorem below is an intrinsic shifted-half-interval conversion theorem,
not a second citation to an ordinary maximum theorem.  Its proof calls the
external theorem only through the index-locked quantities \(C_k\) of
\((E_{\mathrm{KMP}})\) and then works entirely with \(R_{n+1}\),
\eqref{eq:R-recursion}, and the fold-fibre pairing.
\end{theorem}

\begin{proof}
The case \(H_1=1\) is immediate from \(R_2(2)=R_2(3)=1\).  For
\(n\ge2\), the only external input in the height theorem is the ordinary interval
maximum formula imported and normalized in
\Cref{lem:fibonacci-maximum-normalization-conversion}; the rest of the proof is
the local conversion of that ordinary interval maximum to the half-interval
notation used here.  Thus, in the notation
\[
  C_k=\max\{R(N):F_k\le N<F_{k+1}\},
\]
we have \(C_1=C_2=1\), and the imported parity clauses give, for \(t\ge1\),
\[
  C_{2t+1}=F_{t+1},\qquad C_{2t+2}=2F_t .
\]
On \([F_{n+1},F_{n+2})\), the truncated count \(R_{n+1}\) is the ordinary count,
because no part above \(F_{n+1}\) can occur.  Therefore the maximum of
\(R_{n+1}\) on this subinterval is \(C_{n+1}\).

It remains only to check the extra half-interval tail.  Write
\(N=F_{n+2}+u\), where
\(0\le u<2F_{n+1}-F_{n+2}=F_{n-1}\).  By \eqref{eq:R-recursion},
\[
  R_{n+1}(N)=R_n(F_{n+2}+u)+R_n(F_n+u)=R_n(F_n+u),
\]
because the largest sum using \(F_1,\ldots,F_n\) is \(F_{n+2}-2\).  The latter
quantity is bounded by the ordinary interval maximum \(C_n\), and
\(C_n\le C_{n+1}\) from the displayed formula.  Hence the enlarged interval
\([F_{n+1},2F_{n+1})\) has the same maximum as \([F_{n+1},F_{n+2})\), namely
\(H_n=C_{n+1}\).  Substituting \(k=n+1\) in the formula for \(C_k\) gives the two
parity formulas for \(H_n\).  The recurrence and initial values follow directly
from those formulas and \(F_{j+1}=F_j+F_{j-1}\).  Thus the cited Theorem~4.7 is
used exactly once, to supply the ordinary maxima \(C_k\); every passage from
\(C_k\) to \(H_n\), and later from \(H_n\) to intrinsic \(\Fold_m\)-fibre height,
is the internal conversion proved here.
\end{proof}

\begin{theorem}[Exact intrinsic \texorpdfstring{\(\Fold_m\)}{Fold_m} fibre height]
\label{thm:intrinsic-fold-fibre-height}
Let
\[
  M_m=\max_{x\in X_m}|\Fold_m^{-1}(x)|.
\]
Then \(M_m=H_{m+1}\).  Consequently, for \(s\ge1\),
\[
  M_{2s-1}=F_{s+1},\qquad M_{2s}=2F_s .
\]
The first ten values are
\[
  M_1,\ldots,M_{10}=2,2,3,4,5,6,8,10,13,16.
\]
On each parity class the heights satisfy
\[
  M_m=M_{m-2}+M_{m-4}\qquad(m\ge5),
\]
with initial values \(M_1=2,M_2=2,M_3=3,M_4=4\).
The only imported enumerative input is the cited-background, index-locked
ordinary interval maximum \([F_k,F_{k+1})\) of
\Cref{thm:source-audited-ordinary-maximum-import,prop:kmp-indexing-lock-intrinsic-height}; the shifted
half-interval conversion to \(H_{m+1}\) and the fibre identification
\(M_m=H_{m+1}\) are the in-paper steps of
\Cref{thm:half-interval-fibonacci-representation-max,lem:fold-fibre-representation-pairing}.
Equivalently, the theorem is an intrinsic fold-aware conversion from the cited
ordinary maximum formula, and not a new proof of that ordinary formula.
\end{theorem}

\begin{proof}
By \Cref{lem:fold-fibre-representation-pairing}, for \(r=\Val_m(x)\),
\[
  |\Fold_m^{-1}(x)|=R_{m+2}(F_{m+2}+r),
  \qquad 0\le r<F_{m+2}.
\]
Since \(\Val_m\) is a bijection from \(X_m\) to these residues, taking the maximum
over \(x\in X_m\) gives \(M_m=H_{m+1}\).  This step uses no ordinary representation
maximum theorem; it is only the residue-fibre pairing over the shifted interval
\([F_{m+2},2F_{m+2})\).  Substituting \(n=m+1\) in
\Cref{thm:half-interval-fibonacci-representation-max} gives
\(M_{2s-1}=H_{2s}=F_{s+1}\) and \(M_{2s}=H_{2s+1}=2F_s\).  The ten displayed
values and the parity recurrence are the same substitution applied to the listed
initial values and recurrence for \(H_n\).
\end{proof}

\begin{proposition}[Ordinary-maximum import isolation for the height law]
\label{prop:ordinary-maximum-import-isolation}
\label{thm:ordinary-maximum-citation-boundary-for-height}
In the proof of \Cref{thm:intrinsic-fold-fibre-height}, the only external
enumerative assertion is the ordinary half-open interval maximum formula
\((E_{\mathrm{KMP}})\) isolated in
\Cref{thm:source-audited-ordinary-maximum-import,prop:kmp-indexing-lock-intrinsic-height}.
All other steps are internal fold-aware conversions: the representation pairing
of \Cref{lem:fold-fibre-representation-pairing}, the shifted half-interval
conversion of \Cref{thm:half-interval-fibonacci-representation-max}, and the
maximization over the residue-normal-form bijection \(\Val_m\colon X_m\to
\{0,
\ldots,F_{m+2}-1\}\).  Consequently any citation to
\Cref{thm:intrinsic-fold-fibre-height} in this article is a citation to the
intrinsic \(\Fold_m\)-height conversion, not to a newly proved ordinary maximum
theorem for all Fibonacci representations.
\end{proposition}

\begin{proof}
The source-audited import lemma identifies exactly what is taken from
Koc\'abov\'a--Mas\'akov\'a--Pelantov\'a: after the positive-index normalization,
Theorem~4.7 of the cited paper supplies the values of
\(C_k=\max\{R(N):F_k\le N<F_{k+1}\}\).  The proof of
\Cref{thm:half-interval-fibonacci-representation-max} then uses only these
numbers \(C_k\), the recursion \eqref{eq:R-recursion}, and elementary interval
containments to prove the shifted maximum \(H_n\) for the truncated functions
\(R_{n+1}\) on \([F_{n+1},2F_{n+1})\).  That argument does not classify ordinary
maximizers and does not extend the ordinary theorem; it converts the imported
ordinary maximum into the shifted truncated interval needed by the fold map.

The final height theorem applies \Cref{lem:fold-fibre-representation-pairing}:
for \(r=\Val_m(x)\), the fibre size is
\(R_{m+2}(F_{m+2}+r)\).  Since \(\Val_m\) is a bijection onto the displayed
residue interval, maximizing over \(x\in X_m\) is exactly maximizing the shifted
truncated representation count already computed as \(H_{m+1}\).  Thus the only
external input is the ordinary formula for the \(C_k\), and every subsequent
operation is internal to the fold-aware normal-form system.  Any later use of the
height law therefore imports no ordinary maximum theorem beyond the cited
\((E_{\mathrm{KMP}})\) boundary.
\end{proof}


% End original section: sec04_intrinsic_height.tex
% Original section: sec04_intrinsic_pressure.tex
\begin{corollary}[\texorpdfstring{\(L^\infty\)}{L-infinity} fold entropy]
\label{cor:linfty-fold-entropy}
Let \(\varphi=(1+\sqrt5)/2\).  Then
\[
  \lim_{m\to\infty}\frac1m\log M_m=\frac12\log\varphi,
  \qquad
  M_m=\Theta(\varphi^{m/2}).
\]
This is an intrinsic \(q=\infty\) endpoint for the unmodified fold, not a
consequence of any single fixed-\(q\) recurrence in
\Cref{thm:intrinsic-fold-moment-transfer}.
\end{corollary}

\begin{proof}
The exact height law gives \(M_{2s-1}=F_{s+1}\) and \(M_{2s}=2F_s\).  Since
\(F_s=\Theta(\varphi^s)\), both parity subsequences are
\(\Theta(\varphi^{m/2})\), and division of logarithms by \(m\) gives the stated
limit.  The fixed-\(q\) transfer theorem concerns each finite moment degree
separately; the maximum fibre size is the \(q=\infty\) norm and is supplied here
by the separate representation-pairing and half-interval maximum argument.
\end{proof}

\begin{theorem}[Fixed-degree intrinsic pressure and its zero-temperature limit]
\label{thm:zero-temperature-fold-pressure}
\label{thm:height-squeeze-endpoint}
For every fixed \(q\ge1\), the limit
\[
  p(q)=\lim_{m\to\infty}\frac1m\log S_q^{\mathrm{int}}(m)
\]
exists.  Moreover,
\[
  M_m^q\le S_q^{\mathrm{int}}(m)\le F_{m+2}M_m^q
\]
for every \(m\), and hence
\[
  q\frac12\log\varphi\le p(q)
  \le q\frac12\log\varphi+\log\varphi .
\]
Consequently
\[
  \lim_{q\to\infty}\frac{p(q)}q=\frac12\log\varphi .
\]
The existence of \(p(q)\) is intrinsic finite-state information; the displayed
bounds and the normalized \(q\to\infty\) limit use only the exact fibre height.
No fibre-histogram law or large-deviation principle is asserted.
\end{theorem}

\begin{proof}
Existence and the identity \(p(q)=\log\lambda_q\) follow from
\Cref{thm:primitive-intrinsic-moment-transfer}.  The lower bound is the
contribution of a point of maximal fibre size.  The upper
bound uses \(d_m(x)\le M_m\) for all \(x\in X_m\) and
\(|X_m|=F_{m+2}\).  Taking \(m^{-1}\log\), passing to the limit, and using
\Cref{cor:linfty-fold-entropy} together with
\(\lim_{m\to\infty}m^{-1}\log F_{m+2}=\log\varphi\), gives the displayed
pressure bounds.  Dividing by \(q\) and letting \(q\to\infty\) gives the final
limit by the squeeze theorem.
\end{proof}

\begin{theorem}[Diagonal height endpoint for intrinsic fold moments]
\label{thm:diagonal-zero-temperature-fold-endpoint}
\label{thm:diagonal-height-endpoint}
Let \((q_m)_{m\ge1}\) be any sequence of positive integers with
\(q_m\to\infty\).  Then the original Zeckendorf fold fibres satisfy
\[
  \lim_{m\to\infty}\frac{1}{m q_m}
  \log\sum_{x\in X_m}|\Fold_m^{-1}(x)|^{q_m}
  ={\frac12}\log\varphi .
\]
Equivalently, the endpoint detected in \Cref{thm:zero-temperature-fold-pressure}
is only a height-squeeze consequence of the exact maximal-fibre law.  It controls
the normalized diagonal large-moment slope for every diverging moment degree, but
its height-squeeze proof does not assert a fibre histogram, a fixed-degree
pressure limit, or a large-deviation principle.
This is an intrinsic theorem about the unmodified maps \(\Fold_m\), uses no
auxiliary cover fibres, and supplies only the \(q=\infty\) endpoint needed in
\Cref{thm:main-hierarchy}.
\end{theorem}

\begin{proof}
Write
\(S_q^{\mathrm{int}}(m)=\sum_{x\in X_m}|\Fold_m^{-1}(x)|^q\) and
\(M_m=\max_{x\in X_m}|\Fold_m^{-1}(x)|\).  For every positive integer \(q\),
the contribution of a fibre of maximal size and the trivial upper bound on all
fibres give
\[
  M_m^q\le S_q^{\mathrm{int}}(m)\le |X_m|M_m^q=F_{m+2}M_m^q .
\]
Apply this with \(q=q_m\), take logarithms, and divide by \(m q_m\).  We obtain
\[
  \frac1m\log M_m
  \le
  \frac{1}{m q_m}\log S_{q_m}^{\mathrm{int}}(m)
  \le
  \frac1m\log M_m+\frac{1}{m q_m}\log F_{m+2} .
\]
By \Cref{cor:linfty-fold-entropy}, the first term tends to
\((1/2)\log\varphi\).  Also
\(m^{-1}\log F_{m+2}\to\log\varphi\), while \(q_m\to\infty\); hence the final
error term tends to zero.  The squeeze theorem gives the asserted limit.
\end{proof}

The fixed-degree matrices are uniform in the resolution parameter, exactly as
required by \Cref{def:fixed-finite-state-transfer}.  Uniformity in the moment
degree is a different question.  Encoding the whole moment tower by a moment
marker gives an intrinsic obstruction to a single rational-function all-degree
law in that marker.

\begin{theorem}[Intrinsic rational all-degree transfer obstruction]
\label{thm:intrinsic-all-degree-transfer-obstruction}
For \(m\ge1\), set
\[
  G_m(w)=\sum_{q\ge0}S_q^{\mathrm{int}}(m)w^q,
  \qquad
  \mathcal M(z,w)=\sum_{m\ge1}G_m(w)z^m.
\]
Every fixed-\(q\) slice
\(\sum_{m\ge1}S_q^{\mathrm{int}}(m)z^m\) is rational and has the primitive
transfer asymptotic of
\Cref{thm:intrinsic-fold-moment-transfer,thm:primitive-intrinsic-moment-transfer},
but the intrinsic diagonal sequence
\[
  \Delta_m=S_m^{\mathrm{int}}(m)
\]
has no fixed finite-state transfer formula in the sense of
\Cref{def:fixed-finite-state-transfer}.  More strongly, the bivariate series
\(\mathcal M(z,w)\) is not a rational function.

Consequently there do not exist an integer \(r\ge1\), rational-function data
\[
  A(w)\in M_r(\mathbb C(w)),
  \qquad u(w),v(w)\in\mathbb C(w)^r,
\]
and \(m_0\ge1\) such that, as formal power series at \(w=0\),
\[
  G_m(w)=u(w)^{\mathsf T}A(w)^m v(w)
  \qquad(m\ge m_0).
\]
Thus the unmodified Zeckendorf fold has intrinsic fixed finite-state transfer in
every fixed moment degree, but its full intrinsic moment tower admits no single
rational finite-state transfer law.
\end{theorem}

\begin{proof}
A fibre of maximal size contributes
\[
  \Delta_m=S_m^{\mathrm{int}}(m)\ge M_m^m.
\]
By \Cref{cor:linfty-fold-entropy},
\[
  \log M_m=\frac12m\log\varphi+O(1),
\]
and hence
\[
  \log S_m^{\mathrm{int}}(m)
  \ge\frac12m^2\log\varphi-C_0m
\]
for some \(C_0>0\) and all sufficiently large \(m\).  Thus \(\Delta_m\) exceeds
\(C\rho^m\) eventually for every \(C,\rho>0\).  The
exponential budget of \Cref{prop:finite-state-coefficient-budget} proves that
\((\Delta_m)\) has no fixed finite-state transfer formula.

If \(\mathcal M\) were rational, its formal expansion at the origin could be
written as \(P/Q\) with \(P,Q\in\mathbb C[z,w]\) and \(Q(0,0)\ne0\).  It would
therefore be analytic on some closed polydisc about the origin.  Cauchy's
coefficient estimate would give constants \(C,R_z,R_w>0\) such that
\[
  S_q^{\mathrm{int}}(m)\le C R_z^{-m}R_w^{-q}
  \qquad(m\ge1,\ q\ge0).
\]
On the diagonal \(q=m\), this gives an \(O(m)\) upper bound for
\(\log S_m^{\mathrm{int}}(m)\).
This contradicts the preceding quadratic lower bound.  Therefore
\(\mathcal M\) is not rational.

For the transfer consequence, each fixed-resolution series is rational:
\[
  G_m(w)
  =\sum_{x\in X_m}\sum_{q\ge0}\bigl(d_m(x)w\bigr)^q
  =\sum_{x\in X_m}\frac1{1-d_m(x)w}.
\]
If the displayed all-degree transfer law held on a tail, then
\[
  \sum_{m\ge m_0}G_m(w)z^m
  =z^{m_0}u(w)^{\mathsf T}A(w)^{m_0}
    (I-zA(w))^{-1}v(w)
\]
would be rational in \(z,w\).  Adding the finitely many rational initial
terms would make \(\mathcal M(z,w)\) rational, contradicting the first part.
\end{proof}

\begin{theorem}[Height-squeeze firewall and pressure non-determination]
\label{prop:height-squeeze-only-endpoint-firewall}
\label{thm:height-squeeze-pressure-nondetermination}
The pressure bounds and normalized \(q\to\infty\) conclusion in
\Cref{thm:zero-temperature-fold-pressure}, together with the diagonal conclusion
of \Cref{thm:diagonal-zero-temperature-fold-endpoint}, depend only on the
cardinality \(|X_m|=F_{m+2}\) and the maximal intrinsic fibre height
\(M_m=\max_{x\in X_m}|\Fold_m^{-1}(x)|\).  The existence of the fixed-degree
pressure itself instead comes from the primitive intrinsic transfer theorem.
More precisely, for any choice
of integers \(a_m(x)\) with \(1\le a_m(x)\le M_m\) and
\(\max_{x\in X_m}a_m(x)=M_m\), the sums
\[
  T_q(m)=\sum_{x\in X_m}a_m(x)^q
\]
satisfy the same fixed-\(q\) endpoint envelope and the same diagonal
large-moment limit as the intrinsic sums.  These height-squeeze conclusions
therefore do not determine, and are not used here as, a full pressure limit, a
fibre-histogram law, or a large-deviation principle.

The obstruction is sharp already inside the class with the same cardinality and
maximum.  There are two arrays \(a_m,b_m\colon X_m\to\{1,
M_m\}\) with \(\max a_m=\max b_m=M_m\) such that both arrays satisfy the same
endpoint conclusions above, but for every fixed \(q\ge1\)
\begin{align*}
  \lim_{m\to\infty}{\frac1m}\log\sum_{x\in X_m} a_m(x)^q
  &=\max\!\left\{\log\varphi,{\frac q2}\log\varphi\right\},\\
  \lim_{m\to\infty}{\frac1m}\log\sum_{x\in X_m} b_m(x)^q
  &=\left(1+{\frac q2}\right)\log\varphi .
\end{align*}
Thus the height-squeeze data alone cannot be upgraded to a fixed-\(q\) pressure
theorem or to a histogram classification.
\end{theorem}

\begin{proof}
The hypotheses give the same elementary bounds used above:
\[
  M_m^q\le T_q(m)\le |X_m|M_m^q=F_{m+2}M_m^q .
\]
Together with \Cref{cor:linfty-fold-entropy} and
\(m^{-1}\log F_{m+2}\to\log\varphi\), the proof of
\Cref{thm:zero-temperature-fold-pressure} applies verbatim to the lower and
upper exponential envelopes of \(T_q(m)\).  If \(q_m\to\infty\), the same bounds
give
\[
  \frac1m\log M_m
  \le \frac{1}{m q_m}\log T_{q_m}(m)
  \le \frac1m\log M_m+\frac{1}{m q_m}\log F_{m+2},
\]
so the proof of \Cref{thm:diagonal-zero-temperature-fold-endpoint} also applies
verbatim.

It remains to record why this is a genuine firewall rather than hidden
thermodynamics.  Fix one point \(x_m^\ast\in X_m\).  The two arrays
\[
  a_m(x_m^\ast)=M_m,\quad a_m(x)=1\ (x\ne x_m^\ast),
  \qquad
  b_m(x)=M_m\quad(x\in X_m)
\]
have the same cardinality and the same maximum height, and hence the same
height-squeeze endpoints by the first paragraph.  Their empirical distributions
of the normalized log-sizes are nevertheless different: for \(a_m\) the mass at
\(m^{-1}\log M_m\) has weight \(F_{m+2}^{-1}\to0\), while for \(b_m\) all mass is
at \(m^{-1}\log M_m\to(1/2)\log\varphi\).  Their fixed-degree exponential sums
also differ, since
\[
  \sum_x a_m(x)^q=M_m^q+F_{m+2}-1,
  \qquad
  \sum_x b_m(x)^q=F_{m+2}M_m^q .
\]
Using \Cref{cor:linfty-fold-entropy},
\(m^{-1}\log M_m\to(1/2)\log\varphi\), and
\(m^{-1}\log F_{m+2}\to\log\varphi\), these identities give the two displayed
fixed-\(q\) limits.  For \(b_m\) the formula is immediate.  For \(a_m\), the
exponential growth of a sum of two positive sequences is the maximum of their
exponential growth rates.  The empirical distributions are also incompatible:
for \(a_m\) the mass at the maximal normalized log-size has weight
\(F_{m+2}^{-1}\to0\), whereas for \(b_m\) all mass is at that point.  Thus the
height-squeeze endpoint is compatible with incompatible histograms and different
fixed-degree growth envelopes.  Consequently the height-squeeze argument carries
only the maximal-height information stated above; the fixed-degree pressure
requires the primitive transfer matrix, and neither argument supplies a histogram
or large-deviation assertion.
\end{proof}


\begin{theorem}[Intrinsic--cover transfer budget dichotomy]
\label{thm:moment-height-cover-trichotomy}
\label{thm:intrinsic-cover-transfer-budget-dichotomy}
For the Zeckendorf fold, the \(L2\) finite-fibre interface splits into a sharp
intrinsic budget and a separate cover-relative obstruction.
\begin{enumerate}
\item For every fixed finite \(q\), the intrinsic moment sequence
  \((S_q^{\mathrm{int}}(m))_{m\ge1}\) has fixed finite-state transfer.  Hence it
  satisfies a constant-coefficient recurrence, has a rational generating series,
  and has a primitive Perron asymptotic and a pressure
  \(p(q)=\log\lambda_q\).
\item The full intrinsic moment tower has no single rational finite-state
  transfer law: its bivariate generating series
  \(\mathcal M(z,w)=\sum_{m\ge1}\sum_{q\ge0}
  S_q^{\mathrm{int}}(m)z^mw^q\) is not rational.
\item At \(q=\infty\), the intrinsic fibre height is exactly
  \[
    M_{2s-1}=F_{s+1},\qquad M_{2s}=2F_s,
  \]
  so
  \[
    \lim_{m\to\infty}{\frac1m}\log M_m={\frac12}\log\varphi .
  \]
  Moreover the intrinsic large-moment endpoint is
  \((1/2)\log\varphi\) both through the fixed-degree pressures of
  \Cref{thm:zero-temperature-fold-pressure} and in the diagonal sense of
  \Cref{thm:diagonal-zero-temperature-fold-endpoint}.
\item In contrast, in the computable finite-fibre cover category over the same
  normal-form sets \(X_m\), finite-resolution certificates impose no such
  transfer budget: there is a computable one-spike cover preserving all normal
  forms and all \(L1\) readouts whose second cover moment grows faster than every
  exponential sequence and therefore has no fixed finite-state transfer formula.
\item Quantitatively, the preceding cover obstruction cannot be hidden inside the
  intrinsic height scale.  If
  \(H_m^\sharp=\max_{x\in X_m}d_m^\sharp(x)\) is the maximal cover fibre size and
  \(H_m^\sharp\le B\sigma^m M_m\) for some constants \(B>0\), \(\sigma>1\), then
  \((S_2^\sharp(m))\) is exponentially bounded.  Hence any cover whose second
  moment violates every exponential bound must have
  \(H_m^\sharp/M_m\) not exponentially bounded.  The factorial one-spike cover
  realizes this lower bound: its spike height is \(m!+1\), while
  \(M_m=\Theta(\varphi^{m/2})\).
\end{enumerate}
Consequently the \(L2\) theorem chain used by \Cref{thm:main-hierarchy} has an
intrinsic separation between degreewise transfer and a single rational-function
all-degree law in the moment marker.  Its stronger statement that finite-fibre
certificates need not force transfer even in degree two remains cover-relative and requires
certificate spikes beyond the intrinsic height by more than every exponential
factor.  Neither assertion requires a full intrinsic fibre-histogram or
large-deviation classification for the unmodified \(\Fold_m\) fibres.
\end{theorem}

\begin{proof}
For the first item, \Cref{thm:intrinsic-fold-moment-transfer} gives a fixed
finite-state transfer formula for \((S_q^{\mathrm{int}}(m))_{m\ge1}\) in every
fixed degree \(q\).  Applying \Cref{prop:fixed-transfer-interface-budget} to this
sequence gives the constant-coefficient recurrence, rational generating series,
and exponential growth bound, while
\Cref{thm:primitive-intrinsic-moment-transfer} gives the Perron asymptotic and
pressure formula.

The second item is
\Cref{thm:intrinsic-all-degree-transfer-obstruction}: the diagonal coefficients
grow like \(\exp((\frac12\log\varphi+o(1))m^2)\), which is incompatible with a
rational bivariate generating series.

For the third item, \Cref{thm:intrinsic-fold-fibre-height} gives the displayed
closed formula for \(M_m\).  Since
\(F_n=(\varphi^n-(-\varphi)^{-n})/\sqrt5\), both subsequences have logarithmic
slope \((1/2)\log\varphi\): for odd indices,
\(m=2s-1\) and \(m^{-1}\log F_{s+1}\to(1/2)\log\varphi\); for even indices,
\(m=2s\) and \(m^{-1}\log(2F_s)\to(1/2)\log\varphi\).  The large-moment endpoint
assertions are precisely
\Cref{thm:zero-temperature-fold-pressure,thm:diagonal-zero-temperature-fold-endpoint}.

For the fourth item, apply \Cref{thm:L2-finite-fiber-no-transfer-content} with
\(a_m=m!\).  The resulting cover is computable and preserves normal forms and
readouts by construction, but its second moment is
\((m!+1)^2+F_{m+2}-1\), which grows faster than \(C\rho^m\) for every pair of
positive constants \(C,\rho\).  By \Cref{prop:fixed-transfer-interface-budget}, every
fixed finite-state transfer sequence is exponentially bounded; hence this cover
moment has no fixed finite-state transfer formula.

The first three clauses concern the original maps \(\Fold_m\), while the fourth
concerns auxiliary cover
multiplicities over the already fixed normal forms.  Therefore the negative
statement cannot be read as an intrinsic nontransfer theorem without an extra
bridge identifying intrinsic fold multiplicities with the chosen cover
multiplicities; that bridge is excluded by
\Cref{thm:intrinsic-cover-transfer-dichotomy,thm:L2-cover-intrinsic-nonimplication}.

It remains to prove the quantitative height-relative assertion.  For any cover,
\[
  S_2^\sharp(m)=\sum_{x\in X_m}d_m^\sharp(x)^2
  \le |X_m|\,(H_m^\sharp)^2 .
\]
If \(H_m^\sharp\le B\sigma^mM_m\), then \(|X_m|=F_{m+2}\) and
\Cref{thm:intrinsic-fold-fibre-height,cor:linfty-fold-entropy} give constants
\(C>0\) and \(\lambda>1\) such that \(|X_m|M_m^2\le C\lambda^m\).  Hence
\[
  S_2^\sharp(m)\le B^2C(\sigma^2\lambda)^m,
\]
so the second covered moment is exponentially bounded.  The contrapositive gives
the stated lower bound: a cover whose second moment exceeds every exponential
bound must have \(H_m^\sharp/M_m\) exceeding every exponential bound.

For the one-spike witness of \Cref{thm:L2-finite-fiber-no-transfer-content} with
\(a_m=m!\), the maximal cover fibre size is \(H_m^\sharp=m!+1\).  The exact height
law gives \(M_m=\Theta(\varphi^{m/2})\).  Hence there is a constant \(C_0>0\) such
that \(M_m\le C_0\varphi^{m/2}\) for all \(m\).  For any proposed exponential bound
\(C\rho^mM_m\) on \(H_m^\sharp\), choose \(D=C C_0\) and
\(\eta=\rho\varphi^{1/2}\).  Since \(m!/\eta^m\to\infty\), eventually
\(m!>D\eta^m\), and therefore
\[
  H_m^\sharp=m!+1>C\rho^mM_m .
\]
Thus the factorial cover has cover-height excess over the intrinsic maximum that
is not exponentially bounded.  This proves that the bad cover is separated from
the ordinary \(\Fold_m\)-height scale by the exact intrinsic maximum formula, not
by an uncited histogram or thermodynamic input.
\end{proof}


% End original section: sec04_intrinsic_pressure.tex
% Original section: sec05_finite_fibre_strictness.tex
\subsection{Finite-fibre transfer strictness}

\begin{proposition}[Finite-fibre certificates and transfer obstruction]
\label{prop:L2-generic-vs-structural}
Finite-fibre certificates are fixed-resolution finite-atomic data.  A
finite-state transfer formula in the resolution variable is extra structure.
\end{proposition}

\begin{proof}
This is \Cref{prop:L2-novelty-finite-state-separation}.
\end{proof}

\begin{definition}[Computable finite-fibre cover]
\label{def:computable-cover}
A computable finite-fibre cover of the Zeckendorf resolution data is auxiliary
finite certificate data over the fixed normal-form set \(X_m\): over each normal
form one attaches a computably specified nonempty finite set, without changing
\(X_m\), the original maps \(\Fold_m\), or any readout on normal forms.
\end{definition}

\begin{lemma}[Covers preserve readouts]
\label{lem:cover-preserves-protocol}
Computable finite-fibre covers preserve the Zeckendorf normal forms and the
compiled \(L1\) readouts.
\end{lemma}

\begin{proof}
The cover changes only the number of certificate points over a fixed normal
form.  Projection to the normal form is unchanged, so every operation and
readout defined on normal forms is unchanged.
\end{proof}

\begin{theorem}[Certificate--transfer dichotomy for \texorpdfstring{\(L2\)}{L2}]
\label{thm:L2-finite-state-novelty-spine}
Fix the Zeckendorf normal-form system \((X_m,\Fold_m)\).  Finite-fibre
certificate universality at level \(L2\) is fixed-resolution finite-atomic data;
a fixed finite-state transfer law in the resolution parameter is additional
cross-resolution structure.  More precisely, in the computable finite-fibre
cover category of \Cref{def:computable-cover}:
\begin{enumerate}
\item a cover is exactly a computable positive integer fibre-size function
  \(d_m:X_m\to\mathbf N_{\ge1}\), and projection to \(X_m\) preserves all
  Zeckendorf normal forms and compiled \(L1\) readouts;
\item if the second covered collision moments
  \[
    M_m(d)=\sum_{x\in X_m} d_m(x)^2
  \]
  are governed by a fixed finite-state transfer formula, then
  \(M_m(d)\le C\lambda^m\) for some constants \(C,\lambda>0\);
\item for every finite window \(N\), every rational \(A>0\), and every rational
  \(\rho>1\), there is a computable cover, equal to the singleton cover for
  \(m<N\), whose moments satisfy \(M_m(d)>A\rho^m\) for all \(m\ge N\) and are
  not governed by any fixed finite-state transfer formula.
\end{enumerate}
Thus the \(L2\) novelty proved here is finite-fibre certificate universality;
finite-state transfer is a rigid strengthening that can be imposed for the
intrinsic fold moments of \Cref{thm:intrinsic-fold-moment-transfer} but is not
forced by finite-fibre certificates or by any finite initial certificate window.
This is the certificate/transfer boundary used by
\Cref{thm:finite-state-strictness,thm:L2-cover-relative-closure} in the
\(L2\) clause of \Cref{thm:main-hierarchy}.
\end{theorem}

\begin{proof}
For a fixed resolution, \Cref{prop:classical-finite-atomic-skeleton} identifies
finite-fibre certificates with finite atomic certificate data.  In the cover
category this datum is exactly a positive integer fibre-size function over each
normal form: one direction counts the added finite certificate fibre over
\(x\in X_m\), and the other attaches \(\{1,\ldots,d_m(x)\}\) over \(x\).  The
constructions are computable and inverse up to relabelling of certificate
points.  The preservation of normal forms and \(L1\) readouts is
\Cref{lem:cover-preserves-protocol}.  This proves the first assertion.

For the second assertion, a fixed finite-state transfer formula has only
finitely many states and fixed rational transition and output data.  By
\Cref{prop:finite-state-coefficient-budget,prop:fixed-transfer-interface-budget},
every sequence produced by such a fixed transfer formula is bounded by
\(C\lambda^m\) in the resolution parameter for suitable constants \(C\) and
\(\lambda\).

It remains to prove the quantitative non-forcing assertion.  Choose a computable
distinguished normal form \(x_m\in X_m\), for instance the lexicographically first
normal form.  Put \(d_m\equiv1\) for \(m<N\).  For \(m\ge N\), put
\(d_m(x)=1\) for \(x\ne x_m\) and
\[
  d_m(x_m)=1+m!+\left\lceil \sqrt{A\rho^m}\right\rceil .
\]
This is a computable finite-fibre cover of the unchanged normal-form system.
For \(m\ge N\), its second covered moment is
\[
  M_m(d)=\bigl(1+m!+\lceil\sqrt{A\rho^m}\rceil\bigr)^2+|X_m|-1>A\rho^m.
\]
Moreover \(M_m(d)\ge (m!)^2\) for all large \(m\).  If the sequence had a fixed
finite-state transfer formula, the previous paragraph would give
\(M_m(d)\le C\lambda^m\) for some \(C,\lambda\), contradicting the elementary
fact that \((m!)^2/(C\lambda^m)\to\infty\).  Hence no such transfer formula
exists.  Since the construction is the singleton cover on the prescribed window
and every cover projects to the same unmodified \((X_m,\Fold_m)\), finite-state
transfer is not determined by finite-fibre certificates, by finite prefix data,
or by the intrinsic fold map without an additional comparison theorem.  The
positive intrinsic transfer theorem is therefore a separate strengthening, not
the definition of \(L2\).
\end{proof}

\begin{theorem}[Computable covers are computable fibre-size data]
\label{thm:finite-fiber-cover-realization}
Specifying a computable finite-fibre cover is equivalent to specifying a
computable positive integer fibre size over each normal form at each resolution.
\end{theorem}

\begin{proof}
Given a cover, count the finite certificate fibres effectively.  Given
computable positive sizes, attach the auxiliary certificate fibre
\(\{1,\ldots,d\}\) over the corresponding normal form.  These
constructions are inverse up to relabelling of certificate points.
\end{proof}

\begin{theorem}[Extremal cover spine used for \texorpdfstring{\(L2\)}{L2} strictness]
\label{thm:L2-cover-skeleton-strictness}
Fix the Zeckendorf normal-form system \((X_m,\Fold_m)\).  In the class of
computable finite-fibre covers of \Cref{def:computable-cover}, the
finite-state-transfer obstruction proved in this article is witnessed, and at
fixed certificate excess is maximized, by the computable fibre-size function over
the fixed sets \(X_m\).  More precisely, the cover-relative strictness theorem
proved here follows from the following three-part spine:
\begin{enumerate}
\item finite-state transfer sequences are exponentially bounded in the
  resolution parameter by the fixed-transfer budget of
  \Cref{prop:fixed-transfer-interface-budget}; and
\item among covers with a prescribed total excess \(E_m\) above singleton fibres
  at resolution \(m\), the second covered collision moment is at most
  \((E_m+1)^2+|X_m|-1\), with equality precisely for a one-spike cover; and
\item a computable cover may realize this extremal one-spike fibre over a fixed
  normal form at resolution \(m\) for any prescribed computable excess \(E_m\).
\end{enumerate}
Moreover, the same one-spike skeleton can be chosen to agree with any prescribed
finite certificate window before it violates the transfer budget.
No assertion about the intrinsic fibre multiplicities of the original maps
\(\Fold_m\) is used or implied by this skeleton.
\end{theorem}

\begin{proof}
The first part is \Cref{prop:fixed-transfer-interface-budget}: by
\Cref{def:fixed-finite-state-transfer,prop:finite-state-coefficient-budget}, any
fixed finite-state transfer formula gives a sequence bounded by
\(C\lambda^m\) for some constants \(C,\lambda\).  For the extremal assertion,
write the cover sizes at resolution \(m\) as \(1+e_x\), where
\(e_x\ge0\) and \(\sum_x e_x=E_m\).  Then
\[
  \sum_x(1+e_x)^2=|X_m|+2E_m+\sum_xe_x^2
  \le |X_m|+2E_m+E_m^2=(E_m+1)^2+|X_m|-1,
\]
and equality holds exactly when all nonzero excess is concentrated at one normal
form.  For the realizing cover, use \Cref{thm:finite-fiber-cover-realization}:
choose one normal form in each
nonempty \(X_m\), assign its cover fibre size \(E_m+1\), and assign size \(1\) to
all other cover fibres.  This is a computable finite-fibre cover, and by
\Cref{lem:cover-preserves-protocol} the projection to \(X_m\), the normal forms,
and all \(L1\) readouts are unchanged.  Its second collision moment is therefore
the extremal value \((E_m+1)^2+|X_m|-1\); taking a superexponential computable
sequence such as \(E_m=m!\) violates the exponential budget and gives
\Cref{thm:finite-state-strictness}.  If an initial certificate window is compared
with the singleton cover, take fibre size \(1\), equivalently excess \(E_m=0\),
on that window and make the same superexponential choice afterwards; this is
exactly the finite-window obstruction of
\Cref{thm:L2-finite-window-transfer-obstruction}.  Since the construction
changes only cover certificate multiplicities above already fixed normal forms,
it makes no claim about the unmodified fibres of \(\Fold_m\).
\end{proof}

\begin{theorem}[Quantitative finite-state budget obstruction]
\label{thm:L2-quantitative-budget-obstruction}
Let \(A>0\) and \(\rho>1\) be rational numbers, and let \(N\ge1\).  There is a
computable finite-fibre cover of the Zeckendorf normal-form system with the
following properties.
\begin{enumerate}
\item For every \(m<N\) it is the singleton cover, so its finite-fibre
  certificate multiset and all \(L1\) readouts agree with the unmodified
  normal-form presentation at those resolutions.
\item For every \(m\ge N\), its second covered collision moment is larger than
  \(A\rho^m\).
\item The cover changes only one auxiliary certificate fibre over the fixed set
  \(X_m\) at each resolution.  Hence the budget violation is purely
  cover-relative and gives no assertion about the intrinsic multiplicities of
  \(\Fold_m\).
\end{enumerate}
Consequently no finite initial certificate window, and no effective rational
exponential finite-state transfer budget, can characterize \(L2\) finite-fibre
certificates inside the computable cover category.
\end{theorem}

\begin{proof}
For \(m<N\), put singleton fibres over every element of \(X_m\).  For
\(m\ge N\), choose one normal form \(x_m^\ast\in X_m\) and put
\[
  E_m=\left\lceil \sqrt{A\rho^m}\right\rceil,
\]
then attach a fibre of size \(E_m+1\) over \(x_m^\ast\) and singleton fibres over
all other normal forms.  The function \(m\mapsto E_m\) is computable from the
fixed rational constants as an integer-valued sequence; equivalently, compute
\(E_m\) as the least integer \(E\) with \(E^2\ge A\rho^m\).  By
\Cref{thm:finite-fiber-cover-realization} this fibre-size function defines a
computable finite-fibre cover, and by \Cref{lem:cover-preserves-protocol} it
preserves all normal forms and all \(L1\) readouts.

For \(m<N\), the construction is singleton by definition.  For \(m\ge N\), the
second covered moment is
\[
  (E_m+1)^2+|X_m|-1.
\]
Since \(|X_m|\ge1\) and \(E_m^2\ge A\rho^m\), this is strictly larger than
\(A\rho^m\).  Only the single auxiliary fibre over \(x_m^\ast\) is enlarged;
the original sets \(X_m\) and maps \(\Fold_m\) are not changed.  Thus the
construction defeats the specified exponential transfer budget while remaining
entirely in the cover category.
\end{proof}

\begin{theorem}[Main-article boundary for finite-fibre refinements]
\label{thm:L2-main-article-boundary}
Only computable finite-fibre cover data that leave normal forms fixed are used
in this article's \(L2\) strictness theorem.  The unmodified Zeckendorf
\(\Fold_m\) fibre multiplicities are covered instead by the intrinsic transfer
and height package of
\Cref{thm:intrinsic-fold-moment-transfer,thm:intrinsic-second-moment-recurrence,cor:polynomial-intrinsic-fiber-statistics,thm:intrinsic-fold-fibre-height,cor:linfty-fold-entropy,thm:zero-temperature-fold-pressure}.
No full intrinsic fibre-histogram law or large-deviation theorem is used in the
cover strictness chain.
\end{theorem}

\begin{proof}
All results here use \Cref{def:computable-cover}, whose projection to
\(X_m\) is fixed.  The cover-size function is independent finite certificate data
over the already chosen normal forms, while the intrinsic multiplicities are the
fibre sizes of the unmodified maps \(\Fold_m\).  By
\Cref{thm:L2-cover-intrinsic-nonimplication}, a transfer statement or obstruction
for the cover moment has no consequence for the intrinsic fold moment without an
additional comparison theorem.  The intrinsic fold moment and height endpoint are
handled directly by
\Cref{thm:intrinsic-fold-moment-transfer,thm:intrinsic-fold-fibre-height,cor:linfty-fold-entropy,thm:zero-temperature-fold-pressure}; the cover strictness proof supplies no
full intrinsic histogram or large-deviation classification.
\end{proof}

\begin{theorem}[Computable-cover-only boundary of finite-state obstruction]
\label{prop:cover-only-strictness-boundary}
\label{thm:one-spike-finite-state-boundary}
The one-spike construction is sharp for the computable finite-fibre cover model:
it preserves all normal-form readouts and changes only one added finite
certificate fibre at each resolution, yet it can realize any computable
second-moment spike of the form \((a_m+1)^2+F_{m+2}-1\).  The theorem is silent
about intrinsic \(\Fold_m\) fibre multiplicities.
\end{theorem}

\begin{proof}
This is the construction in \Cref{thm:L2-finite-fiber-no-transfer-content}.
Only one fibre is enlarged, and \Cref{lem:cover-preserves-protocol} proves that
readouts are preserved.
\end{proof}

\begin{theorem}[Intrinsic height is separate from cover spikes]
\label{thm:intrinsic-fold-asymptotics-split}
The factorial one-spike obstruction does not assert a full intrinsic histogram or
large-deviation law for the original Zeckendorf fibres.  It is a statement about
computable finite-fibre covers, so no such classification theorem for
\(\Fold_m\) fibre multiplicities follows from
\Cref{thm:finite-state-strictness}.  The positive intrinsic fixed-\(q\) transfer,
recurrence, height, and zero-temperature endpoint statements used in the article
are exactly those proved in
\Cref{thm:intrinsic-fold-moment-transfer,thm:intrinsic-second-moment-recurrence,cor:polynomial-intrinsic-fiber-statistics,thm:intrinsic-fold-fibre-height,cor:linfty-fold-entropy,thm:zero-temperature-fold-pressure}.
\end{theorem}

\begin{proof}
The proof of \Cref{thm:L2-finite-fiber-no-transfer-content} deliberately alters
cover multiplicities.  It leaves the original fold map untouched.  The formal
nonimplication is \Cref{thm:L2-cover-intrinsic-nonimplication}.
\end{proof}


\begin{theorem}[Cover-relative closure of the \texorpdfstring{\(L2\)}{L2} strictness clause]
\label{thm:L2-cover-relative-closure}
The \(L2\) strictness clause used by \Cref{thm:main-hierarchy} is fully proved
inside the computable finite-fibre cover category of \Cref{def:computable-cover}.
It neither requires nor implies any intrinsic nonrecurrence, finite-state-failure,
or asymptotic growth classification for the multiplicity sequence of the original
maps \(\Fold_m\).  Positive intrinsic finite-state transfer for those original
maps is supplied separately by \Cref{thm:intrinsic-fold-moment-transfer}.
\end{theorem}

\begin{proof}
The positive finite-fibre part of \(L2\) is fixed-resolution certificate data:
\Cref{prop:classical-finite-atomic-skeleton,prop:zeckendorf-cert-univ} give the
finite-atomic certificates from finite fibres at each fixed resolution.  A
finite-state transfer law is additional cross-resolution
structure.  By
\Cref{def:fixed-finite-state-transfer,prop:fixed-transfer-interface-budget},
every fixed finite-state transfer sequence has an exponential bound in the
resolution parameter.

Now work in the cover category.  By \Cref{thm:finite-fiber-cover-realization}, a
computable cover is exactly a computable positive integer fibre-size function
over the unchanged normal-form sets \(X_m\), and by
\Cref{lem:cover-preserves-protocol} such a cover preserves the Zeckendorf normal
forms and all \(L1\) readouts.  The one-spike skeleton of
\Cref{thm:L2-cover-skeleton-strictness} gives the extremal obstruction at each
fixed certificate excess, and the quantitative budget form
\Cref{thm:L2-quantitative-budget-obstruction} says more: after any prescribed
finite window, and against any effective rational exponential bound \(A\rho^m\), a
computable one-spike cover preserves the protocol while forcing the second
covered moment above that bound.  In particular, taking a superexponential
computable spike such as \(m!\) gives a computable cover whose second collision
moment violates every exponential bound, and therefore cannot come from a fixed
finite-state transfer formula in the sense of \Cref{def:fixed-finite-state-transfer}.
This proves the cover-relative strictness theorem.

The proof has not used the intrinsic fibre sizes of the original maps
\(\Fold_m\).  Indeed \Cref{thm:L2-cover-intrinsic-nonimplication} records the
formal nonimplication: altering auxiliary cover multiplicities over fixed normal
forms gives no comparison theorem for unmodified fold fibres.  Hence the main
article's \(L2\) strictness is exactly cover-relative.  The positive intrinsic
transfer package is the separate theorem package in \Cref{sec:L2}, not a
consequence of the one-spike cover obstruction.
\end{proof}

\begin{theorem}[Cover-relative finite-fibre certificates do not force finite-state transfer]
\label{thm:finite-state-strictness}
\label{thm:cover-relative-finite-fibre-certificates-do-not-force-finite-state-transfer}
There is a computable finite-fibre cover of the Zeckendorf \(L1\) protocol whose
finite-fibre certificates exist at every resolution but whose second collision
moment over the fixed normal-form sets \(X_m\) has no fixed finite-state transfer
formula in the sense of \Cref{def:fixed-finite-state-transfer}.  This is a
cover-relative strictness theorem in the computable
finite-fibre cover category; it is not an intrinsic finite-state failure theorem
for the original \(\Fold_m\) fibre multiplicities and does not contradict the
intrinsic height law of \Cref{thm:intrinsic-fold-fibre-height}.  More sharply,
the obstruction is quantitative: after any prescribed finite certificate window
and against any effective rational exponential transfer budget, a one-spike computable
cover can preserve all normal forms and all \(L1\) readouts while forcing the
second covered moment beyond that budget.  Over the same unmodified Zeckendorf
data \((X_m,\Fold_m)\) there are therefore computable covers whose second
collision moments have opposite finite-state-transfer behaviour, as isolated in
\Cref{prop:cover-category-invariance-finite-state-obstruction}.  Hence the
obstruction is not a property determined by the intrinsic fold map alone.
\end{theorem}

\begin{proof}
Apply the cover skeleton \Cref{thm:L2-cover-skeleton-strictness} with
\(a_m=m!\), equivalently apply \Cref{thm:L2-finite-fiber-no-transfer-content} to
that one-spike sequence.  The cover is computable and preserves readouts by
\Cref{lem:cover-preserves-protocol}.  Its second moment grows faster than every
exponential, whereas every fixed finite-state transfer sequence is exponentially
bounded by \Cref{prop:fixed-transfer-interface-budget}.  The obstruction is
therefore cover-relative: it concerns the added finite certificate fibres over
the unchanged normal-form sets, not intrinsic multiplicities of the original
\(\Fold_m\) fibres.

The sharper two-cover assertion is
\Cref{prop:cover-category-invariance-finite-state-obstruction,thm:L2-two-cover-intrinsic-boundary}.
Its singleton cover has moment \(F_{m+2}\), which has a fixed two-state Fibonacci
transfer formula, while its factorial one-spike cover has moment
\((m!+1)^2+F_{m+2}-1\), which violates the exponential budget just cited.  Since
both covers project to the same unmodified pair \((X_m,\Fold_m)\), finite-state
transfer or failure for these cover moments cannot be read off from the intrinsic
fold map without an additional bridge identifying intrinsic fold multiplicities
with the chosen cover multiplicities.  Such a bridge is outside the theorem chain
by \Cref{thm:L2-cover-intrinsic-nonimplication}.
\end{proof}

\begin{proposition}[Finite prefixes do not certify transfer]
\label{prop:finite-prefix-no-transfer-certification}
In the computable finite-fibre cover category, no bounded initial window of
finite-fibre certificate data can certify finite-state transfer for all later
resolutions.
\end{proposition}

\begin{proof}
Given any bound \(M\), choose a one-spike cover with \(a_m=1\) for
\(m\le M\) and \(a_m=m!\) for \(m>M\).  It agrees with the unspiked data on the
initial window but violates the finite-state exponential budget later.
\end{proof}

\begin{theorem}[Budget spine used for \texorpdfstring{\(L2\)}{L2} strictness]
\label{thm:L2-budget-spine}
The cover-relative \(L2\) strictness obstruction proved in this article is a
growth-budget obstruction: finite-state transfer sequences obey the exponential
bound of \Cref{prop:fixed-transfer-interface-budget}, while computable
finite-fibre covers can violate it, and the superexponential cover obstructions
used here have the one-spike skeleton of
\Cref{thm:L2-cover-skeleton-strictness}.
\end{theorem}

\begin{proof}
The budget is \Cref{prop:fixed-transfer-interface-budget}, which applies the
definition \Cref{def:fixed-finite-state-transfer} through
\Cref{prop:finite-state-coefficient-budget}; the violating cover is
\Cref{thm:finite-state-strictness}.  The one-spike bad-example skeleton used for
the cover-relative strictness theorem is \Cref{thm:L2-cover-skeleton-strictness}.
\end{proof}

\begin{corollary}[Finite-state growth budget]
\label{cor:finite-state-growth-budget}
Every fixed finite-state transfer formula gives at most exponential growth in
the resolution parameter.
\end{corollary}

\begin{proof}
This is the exponential-bound clause of
\Cref{prop:fixed-transfer-interface-budget}.
\end{proof}




% End original section: sec05_finite_fibre_strictness.tex
% Original section: sec05_hierarchy_strictness.tex
\section{Strictness of the hierarchy}\label{sec:strictness}

This section proves the two strictness mechanisms used by
\Cref{thm:main-hierarchy}.  The primary Richardson theorem is the verified-triple
statement: any finite pure EML triple whose \(\pi\), sine, and punctured
absolute-value rows have been proved transports Richardson's undecidable
identity-to-zero problem.  The paper proves the punctured absolute-value row
internally; the printed pair \((U_{47},U_{111})\) appears only as the conditional
cited specialization \((H_{\mathcal W})\) from \Cref{ass:HW}, and the
endpoint-root string \(U_{126}\) is not a row of the Richardson package.

\subsection{The Richardson obstruction}

Let \(\mathcal R\) be Richardson's real elementary language from
\cite[pp.~514, 520, Theorem~2]{Richardson1968}: one real variable, rational
constants, \(\log2\), \(\pi\), \(+,-,\times\), composition, \(\exp\), \(\sin\),
and a unary symbol \(\mu\) satisfying \(\mu(x)=|x|\) for \(x\ne0\).  We use this
absolute-value primitive exactly in that punctured sense.  For an expression
\(E\), let \(D_E\) denote the recursively defined real domain on which all
ordinary elementary subterms are defined and, whenever a subterm \(\mu(F)\) is
evaluated, the value \(F(x)\) is nonzero.

This deletion convention is part of the source problem used below.  It does not
change any constant or operation in Richardson's language; it only makes explicit
that the reduction never asks for a value of the auxiliary absolute-value symbol
at an argument equal to zero.

\begin{theorem}[Richardson's punctured identity-to-zero theorem]
\label{thm:richardson-punctured-source-problem}
\label{lem:richardson-punctured-mu-reduction}
\label{lem:richardson-fragment}
The following identity-to-zero decision problem is undecidable: given a
Richardson expression \(E\) built from exactly the constants and operations listed
above, with its recursively determined partial real domain \(D_E\), determine
whether
\[
  E(x)=0\qquad\hbox{for every }x\in D_E .
\]
This is the partial-domain form of Richardson's identity problem needed below:
identity means equality of partial real functions on the displayed source domain,
and the value of \(\mu\) at zeros of its argument is never queried.
\end{theorem}

\begin{proof}
Richardson works with partial real functions and takes equality of expressions to
mean equality on their common domains \cite[p.~514]{Richardson1968}.  His
condition (2) assumes a unary function \(\mu\) satisfying only
\(\mu(t)=|t|\) for \(t\ne0\) \cite[p.~514]{Richardson1968}.  We recall the short
reduction from the less-than problem in his proof, because it is exactly the
punctured convention used here.

By Richardson's Theorem~1, under condition (1) there is no algorithm deciding,
for an expression \(A\), whether \(A(x)<0\) for some real \(x\)
\cite[p.~520, Theorem~1]{Richardson1968}.  In the proof of Theorem~2 he applies
this to the effectively produced subelementary functions \(G(n,x)\), and forms
\[
  B_n(x)=\mu(G(n,x)-1)-(G(n,x)-1).
\]
Let
\[
  D_n=\{x\in\operatorname{dom}G(n,\cdot):G(n,x)-1\ne0\}.
\]
On \(D_n\), condition (2) gives
\[
  B_n(x)=|G(n,x)-1|-(G(n,x)-1).
\]
Hence \(B_n(x)=0\) for all \(x\in D_n\) if and only if no real \(x\) satisfies
\(G(n,x)<1\): if \(G(n,x)>1\), the displayed expression is zero, while if
\(G(n,x)<1\), it is \(2(1-G(n,x))>0\).  Points with \(G(n,x)=1\) are precisely
the points deleted from \(D_n\), so no value of \(\mu(0)\) is used.  Richardson's
Theorem~2 proves that deciding the corresponding identities would decide this
undecidable less-than problem \cite[p.~520, Theorem~2]{Richardson1968}.  Since the
map \(n\mapsto B_n\) is effective, a decision procedure for the displayed
punctured identity-to-zero problem would decide Richardson's undecidable problem,
which is impossible.
\end{proof}

\begin{proposition}[Stable Richardson source check]
\label{prop:stable-richardson-source-check}
The punctured domain problem in
\Cref{thm:richardson-punctured-source-problem} is exactly the form of
Richardson's Theorem~2 used in this paper.  The source treats expressions as
partial real functions on common domains, assumes the language contains the
constants and closure operations listed in condition~(1), and in condition~(2)
requires only a unary function \(\mu\) with
\(\mu(x)=|x|\) for \(x\ne0\) \cite[pp.~514, 520, Theorem~2]{Richardson1968}.
Consequently an identity algorithm for the recursive domains \(D_E\) defined
above would decide Richardson's theorem, while no value of \(\mu(0)\) is part of
the reduction.
\end{proposition}

\begin{proof}
The bibliographic record for \cite{Richardson1968} pins the article by DOI,
journal, volume, issue, and pages \(514\)--\(520\); the Cambridge Core metadata
for the same DOI identifies the article as Richardson's paper and records those
pages together with the abstract's partial-function convention.  On p.~514,
Richardson sets up a class of expressions denoting real single-valued partially
defined functions and compares expressions on their common domains.  The same
page states the two hypotheses used for Theorem~2: condition~(1) gives the
rational constants, \(\log2\), \(\pi\), arithmetic closure, composition,
\(\exp\), and \(\sin\); condition~(2) adds a function \(\mu\) satisfying
\(\mu(x)=|x|\) for nonzero \(x\).  These are exactly the primitives listed in
\(\mathcal R\), except that this manuscript makes the common-domain convention
explicit by deleting from \(D_E\) every point at which a \(\mu\)-argument is zero.

The reduction in the proof of
\Cref{thm:richardson-punctured-source-problem} is therefore only a spelling-out
of Richardson's Theorem~2 under this explicit domain notation.  For the terms
\(B_n(x)=\mu(G(n,x)-1)-(G(n,x)-1)\), all evaluations occur on the set where
\(G(n,x)-1\ne0\).  On that set condition~(2) gives the ordinary absolute value,
so the zero identity of \(B_n\) is equivalent to the corresponding less-than
instance used by Richardson.  If \(G(n,x)-1=0\), the point has been removed from
\(D_{B_n}\) and no value of \(\mu(0)\) is queried.  This is the only place where
this paper changes notation relative to the source: Richardson's common-domain
partial-function convention is retained, but written as the recursive punctured
sets \(D_E\).  The page and theorem references used for this interface are the
p.~514 definition-and-hypotheses passage and the p.~520 statement and proof of
Theorem~2; no stronger total-domain absolute-value theorem is being imported.
Hence a decision procedure for the displayed punctured identity problem would be
a decision procedure for the problem Richardson proves undecidable, and the
manuscript's \(D_E\) formulation is a conservative explicit-domain restatement of
the cited theorem.
\end{proof}

\begin{theorem}[Abstract punctured normalizer obstruction]
\label{thm:abstract-normalizer-obstruction}
Let \(\tau\colon\mathcal R\to\mathcal T\) be an effective translation into a
computable partial real term class such that, for every Richardson expression
\(E\), the translated term \(\tau(E)\) is defined on \(D_E\) and satisfies
\(\llbracket\tau(E)\rrbracket(x)=E(x)\) for all \(x\in D_E\).  Then the translated
source-indexed fragment
\[
  E\longmapsto (\tau(E),D_E,\tau(0))
\]
admits no total computable normalizer into finite objects with decidable equality
that is equality-reflecting and equality-preserving for the transported partial
real functions after restriction to the source domain \(D_E\): for every
Richardson expression \(E\), equality of the two finite normal forms for the
source-indexed objects \((\tau(E),D_E)\) and \((\tau(0),D_E)\) must hold if and
only if \(\llbracket\tau(E)\rrbracket=\llbracket\tau(0)\rrbracket\) on \(D_E\).
\end{theorem}

\begin{proof}
If such a normalizer existed, compute \(\tau(E)\) and compare its normal form
with that of \(\tau(0)\) inside the same source-indexed object carrying the domain
\(D_E\).  By the hypothesis on the normalizer, equality of these finite normal
forms is equivalent to
\(\llbracket\tau(E)\rrbracket(x)=\llbracket\tau(0)\rrbracket(x)\) for every
\(x\in D_E\).  By value preservation on \(D_E\), this is equivalent to
\(E(x)=0\) for every \(x\in D_E\), contradicting
\Cref{thm:richardson-punctured-source-problem}.
\end{proof}

\begin{lemma}[Transported zero-test obstruction]
\label{lem:transported-partial-normalizer-obstruction}
The conclusion of \Cref{thm:abstract-normalizer-obstruction} holds for every EML
translation that is defined and value-preserving on each punctured Richardson
domain \(D_E\).
\end{lemma}

\begin{proof}
Apply \Cref{thm:abstract-normalizer-obstruction} with \(\mathcal T\) equal to the
translated EML fragment.
\end{proof}

\begin{lemma}[Printed EML core witnesses]
\label{lem:eml-printed-core}
In the principal-branch EML semantics of \Cref{def:eml}, the following pure EML
macros represent the indicated real functions:
\[
  E(u)=\eml(u,1)=e^u,
  \qquad
  L(u)=\eml\bigl(1,E(\eml(1,u))\bigr)=\log u\quad(u>0),
\]
\[
  D(a,b)=\eml(L(a),E(b))=a-b\quad(a>0,
b\in\RR),
\]
\[
  P(u)=\eml(u,E(u))=e^u-u>0,
  \quad
  N(u)=D(P(u),E(u))=-u,
\]
\[
  A(u,v)=D(E(u),D(P(u),v))=u+v,
  \qquad S(u,v)=A(u,N(v))=u-v .
\]
\end{lemma}

\begin{proof}
The EML semantics are \(\eml(x,y)=e^x-\operatorname{Log}y\) when the principal
logarithm is defined.  Since the second input of \(E(u)\) is \(1\),
\(E(u)=e^u\).  If \(u>0\), then \(\eml(1,u)=e-\log u\) and \(E(\eml(1,u))=e^e/u>0\), so
\(L(u)=e-\log(e^e/u)=\log u\).  Thus
\(D(a,b)=\exp(\log a)-\log(e^b)=a-b\) for \(a>0\).  The function
\(P(u)=e^u-u\) has minimum \(1\) at \(u=0\), hence is positive.  Substituting in
\(D\) gives \(N(u)=-u\), then \(A(u,v)=u+v\), and finally \(S(u,v)=u-v\).
\end{proof}

\begin{theorem}[Self-contained core EML witness package]
\label{lem:eml-real-witness-library}
The macros of \Cref{lem:eml-printed-core} extend internally to finite pure EML
terms representing multiplication, division by a nonzero real, rational
constants, \(\log2\), and the positive square root:
\[
\begin{aligned}
  T(u,v)&=E(A(L(u),L(v)))=uv\quad(u,v>0),\\
  R_+(u)&=E(N(L(u)))=1/u\quad(u>0),
\end{aligned}
\]
\[
  Q(u)=E(u)>0,
  \qquad P(u)=E(u)-u>0,
\]
\[
\begin{aligned}
  M(u,v)=S\Bigl(&A\bigl(T(Q(u),Q(v)),T(P(u),P(v))\bigr),\\
  &A\bigl(T(Q(u),P(v)),T(P(u),Q(v))\bigr)\Bigr)=uv,
\end{aligned}
\]
for all real \(u,v\).  Indeed all four inputs to \(T\) in this displayed
definition are positive.  With this expanded pure-EML multiplication term,
\[
  J(u,v)=M\bigl(u,M(v,R_+(M(v,v)))\bigr)=u/v\quad(v\ne0),
\]
and
\[
  G(u)=E(J(L(u),A(1,1)))=\sqrt u\quad(u>0).
\]
The value at \(0\) for the nonnegative square root, the constant \(\pi\), and
\(\sin\) on \(\RR\) are not supplied by this core package.
\end{theorem}

\begin{proof}
For positive inputs, \(T(u,v)=\exp(\log u+\log v)=uv\) and
\(R_+(u)=\exp(-\log u)=1/u\).  Since \(Q(u)>0\) and \(P(u)>0\), all four calls
to \(T\) in the definition of \(M\) are made on its positive domain.  Using
\(Q(u)-P(u)=u\), the displayed term gives
\[
  M(u,v)=Q(u)Q(v)+P(u)P(v)-Q(u)P(v)-P(u)Q(v)=uv .
\]
If \(v\ne0\), then \(M(v,v)=v^2>0\), so
\(J(u,v)=u\,v/v^2=u/v\).  Finally
\(G(u)=\exp((\log u)/2)=\sqrt u\) for \(u>0\).  Rational constants are obtained
from \(1\) by repeated addition, negation, and division by nonzero integers;
\(\log2=L(A(1,1))\).
\end{proof}

\begin{theorem}[Punctured Richardson root from the positive EML core]
\label{thm:positive-core-nonnegative-root-boundary}
\label{thm:eml-punctured-richardson-core}
The self-contained core EML package supplies exactly the square-root input needed
for Richardson's theorem in its punctured absolute-value form.  More precisely,
let \(G(t)\) be the positive-domain term of
\Cref{lem:eml-real-witness-library}.  Then the pure EML macro
\[
  \Abs(u)=G(M(u,u))
\]
is defined for every real \(u\ne0\) and satisfies
\[
  \Abs(u)=|u|\qquad(u\ne0).
\]
No value at \(u=0\) is asserted.  Thus the Richardson reduction used in this
paper requires only a punctured absolute-value witness, not a pure EML term for
\(t\mapsto\sqrt t\) at the endpoint \(t=0\).
\end{theorem}

\begin{proof}
By \Cref{lem:eml-real-witness-library}, the multiplication macro satisfies
\(M(u,u)=u^2\) for all real \(u\), and the square-root macro \(G(t)\) is defined
and equal to \(\sqrt t\) for \(t>0\).  If \(u\ne0\), then \(u^2>0\), so the
composition \(G(M(u,u))\) is in the proved positive domain of \(G\).  Therefore
\[
  \Abs(u)=G(M(u,u))=\sqrt{u^2}=|u|\qquad(u\ne0).
\]
At \(u=0\) the inner value is \(M(0,0)=0\), outside the proved domain of \(G\), so
no endpoint value is supplied.  Richardson's condition (2), recalled in
\Cref{lem:richardson-punctured-mu-reduction}, asks only for a function agreeing
with \(|x|\) when \(x\ne0\).  Hence the punctured witness above is sufficient for
Richardson's identity theorem and avoids the impossible endpoint-square-root
requirement.
\end{proof}

\begin{lemma}[Conditional printed EML witnesses for \texorpdfstring{\(\pi\)}{pi} and sine]
\label{lem:eml-pi-sine-witnesses}
Let \(W_{\pi}=U_{47}\) and \(W_{\sin}=U_{111}\) be the finite pure EML strings
printed in \texttt{ancillary/eml\_witness\_catalogue.txt}.  Under the external
hypothesis \((H_{\mathcal W})\) of \Cref{ass:HW}, these two strings have the
real-domain identities
\[
  W_{\pi}=\pi,
  \qquad
  W_{\sin}(x)=\sin x\quad(x\in\RR).
\]
Without \((H_{\mathcal W})\), the catalogue and its checker certify only finite
pure-EML syntax and do not prove these two semantic identities.
\end{lemma}

\begin{proof}
The catalogue gives finite strings in the pure EML grammar, and the companion
checker \texttt{ancillary/check\_eml\_witnesses.py} verifies, by a bounded parse,
that the two displayed right-hand sides are finite expressions over the single
binary operation \(\eml\), the constant \(1\), and the input leaf for the unary
sine witness.  The checker explicitly records that it proves no real-domain
semantic identities.  The displayed equalities are therefore exactly the first
two clauses of \((H_{\mathcal W})\), restricted to \(U_{47}\) and \(U_{111}\).
The third printed row \(U_{126}\) from \Cref{ass:HW} is not used here and is not
transported into the Richardson triple, because the article's punctured
absolute-value row is the internally proved term \(G(M(u,u))\) of
\Cref{thm:eml-punctured-richardson-core}.  No conclusion about \(\pi\), sine, or
the endpoint value of a square root is drawn from finite syntax alone.
\end{proof}

\begin{theorem}[Printed Richardson gate rigidity]
\label{thm:printed-richardson-gate-rigidity}
\label{prop:printed-pi-sine-instantiation-criterion}
\label{prop:no-unconditional-printed-pi-sine-instance}
For the exact printed strings \(U_{47}\) and \(U_{111}\), the Richardson
obstruction in
\Cref{thm:richardson-witness-row-rigidity,thm:separation,thm:proc-ams-relative-richardson-theorem}
has the concrete printed-string specialization if and only if the two semantic
identities
\[
  U_{47}=\pi,
  \qquad
  U_{111}(x)=\sin x\quad(x\in\RR)
\]
are explicit hypotheses or are proved for those exact strings in the
principal-branch EML semantics.  Equivalently, the printed specialization is
conditional under \((H_{\mathcal W})\) of \Cref{ass:HW}, or under an independent
proof of exactly those two displayed identities; it is not an unconditional
consequence of this manuscript.  The finite catalogue and parser certificate
alone supply only membership in the pure EML grammar, so they do not by
themselves give an unconditional printed-string Richardson instance; nor does
the separate cited \(\mathcal B_{36}\) representative-existence import, which is
not a semantic verification of these two printed strings.  The printed endpoint
root \(U_{126}\) has no role in this gate.  Consequently every theorem,
corollary, or proof in this paper that names the concrete pair
\((U_{47},U_{111})\) as the \(\pi\)/sine rows of a Richardson triple is to be read
under \((H_{\mathcal W})\) or under a separate proof of exactly the two displayed
identities; without one of those inputs the statement reverts to the abstract
verified-triple interface.  This gate is rigid under the main theorem chain: the
finite syntax catalogue, the cited \(\mathcal B_{36}\) representative-existence
theorem, and the internally proved punctured absolute-value row cannot replace
either of the two printed real-domain identities.
In particular, a reference to \Cref{thm:separation} or
\Cref{thm:proc-ams-relative-richardson-theorem} alone never authorizes replacing
the abstract verified \(\pi\)- and sine-witness rows by the printed pair
\((U_{47},U_{111})\); that substitution is exactly the conditional branch just
described.
\end{theorem}

\begin{proof}
If the two displayed identities are supplied, combine them with the internally
proved punctured absolute-value identity
\(\Abs(u)=|u|\) for \(u\ne0\) from
\Cref{thm:eml-punctured-richardson-core}.  The triple
\((U_{47},U_{111},\Abs)\) then satisfies the three hypotheses of the relative
Richardson interface
\Cref{thm:replacement-triple-richardson-interface,thm:relative-finite-richardson-fragment},
and hence gives the printed-string specialization of the obstruction.

Conversely, \Cref{thm:richardson-witness-row-rigidity} proves that the concrete
printed-string specialization of that structural translation is value-preserving
only if the \(\pi\)-row and sine row hold for the substituted witnesses.  Thus
the use of \(U_{47}\) in the \(\pi\)-clause and \(U_{111}\) in the sine clause
requires exactly the first two displayed semantic hypotheses in
\Cref{lem:richardson-translation,thm:punctured-domain-richardson-reduction};
without them the induction proving value preservation has no base case for
\(\pi\) and no unary-function step for sine.  The ancillary catalogue records
only finite grammar strings, and the parser checks only that grammar.  Such a
syntactic certificate is compatible with many possible semantic values of a
formal string before principal-branch evaluation is proved, and therefore cannot
imply either real-domain identity.  Thus the exact printed strings instantiate
the Richardson obstruction precisely under the displayed semantic identities,
which are the \(U_{47},U_{111}\) part of \((H_{\mathcal W})\).  The external
\(\mathcal B_{36}\) import is deliberately not used in this converse: it gives a
version-pinned representative-existence assertion for the calculator basis, not
a proof that these two source-locked printed strings have the displayed
real-domain values.
The internal punctured absolute-value row supplies only the Richardson
\(\mu\)-clause, so it cannot repair a missing \(\pi\)-row or sine row.  Thus the
three possible inputs named in the statement affect disjoint primitive tests:
finite syntax gives only the existence of strings, the external
\(\mathcal B_{36}\) import gives only representative existence for the calculator
basis, and \(\Abs\) gives only \(|u|\) off zero.  None of these inputs proves the
two displayed printed identities, and the printed-pair branch is unavailable
exactly when those identities are unavailable.  Since
\Cref{thm:separation,thm:proc-ams-relative-richardson-theorem} are stated for
verified triples, citing them without the printed real-domain rows supplies only
the abstract scheme and not the concrete \((U_{47},U_{111})\) substitution.
\end{proof}

\begin{corollary}[No unconditional printed Richardson triple]
\label{cor:no-unconditional-printed-richardson-triple}
\label{thm:no-unconditional-concrete-eml-richardson-instance}
The manuscript contains no unconditional concrete Richardson instance with the
printed triple \((U_{47},U_{111},U_{126})\), nor with its punctured replacement
\((U_{47},U_{111},\Abs)\).  Unconditionally, the only concrete member of a
Richardson replacement triple proved here is the punctured absolute-value row
\(\Abs(u)=|u|\) for \(u\ne0\).  A concrete printed specialization using
\(U_{47}\) and \(U_{111}\) is obtained exactly after adding the two real-domain
identities \(U_{47}=\pi\) and \(U_{111}(x)=\sin x\) for all real \(x\).
\end{corollary}

\begin{proof}
The endpoint triple \((U_{47},U_{111},U_{126})\) would require the printed
\(\pi\)-identity, the printed sine identity, and an endpoint square-root identity.
The article proves none of these three identities from the printed catalogue:
\Cref{lem:eml-pi-sine-witnesses} states the \(U_{47}\) and \(U_{111}\) identities
only under \((H_{\mathcal W})\), and the endpoint-root identity is not used in the
punctured Richardson reduction.  The punctured replacement
\((U_{47},U_{111},\Abs)\) removes the endpoint-root requirement, but still leaves
exactly the two printed semantic identities for \(U_{47}\) and \(U_{111}\).
By \Cref{prop:printed-pi-sine-instantiation-criterion}, the catalogue and parser
certificate alone do not supply those two identities; they certify only finite
pure-EML syntax.  Conversely, once those two identities are supplied,
\Cref{thm:eml-punctured-richardson-core} supplies the third, punctured row
\(\Abs(u)=|u|\) for \(u\ne0\), so
\Cref{thm:replacement-triple-richardson-interface} gives the concrete printed
punctured specialization.  Hence the unconditional theorem in this manuscript is
the verified-triple interface together with the internally proved absolute-value
row, not a concrete printed \(U_{47}\)/\(U_{111}\) instance.
\end{proof}
\begin{lemma}[Quantitative core witness table]
\label{lem:eml-core-witness-table}
The EML part of the Richardson translation uses the following finite witness
table.  Each entry is a pure EML macro from
\Cref{lem:eml-printed-core,lem:eml-real-witness-library,lem:eml-pi-sine-witnesses};
the domain column is the exact real domain on which the identity is proved or,
for \(W_{\pi}\) and \(W_{\sin}\), externally assumed by \((H_{\mathcal W})\).
Equivalently, in the unconditional replacement-triple version one replaces only
the two conditional rows by finite pure EML terms \(\mathsf P\) and \(\mathsf S\)
whose identities \(\mathsf P=\pi\) and \(\mathsf S(x)=\sin x\) on \(\RR\) have
already been proved.
\[
\begin{array}{@{}c@{\quad}c@{\quad}c@{}}
\hbox{constructor} & \hbox{witness} & \hbox{domain}\\ \hline
0,1,\mathbb Q,\log2 &
\hbox{finite arithmetic from }1,\ L(A(1,1)) & \RR\hbox{ for constants}\\
u+v & A(u,v) & \RR^2\\
u-v & S(u,v) & \RR^2\\
uv & M(u,v) & \RR^2\\
u/v & J(u,v) & \{(u,v):v\ne0\}\\
e^u & E(u) & \RR\\
\pi & W_{\pi} & \RR^0\\
\sin u & W_{\sin}(u) & \RR\\
|u|\hbox{ off }0 & \Abs(u)=G(M(u,u)) & \RR\setminus\{0\}
\end{array}
\]
This is the complete finite witness table used by the conditional printed-string
translation and, after replacing the two conditional rows by any verified pair
\((\mathsf P,\mathsf S)\), by the unconditional replacement-triple translation in
\Cref{lem:richardson-translation,thm:richardson-translation-budget}.
The full nonnegative-square-root row \(\sqrt u\) on \([0,\infty)\) is not part of
the Richardson input in this paper.
\end{lemma}

\begin{proof}
The arithmetic, exponential, division, rational-constant, and \(\log2\) rows are
proved in \Cref{lem:eml-printed-core,lem:eml-real-witness-library}.  The \(\pi\)
and sine rows are exactly the conditional rows of
\Cref{lem:eml-pi-sine-witnesses}.  The punctured absolute-value row is proved in
\Cref{thm:eml-punctured-richardson-core}.  These are Richardson's condition (1)
and condition (2) as used in \Cref{lem:richardson-fragment} once the two
conditional rows have been supplied: the absolute-value symbol is required to
agree with \(|x|\) only away from zero.  If \(W_{\pi}\) and \(W_{\sin}\) are
replaced by any verified \((\mathsf P,\mathsf S)\), the same proof reads word for
word with those two proved identities in place of \((H_{\mathcal W})\).
\end{proof}

\begin{theorem}[Verified-triple Richardson instantiation]
\label{thm:verified-pair-richardson-instantiation}
\label{thm:verified-replacement-richardson-instantiation}
Let \((\mathsf P,\mathsf S)\) be finite pure EML terms satisfying
\[
  \mathsf P=\pi,
  \qquad \mathsf S(x)=\sin x\quad(x\in\RR)
\]
in the principal-branch real semantics.  With the internally proved punctured
absolute-value term
\[
  \mathsf A(u)=\Abs(u)=G(M(u,u)),
\]
the triple \((\mathsf P,\mathsf S,\mathsf A)\) gives a Richardson-undecidable pure
EML fragment only in the following identity-to-zero, source-domain sense.  More
precisely, there is an effective translation
\(\tau_{\mathsf P,\mathsf S,\Abs}\) from Richardson's punctured language to pure
EML terms such that, for every Richardson expression \(E\),
\[
  E=0\hbox{ on }D_E
  \quad\Longleftrightarrow\quad
  \tau_{\mathsf P,\mathsf S,\Abs}(E)=
  \tau_{\mathsf P,\mathsf S,\Abs}(0)
  \hbox{ on }D_E .
\]
Consequently this translated source-indexed fragment has no zero-identity
algorithm, no computable semantic normalizer into finite objects with decidable
equality that reflects and preserves equality of the transported partial real
functions on each displayed domain \(D_E\), and no equality-reflecting semantic
compiler into finite decidable protocol objects.  The exact printed pair
\((U_{47},U_{111})\) is one such pair only under \((H_{\mathcal W})\) or under
separate analytic proofs of the two displayed identities for those exact strings;
without such hypotheses this theorem remains the replacement-triple statement and
has no printed-string specialization.  Equivalently, the load-bearing input is
the finite verified triple \((\mathsf P,\mathsf S,\Abs)\); the pair notation records
only that this paper constructs the third, punctured absolute-value row
internally.
\end{theorem}

\begin{proof}
By \Cref{thm:eml-punctured-richardson-core},
\(\Abs(u)=|u|\) for every nonzero real \(u\).  Thus the three finite terms
\((\mathsf P,\mathsf S,\Abs)\) satisfy exactly the three semantic rows
\[
  \mathsf P=\pi,
  \qquad \mathsf S(x)=\sin x\quad(x\in\RR),
  \qquad \Abs(x)=|x|\quad(x\ne0)
\]
required in \Cref{thm:replacement-triple-richardson-interface}.  Applying that
theorem gives the effective translation and the displayed equivalence on the
recursive punctured domains \(D_E\).  Richardson's punctured identity-to-zero
problem is undecidable by \Cref{thm:richardson-punctured-source-problem}; hence a
zero-identity algorithm for the translated source-indexed fragment would decide
that source problem.  A computable finite semantic normalizer, or an
equality-reflecting compiler into finite decidable protocol objects, would decide
the same source problem by normalizing or compiling the two translated
source-restricted terms in the displayed equivalence and comparing the resulting
finite objects; this use is exactly the two-sided reflection/preservation
hypothesis stated above, not a merely syntactic normalization procedure.  The
final printed-pair sentence is exactly
\Cref{prop:printed-pi-sine-instantiation-criterion}: the two printed rows supply
\(\mathsf P\) and \(\mathsf S\) only when their real-domain identities are proved,
which in this manuscript is the optional hypothesis \((H_{\mathcal W})\).
\end{proof}

\begin{corollary}[Printed Richardson specialization is conditional]
\label{cor:printed-richardson-specialization-conditional}
The printed strings \((U_{47},U_{111})\) are not part of the unconditional
Richardson obstruction proved in this article.  They instantiate
\Cref{thm:verified-replacement-richardson-instantiation} exactly after the two
additional semantic identities
\[
  U_{47}=\pi,
  \qquad U_{111}(x)=\sin x\quad(x\in\RR)
\]
have been assumed under \((H_{\mathcal W})\) or proved separately.  Without those
identities the reader-facing theorem
is the verified-triple statement, with arbitrary finite proved witnesses
\((\mathsf P,\mathsf S)\) and the internally proved punctured row \(\Abs\) forming
the triple \((\mathsf P,\mathsf S,\Abs)\).
\end{corollary}

\begin{proof}
The theorem just proved uses the printed names only through the two displayed
identities.  By \Cref{lem:eml-pi-sine-witnesses}, the catalogue gives finite pure
EML syntax for those names, but without \((H_{\mathcal W})\) it does not prove
the two semantic identities.  Conversely, once the two identities are supplied,
the pair \((U_{47},U_{111})\) satisfies the hypotheses of
\Cref{thm:verified-replacement-richardson-instantiation}.  The absolute-value row
is independent of the printed pair and is supplied internally by
\Cref{thm:eml-punctured-richardson-core}.  Thus the equivalence between the
printed specialization and the extra semantic identities is exact.
\end{proof}

\begin{theorem}[Conditional EML-Richardson witness package]
\label{thm:no-core-only-richardson-translation}
\label{prop:printed-witness-certificate-boundary}
\label{thm:printed-richardson-witnesses-interface}
\label{thm:exact-richardson-witness-conditionality}
\label{thm:printed-string-real-domain-upgrade-criterion}
\label{thm:finite-eml-richardson-trichotomy}
\label{thm:no-unconditional-printed-root-richardson-instance}
\label{thm:separation-is-interface-only}
The principal-branch EML terms displayed in
\Cref{lem:eml-core-witness-table} form a Richardson witness package conditional
on the two source-locked printed-string assertions
\[
  W_{\pi}=U_{47}=\pi,
  \qquad
  W_{\sin}=U_{111}\text{ and }W_{\sin}(x)=\sin x\quad(x\in\RR).
\]
These two assertions are exactly the optional \((H_{\mathcal W})\) clauses used
for a printed-name specialization, unless they are proved separately for the
exact printed strings; neither the printed names nor their parser certificate is
an implicit proof of the displayed identities.  All other rows needed
for rational constants, \(\log2\), arithmetic, \(\exp\), and the punctured
absolute-value witness are proved inside the article.  The endpoint-root string
\(U_{126}\) and the endpoint identity for \(\sqrt t\) are not used in this
punctured Richardson package.  The multiple labels on this theorem are
dependency-cut aliases for these conditionality and boundary clauses; the single
load-bearing assertion is the displayed conditional witness package.
\end{theorem}

\begin{proof}
Every nontrigonometric, non-\(\pi\) row in \Cref{lem:eml-core-witness-table} has
already been proved on its stated real domain.  The two remaining rows are
precisely the two semantic identities named in the statement; they are not
deduced from the finite printed syntax, from the \(\mathcal B_{36}\) import, or
from the endpoint-root row.  Richardson's
Theorem~2, as recorded in \Cref{lem:richardson-fragment}, requires rational
constants, \(\log2\), \(\pi\), arithmetic, \(\exp\), sine, and a unary symbol
agreeing with absolute value off zero.  Under the two printed-string identities,
the table supplies the \(\pi\) and sine rows, while
\Cref{thm:eml-punctured-richardson-core} supplies the absolute-value row without
asserting a square-root value at zero.  No step invokes \(U_{126}\) or an
endpoint-square-root identity.
\end{proof}

\begin{lemma}[Zero-domain-free EML absolute-value witness]
\label{lem:eml-zero-absolute-witness}
For every real-valued translated subexpression \(F\), the term
\[\Abs(F)=G(M(F,F))\]
is defined and equal to \(|F|\) at every real input where \(F\) is defined and
\(F\ne0\).  This is the exact punctured absolute-value clause in Richardson's
condition (2).
\end{lemma}

\begin{proof}
At such an input, \Cref{lem:eml-real-witness-library} gives
\(M(F,F)=F^2>0\).  The positive-domain identity for \(G\) then gives
\(G(F^2)=\sqrt{F^2}=|F|\).  No claim is made at zeros of \(F\), and no such claim
is needed by \Cref{lem:richardson-punctured-mu-reduction}.
\end{proof}

\begin{corollary}[Finite Richardson witness library]
\label{cor:finite-richardson-witness-library}
The witness table of \Cref{lem:eml-core-witness-table} is finite and effective:
from a Richardson expression one can compute the corresponding pure EML term by
structural recursion, replacing each primitive by the displayed witness.
\end{corollary}

\begin{proof}
There are finitely many primitive clauses in the table, and each clause is a
finite macro over \((\eml,1)\).  Expanding a macro at each node of a finite
Richardson expression is a terminating structural recursion, so the translation
is computable.
\end{proof}

\begin{theorem}[Replacement-triple Richardson fragment]
\label{thm:largest-verified-richardson-fragment}
\label{thm:replacement-triple-richardson-interface}
Let \(\mathcal R_{\mathrm{core}}\) be Richardson's language with rational
constants, \(\log2\), \(\pi\), \(+,-,\times\), composition, \(\exp\), \(\sin\), and
a unary symbol \(\mu\) satisfying \(\mu(x)=|x|\) for \(x\ne0\).  Let
\((\mathsf P,\mathsf S,\mathsf A)\) be finite pure EML terms satisfying
\[
  \mathsf P=\pi,
  \qquad \mathsf S(x)=\sin x\quad(x\in\RR),
  \qquad \mathsf A(x)=|x|\quad(x\ne0).
\]
Then the principal-branch EML core gives an effective translation of
\(\mathcal R_{\mathrm{core}}\) into pure EML terms, using
\((\mathsf P,\mathsf S,\mathsf A)\) for \((\pi,\sin,\mu)\), that preserves the
Richardson real-function identities on the punctured domains required by the
theorem.  The identity-to-zero problem for the source fragment is undecidable,
and therefore the transported source-indexed EML equality problem
\[
  E\longmapsto
  \bigl(\tau_{\mathsf P,\mathsf S,\mathsf A}(E),D_E,
  \tau_{\mathsf P,\mathsf S,\mathsf A}(0)\bigr)
\]
is undecidable.  This is the unconditional submission theorem: the symbols
\(\mathsf P\), \(\mathsf S\), and \(\mathsf A\) are replacement witnesses whose
three semantic rows are part of the hypothesis.  No conclusion of this theorem
names the printed strings \((U_{47},U_{111})\) unless those exact strings have
first been proved, or explicitly assumed under \((H_{\mathcal W})\), to satisfy
\(U_{47}=\pi\) and \(U_{111}(x)=\sin x\) on \(\RR\).
\end{theorem}

\begin{proof}
The source language is precisely the class covered by Richardson's conditions
(1) and (2), so undecidability of identity to zero in the source is
\Cref{lem:richardson-fragment}.  The finite EML translation exists by the same
structural recursion as \Cref{cor:finite-richardson-witness-library}: use the
proved core rows for rational constants, \(\log2\), arithmetic, and \(\exp\), and
use the three displayed witnesses for \(\pi\), \(\sin\), and \(\mu\).  An induction
on Richardson expressions proves value preservation wherever the source
expression is defined and every encountered \(\mu\)-argument is nonzero.  This is
exactly the punctured absolute-value convention in Richardson's theorem.  Thus a
zero-identity algorithm for the transported source-indexed translated EML
fragment would decide the undecidable source identity-to-zero problem.  The
proof has used only the three displayed row identities and the internally proved
core rows.  Therefore replacing \((\mathsf P,\mathsf S,\mathsf A)\) by any other
finite triple with the same three verified rows gives the same reduction, while
replacing it by a merely printed syntax triple without those rows gives no
value-preservation induction and hence no instance of the theorem.  This proves
the stated replacement-witness boundary for \((U_{47},U_{111})\).
\end{proof}

\begin{corollary}[Unconditional Richardson replacement interface]
\label{cor:unconditional-richardson-replacement-interface}
\label{cor:printed-richardson-names-not-load-bearing}
After removing the optional hypothesis \((H_{\mathcal W})\), every Richardson
obstruction used in
\Cref{thm:obstruction-spine,thm:proc-ams-core-obstruction-spine,thm:main-hierarchy}
factors through a finite verified replacement triple
\((\mathsf P,\mathsf S,\mathsf A)\) satisfying the three rows in
\Cref{thm:replacement-triple-richardson-interface}.  With the internally proved
choice \(\mathsf A=\Abs\), this is equivalently a theorem relative to any finite
proved \(\pi\)- and sine-witness pair \((\mathsf P,\mathsf S)\).  The exact printed
names \((U_{47},U_{111})\) are not load-bearing in this unconditional interface;
they are admissible replacements only after the two real-domain identities for
those exact strings are supplied as hypotheses or proved separately.
\end{corollary}

\begin{proof}
The theorem just proved is stated for arbitrary finite pure EML terms
\((\mathsf P,\mathsf S,\mathsf A)\) satisfying the three semantic rows, and its
proof uses no source-locking information beyond those rows.  Taking
\(\mathsf A=\Abs\) is allowed by \Cref{thm:eml-punctured-richardson-core}, so any
proved pair \((\mathsf P,\mathsf S)\) for \(\pi\) and sine supplies a verified
triple.  Conversely, if \((U_{47},U_{111})\) is inserted without the two displayed
semantic identities, the hypotheses of
\Cref{thm:replacement-triple-richardson-interface} are not met.  Thus deleting
\((H_{\mathcal W})\) deletes exactly the named printed-pair specialization and
leaves the replacement-triple theorem, which is the Richardson clause cited by
the main hierarchy.
\end{proof}

\begin{theorem}[Relative finite Richardson fragment]
\label{thm:relative-finite-richardson-fragment}
The Richardson normalizer obstruction is uniform over finite pure EML triples
\((\mathsf P,\mathsf S,\mathsf A)\) with proved identities
\[
  \mathsf P=\pi,
  \qquad \mathsf S(x)=\sin x\quad(x\in\RR),
  \qquad \mathsf A(x)=|x|\quad(x\ne0).
\]
The internally proved term \(\Abs(u)=G(M(u,u))\) supplies the third row for every
such pair \((\mathsf P,\mathsf S)\).  The printed pair \((U_{47},U_{111})\) is not
used unless \((H_{\mathcal W})\) is explicitly assumed or the same \(\pi\) and
sine identities are independently proved for those exact strings, and
\(U_{126}\) is not a witness row in the punctured Richardson interface.
\end{theorem}

\begin{proof}
Apply \Cref{thm:largest-verified-richardson-fragment}.  Its proof uses only the
three listed identities in the clauses for \(\pi\), sine, and the punctured
absolute value; the remaining clauses are the internally proved EML core.  The
specific internal absolute-value term satisfies the third identity by
\Cref{thm:eml-punctured-richardson-core}.  If the printed pair is proposed as a
replacement for \((\mathsf P,\mathsf S)\), the same \(\pi\) and sine identities
must be proved for those two exact strings; provenance alone supplies no
semantic equality.  The absolute-value row is the internally proved punctured
term \(\Abs\), not an endpoint-root row.
\end{proof}

\begin{lemma}[Structural-recursion Richardson translation with domain clauses]
\label{lem:richardson-translation}
Fix finite pure EML terms \((\mathsf P,\mathsf S,\mathsf A)\) satisfying
\[
  \mathsf P=\pi,
  \qquad \mathsf S(x)=\sin x\quad(x\in\RR),
  \qquad \mathsf A(x)=|x|\quad(x\ne0).
\]
There is a computable translation \(\tau_{\mathsf P,\mathsf S,\mathsf A}\) from
\(\mathcal R_{\mathrm{core}}\) to pure EML terms such that every translated
primitive has the stated real-domain identity.  Consequently, for each
Richardson expression \(E\), the translated term represents the same real partial
function as \(E\) on the punctured Richardson domain relevant to the identity
problem.  The unary \(\mu\)-clause is transported only through the punctured row
\(\mathsf A(u)=|u|\) for \(u\ne0\); no endpoint square-root term, including the
printed row \(U_{126}\), is a generator in this translation.
\end{lemma}

\begin{proof}
Define \(\tau_{\mathsf P,\mathsf S,\mathsf A}\) recursively.  Rational constants
and \(\log2\) are sent to their finite arithmetic terms; the variable is sent to
the EML input; \(+,-,\times\), composition, and \(\exp\) are sent to the
corresponding internally proved rows of \Cref{lem:eml-core-witness-table};
\(\pi\) is sent to \(\mathsf P\); \(\sin\) is sent to \(\mathsf S\); and
\(\mu(F)\) is sent to \(\mathsf A(\tau_{\mathsf P,\mathsf S,\mathsf A}(F))\).  The
recursion is computable because the source expression tree and the three
replacement terms are finite.  The value-preservation assertion follows by
induction on the source expression.  The only non-total clause is \(\mu\), where
the induction uses the punctured identity \(\mathsf A(u)=|u|\) for \(u\ne0\),
exactly matching Richardson's hypothesis on \(\mu\).  Since the source language
has no endpoint square-root primitive and the target image of \(\mu\) is the
chosen punctured row \(\mathsf A\), no value of any root term at \(0\) is read by
the recursion.
\end{proof}


% End original section: sec05_hierarchy_strictness.tex
% Original section: sec05_richardson_interfaces.tex
\begin{theorem}[Punctured-domain many-one reduction]
\label{thm:punctured-domain-richardson-reduction}
\label{thm:richardson-punctured-domain-convention}
Fix finite pure EML terms \((\mathsf P,\mathsf S,\mathsf A)\) satisfying
\[
  \mathsf P=\pi,
  \qquad \mathsf S(x)=\sin x\quad(x\in\RR),
  \qquad \mathsf A(x)=|x|\quad(x\ne0).
\]
For each Richardson expression \(E\), define its punctured Richardson domain
\(D_E\subseteq\RR\) recursively as the set of real inputs for which every
ordinary elementary subterm is defined and, whenever a subterm \(\mu(F)\) is
evaluated, the value of \(F\) is nonzero.  Equivalently, the transported EML
instance is the source-indexed partial function obtained by restricting the
translated term to this same recursively defined set \(D_E\), not by enlarging it
to any principal-branch domain on which the translated syntax happens to have a
value.  Then
\(\tau_{\mathsf P,\mathsf S,\mathsf A}(E)\) and
\(\tau_{\mathsf P,\mathsf S,\mathsf A}(0)\) are defined on \(D_E\), the first term
equals \(E\) there, and the second term equals \(0\) there.  In particular,
\[
  E=0\hbox{ on }D_E
  \quad\Longleftrightarrow\quad
  \tau_{\mathsf P,\mathsf S,\mathsf A}(E)=
  \tau_{\mathsf P,\mathsf S,\mathsf A}(0)
  \hbox{ on }D_E .
\]
Zeros of absolute-value arguments therefore delete the same input points from the
source-indexed transported equality problem on both sides and do not invalidate
the many-one reduction.  The endpoint row \(U_{126}\) is not used in this
punctured obstruction: replacing \(\mathsf A\) by the internally proved
absolute-value row \(\Abs\), or by any separately verified punctured
absolute-value row, is the whole root interface required here.
\end{theorem}

\begin{proof}
Proceed by structural induction on \(E\), proving simultaneously that the source
subterm and its translated EML term agree on the recursively transported domain
and that the transported domain is the one assigned by the same recursion.
Constants, the variable, arithmetic, composition, and \(\exp\) use the internally
proved core identities and their ordinary real domains.  The \(\pi\) and \(\sin\)
clauses use the first two displayed hypotheses.  For a subterm \(\mu(F)\), the
recursive definition of \(D_{\mu(F)}\) requires \(x\in D_F\) and \(F(x)\ne0\).  By
the induction
hypothesis,
\(\tau_{\mathsf P,\mathsf S,\mathsf A}(F)(x)=F(x)\), so the translated argument to
\(\mathsf A\) is nonzero exactly on the same points.  The third displayed
hypothesis gives
\[
  \mathsf A(\tau_{\mathsf P,\mathsf S,\mathsf A}(F)(x))
  =|F(x)|=
  \mu(F)(x)
\]
for every \(x\in D_{\mu(F)}\).  If \(F(x)=0\), the recursive convention removes
\(x\) from the punctured domain before either the Richardson \(\mu\)-clause or the
EML \(\mathsf A\)-clause is used; no value at that point is needed, and no
endpoint-root semantics can repair or change the deleted point.  Thus each
constructor preserves both the relevant domain and the value on that domain.
Taking \(E=0\) gives the displayed equivalence, which is an effective many-one
reduction because the translation is computable.
\end{proof}
\begin{theorem}[Finite-budget Richardson translation]
\label{thm:richardson-translation-budget}
\label{lem:richardson-translation-bookkeeping}
Fix finite pure EML replacement terms \((\mathsf P,\mathsf S,\mathsf A)\) and a
finite EML witness table for the arithmetic, exponential, rational-constant,
\(\log2\), \(\pi\), sine, and punctured absolute-value rows used in
\Cref{lem:richardson-translation}.  Let \(m_R\) be the number of pure EML nodes
in the replacement macro for a Richardson constructor \(R\), counted after
expanding all internal arithmetic macros and counting each argument placeholder
once.  Let \(d_R\) be the number of argument-placeholder occurrences in that
macro, and set \(B=\max_R m_R\) and \(d=\max_R d_R\).
Then \(B\) and \(d\) are finite and effectively computable from the printed
witness table.  If a Richardson expression \(E\) has \(n\) source-constructor
nodes and no source subexpression is shared, then the fully expanded pure EML
syntax of \(\tau_{\mathsf P,\mathsf S,\mathsf A}(E)\) has at most
\[
  B(1+d+d^2+\cdots+d^n)\le B(n+1)\max(1,d)^n
\]
pure EML nodes.  Consequently the translation is a uniform finite-substitution
procedure: every output term is obtained by at most this many explicit node
insertions, the bound depends only on the fixed witness table and \(n\), and the
punctured-domain equality reduction of
\Cref{thm:punctured-domain-richardson-reduction} is a genuine finite many-one
reduction rather than an appeal to an unexpanded compiler.
\end{theorem}

\begin{proof}
There are only finitely many source constructors and, by hypothesis, each row in
the witness table is a printed finite pure EML macro after expanding the internal
arithmetic abbreviations.  Hence the two maxima defining \(B\) and \(d\) are
finite; they are computable by parsing the finite strings and counting pure EML
nodes and placeholder occurrences.

We prove the node bound by induction on the source tree.  Let \(N(F)\) be the
number of pure EML nodes in the fully expanded translation of a source
subexpression \(F\), and let \(h(F)\) be the number of source-constructor nodes in
\(F\).  If \(F\) is obtained by applying a constructor \(R\) to immediate
subexpressions \(F_1,\ldots,F_r\), then full expansion substitutes a separate
copy of \(\tau(F_i)\) into each placeholder occurrence assigned to the \(i\)-th
argument.  If \(q_{R,i}\) denotes the number of such occurrences, then
\(\sum_i q_{R,i}\le d\), and therefore
\[
  N(F)\le B+\sum_i q_{R,i}N(F_i)
       \le B+d\max_i N(F_i)
\]
with the evident convention for nullary constructors.  Since each proper
subexpression has height at most \(h(F)-1\), iteration along any root-to-leaf path
gives
\[
  N(F)\le B(1+d+d^2+\cdots+d^{h(F)}) .
\]
Because the height of a tree with \(n\) constructor nodes is at most \(n\), this
implies the displayed estimate for \(E\).  The recursive construction also gives
an explicit finite substitution algorithm: visit the source tree from the
leaves upward, replace the constructor at each node by its fixed macro, and copy
the already constructed child strings into the finitely many placeholders.  The
algorithm performs no semantic search and uses no external catalogue beyond the
fixed finite witness table.  Combining this finite construction with the
domain-preservation theorem \Cref{thm:punctured-domain-richardson-reduction}
therefore yields the stated finite many-one reduction.
\end{proof}

\begin{theorem}[Witness-row rigidity and defect classification]
\label{thm:richardson-witness-row-rigidity}
\label{thm:richardson-row-necessity-and-sufficiency}
Let \((\mathsf P,\mathsf S,\mathsf A)\) be finite pure EML terms.  The structural
translation \(\tau_{\mathsf P,\mathsf S,\mathsf A}\), with \(\mathsf P\),
\(\mathsf S\), and \(\mathsf A\) substituted respectively for Richardson's
\(\pi\), \(\sin\), and punctured absolute-value primitive, gives a
value-preserving many-one reduction from Richardson's punctured
identity-to-zero problem to a source-indexed pure EML equality problem if and
only if the three witness rows
\[
  \mathsf P=\pi,
  \qquad \mathsf S(x)=\sin x\quad(x\in\RR),
  \qquad \mathsf A(x)=|x|\quad(x\ne0)
\]
hold in the principal-branch EML semantics.  Consequently the exact printed
pair \((U_{47},U_{111})\) is not merely missing a parser certificate: the two
real-domain identities for those exact strings, either as the explicit
\((H_{\mathcal W})\) hypothesis or as separately proved facts, are logically
necessary and sufficient for the printed-pair specialization of the Richardson
obstruction, once the internally proved punctured row \(\mathsf A=\Abs\) is used.

Equivalently, the only possible failures of a finite source-locked Richardson
package are the three row defects
\[
  \mathsf P\ne\pi,
  \qquad \mathsf S\not\equiv\sin\hbox{ on }\RR,
  \qquad \mathsf A\not\equiv |-|\hbox{ on }\RR\setminus\{0\}.
\]
Here a row defect means that the displayed term is either undefined at some
required real input or defined there with the wrong value.  If any one of these
defects occurs, the structural translation already fails to be value-preserving
on the corresponding primitive test expression
  \(\pi\), \(\sin(x)\), or \(\mu(x)\).  Equivalently, the bad examples for this
  interface form the three classifiable row-defect skeletons: a wrong constant
  row, a wrong unary sine row, or a wrong punctured absolute-value row.  Hence no
  equality algorithm, semantic
normalizer, or protocol compiler obstruction in this paper may be attributed to
that defective finite triple.  If none occurs, \Cref{thm:separation} gives the
full punctured Richardson obstruction for that triple.
\end{theorem}

\begin{proof}
If the three displayed rows hold, \Cref{thm:punctured-domain-richardson-reduction}
gives value preservation on every punctured Richardson domain and hence the
many-one reduction to the transported EML equality problem.

Conversely suppose the indicated structural translation is value-preserving for
Richardson expressions on their punctured domains.  Apply it first to the
constant expression \(\pi\).  Its domain is all of \(\RR\), and the translated
term is \(\mathsf P\); value preservation gives \(\mathsf P=\pi\).  Apply it next
to the expression \(\sin(x)\).  Its domain is again all of \(\RR\), and the
translated term is \(\mathsf S(x)\); value preservation gives
\(\mathsf S(x)=\sin x\) for every real \(x\).  Finally fix an arbitrary nonzero
  real number \(a\) and apply the translation to the one-variable Richardson
  expression \(\mu(x)\) at the input value \(x=a\).  The punctured domain of
  \(\mu(x)\) is \(\RR\setminus\{0\}\), so \(a\) is an allowed input, the translated value is
\(\mathsf A(a)\), and Richardson's primitive has value \(|a|\).  Thus
\(\mathsf A(a)=|a|\) for every \(a\ne0\).  These three test expressions are
primitive clauses of the same structural recursion used in every larger
Richardson expression, so no weaker row data can supply the stated
value-preserving reduction.

Taking \(\mathsf A=\Abs\), the third row is supplied by
\Cref{thm:eml-punctured-richardson-core}.  Therefore the printed-pair
specialization with \((U_{47},U_{111})\) exists exactly when the first two tests
above hold for those printed terms; finite syntax alone does not enter the
argument except to make the translated EML terms finite.

The defect classification is the contrapositive of the three primitive tests
just used, together with the sufficiency direction.  If the \(\pi\)-row is
defective, value preservation fails on the constant Richardson expression
\(\pi\).  If the sine row is defective, choose a real input where
\(\mathsf S\) is undefined or has value different from \(\sin\); value
preservation fails for the primitive expression \(\sin(x)\) at that input.  If
the absolute-value row is defective, choose a nonzero real input where
\(\mathsf A\) is undefined or has value different from \(|-|\); value
preservation fails for the primitive punctured expression \(\mu(x)\) there.  In
each defective case the many-one reduction is not the Richardson reduction
proved in this manuscript, so the downstream zero-test, normalizer, and compiler
obstructions cannot be invoked for that finite triple.  In the no-defect case
the three displayed identities hold, and \Cref{thm:separation} supplies exactly
those downstream obstructions.
\end{proof}

\begin{theorem}[Primary verified replacement-triple Richardson obstruction]
\label{thm:proc-ams-relative-richardson-theorem}
\label{thm:relative-richardson-main-spine-closure}
\label{thm:interface-only-richardson-spine}
\label{thm:printed-root-richardson-conditional-clause}
\label{thm:no-unconditional-printed-richardson-edge}
\label{thm:separation}
\label{thm:no-principal-branch-eml-nonnegative-root}
Fix the principal-branch pure EML semantics of \Cref{def:eml}.  The primary,
unconditional Richardson theorem used by this paper is the following
replacement-triple scheme: for any finite pure EML triple
\((\mathsf P,\mathsf S,\mathsf A)\) satisfying
\[
  \mathsf P=\pi,
  \qquad \mathsf S(x)=\sin x\quad(x\in\RR),
  \qquad \mathsf A(x)=|x|\quad(x\ne0).
\]
Richardson's punctured identity-to-zero problem many-one reduces to the
source-indexed partial equality problem for the pure EML fragment translated
using this verified triple; the source domain \(D_E\) is part of the instance and
equality is equality of the transported partial real functions on that domain.
Consequently that transported fragment admits no zero-identity algorithm, no
computable semantic normalizer into finite objects with decidable equality that
reflects and preserves equality on the transported source domains, and no
equality-reflecting semantic compiler into finite decidable protocol objects.

The theorem has an unconditional scheme instance for every triple once the three
displayed real-domain identities for that triple are available as proved
hypotheses.  In particular the internally proved punctured row \(\Abs(u)=|u|\) from
\Cref{thm:eml-punctured-richardson-core} may serve as \(\mathsf A\).  The exact
printed pair \((U_{47},U_{111})\) gives the printed-string specialization of this
theorem only after the two real-domain identities
\(U_{47}=\pi\) and \(U_{111}(x)=\sin x\) for \(x\in\RR\) have been supplied in
the principal-branch semantics.  Equivalently, a citation to this theorem with
the printed pair is shorthand for the optional hypothesis \((H_{\mathcal W})\)
or for independently verified replacements of exactly those two identities.  The
finite syntax catalogue by itself is not one of the verified triples covered by
the theorem.  Thus every submission-facing use of the printed pair is a
conditional branch, whereas every submission-facing unconditional Richardson
conclusion factors through the abstract verified replacement triple above.
Without \((H_{\mathcal W})\), or without separate analytic proofs of those exact
two printed rows, the strings \(U_{47}\) and \(U_{111}\) do not occur in the
unconditional conclusion of this theorem.

Moreover, no finite pure EML term defined on a right-neighbourhood of \(0\) can
represent \(t\mapsto\sqrt t\) on \(\RR_{\ge0}\).  This endpoint obstruction does
not affect the preceding Richardson reduction, because that reduction uses only
the punctured absolute-value identity \(\mathsf A(u)=|u|\) for \(u\ne0\).  The
multiple labels on this theorem are dependency-cut aliases for the same verified
triple interface and its printed-string boundary.
\end{theorem}

\begin{proof}
By \Cref{lem:richardson-translation,thm:richardson-translation-budget}, the
translation \(\tau_{\mathsf P,\mathsf S,\mathsf A}\) is effective.  By
\Cref{thm:punctured-domain-richardson-reduction}, for every Richardson expression
\(E\),
\[
  E=0\hbox{ on }D_E
  \quad\Longleftrightarrow\quad
  \tau_{\mathsf P,\mathsf S,\mathsf A}(E)=
  \tau_{\mathsf P,\mathsf S,\mathsf A}(0)
  \hbox{ on }D_E .
\]
Since Richardson's punctured identity-to-zero problem for
\(\mathcal R_{\mathrm{core}}\) is undecidable by
\Cref{thm:richardson-punctured-source-problem}, this is a many-one reduction.  A
zero-identity algorithm on the source-indexed translated fragment would decide
the Richardson problem.  A computable semantic normalizer into finite objects
with decidable equality would decide the same problem only if it reflects and
preserves equality of the source-indexed partial real functions: normalize the
restrictions of \(\tau_{\mathsf P,\mathsf S,\mathsf A}(E)\) and
\(\tau_{\mathsf P,\mathsf S,\mathsf A}(0)\) to the same domain \(D_E\) and compare
the results.  This is the transported-normalizer obstruction of
\Cref{lem:transported-partial-normalizer-obstruction}.  Likewise, an
equality-reflecting semantic compiler into any finite decidable protocol object
would decide the problem by compiling the two translated terms and checking
target equality.

The specialization statement is the rigidity assertion of
\Cref{thm:richardson-witness-row-rigidity} applied to the printed pair.  The
preceding argument used only the three displayed identities for the chosen
triple \((\mathsf P,\mathsf S,\mathsf A)\), so any triple for which those
identities are hypotheses of the theorem is an instance of the interface.  The term
\(\Abs=G(M(-,-))\) satisfies the third identity by
\Cref{thm:eml-punctured-richardson-core}.  For the printed names
\(U_{47}\) and \(U_{111}\), however,
\Cref{prop:printed-pi-sine-instantiation-criterion,cor:no-unconditional-printed-richardson-triple}
prove exactly that the catalogue supplies finite syntax but not the two semantic
base identities for \(\pi\) and sine.  Thus those printed strings are not part of
the unconditional theorem scheme; they enter only under the corresponding
verified-row hypothesis, namely the \(U_{47},U_{111}\) part of
\((H_{\mathcal W})\), or under separate proofs of the same two identities.

It remains only to record why this does not smuggle in an endpoint square-root
term.  Let \(T(t)\) be a finite pure EML term defined for all
\(t\in[0,\epsilon)\).  After possibly shrinking \(\epsilon\), the real-domain
value of every subterm of \(T\) is real-analytic near \(0\): constants and the
input are analytic, and if a subterm is \(\eml(A(t),B(t))\), then definedness at
\(0\) gives \(B(0)\ne0\), so \(B\) stays in a single principal-logarithm branch
neighbourhood and \(e^{A(t)}-\operatorname{Log}B(t)\) is analytic.  Hence a term
equal to \(\sqrt t\) on \([0,\epsilon)\) would be analytic at \(0\).  Its Taylor
series would have the form \(a_0+a_1t+a_2t^2+\cdots\); equality with \(\sqrt t\)
for \(t>0\) forces \(a_0=0\), and division by \(t\) gives a bounded analytic
left side equal to \(t^{-1/2}\), a contradiction as \(t\downarrow0\).  Thus the
endpoint root is impossible, while the punctured term \(G(M(u,u))\) remains
valid for every nonzero real argument and is exactly the condition Richardson
uses.
\end{proof}

\begin{proposition}[Three-witness stability of the Richardson interface]
\label{prop:three-witness-stability}
\label{prop:three-witness-richardson-stability}
Let \((\mathsf P,\mathsf S,\mathsf A)\) be finite pure EML terms satisfying the
three real-domain identities
\[
  \mathsf P=\pi,
  \qquad \mathsf S(x)=\sin x\quad(x\in\RR),
  \qquad \mathsf A(x)=|x|\quad(x\ne0).
\]
Replacing the displayed witnesses \((W_{\pi},W_{\sin},\Abs)\) by
\((\mathsf P,\mathsf S,\mathsf A)\) gives an effective translated Richardson
fragment with exactly the obstruction asserted in \Cref{thm:separation}: no
zero-identity algorithm, no computable semantic normalizer into finite decidable
objects that reflects and preserves equality on the transported source domains,
and no equality-reflecting semantic compiler into finite decidable protocol
objects.  Hence the connection from \Cref{thm:separation} to
\Cref{cor:strict-L0-L1,thm:obstruction-spine,thm:main-hierarchy} depends only on
these three identities, not on the printed syntax of the chosen witnesses.  The
printed pair \((U_{47},U_{111})\) instantiates the first two rows only after those
two identities have been proved for the exact strings; the third row in the
submission theorem is the punctured term \(\Abs\), not \(U_{126}\).
\end{proposition}

\begin{proof}
Use the translation \(\tau_{\mathsf P,\mathsf S,\mathsf A}\) of
\Cref{lem:richardson-translation}, using
the unchanged clauses for rational constants, \(\log2\), the variable,
\(+,-,\times\), composition, and \(\exp\), and using \(\mathsf P\),
\(\mathsf S\), and \(\mathsf A\) in the clauses for \(\pi\), \(\sin\), and the
Richardson unary symbol \(\mu\).  The recursion is effective because the source
expression tree and the three replacement terms are finite.  An induction on
Richardson expressions gives value preservation on the real domains relevant to
Richardson's theorem: all unchanged clauses are those already verified in
\Cref{lem:eml-core-witness-table}; the three changed clauses are exactly the
three displayed hypotheses; and the only partial clause is the \(\mu\)-clause,
where Richardson requires agreement with \(|x|\) only for nonzero arguments.

Thus, for every expression \(E\) in \(\mathcal R_{\mathrm{core}}\),
\Cref{thm:punctured-domain-richardson-reduction} gives the equivalence of
identity to zero on the punctured Richardson domain and identity to the
translated zero term on the same domain.  Since the punctured identity-to-zero
problem for \(\mathcal R_{\mathrm{core}}\) is undecidable by
\Cref{thm:richardson-punctured-source-problem}, this is a many-one reduction to
the source-indexed replacement EML fragment.  A zero-identity algorithm for that
fragment would decide Richardson identity-to-zero; a computable semantic
normalizer into finite objects with decidable equality would decide it only under
the stated equality-reflection and preservation condition on the
source-restricted partial real functions, by normalizing the two source-restricted
translated terms and comparing the normal forms; and an equality-reflecting
semantic compiler into finite decidable protocol objects would decide it by
compiling the two source-restricted translated terms and comparing the target
objects.  These are precisely the
obstruction consequences used downstream in
\Cref{cor:strict-L0-L1,thm:obstruction-spine,thm:main-hierarchy}.  Finally, the
argument has used only the three displayed semantic identities.  Therefore a
specific printed pair, including \((U_{47},U_{111})\), supplies the first two rows
exactly when those identities have been proved for that pair; its name or
provenance supplies no additional semantic content.  The third row in the
submission theorem is supplied by the punctured term \(\Abs\), unless a different
finite punctured absolute-value witness is separately proved.
\end{proof}
The submission Richardson scope is therefore the following.  The Proc. AMS core
Richardson theorem stream is unconditional as a scheme over verified finite
punctured triples \((\mathsf P,\mathsf S,\mathsf A)\): once the three displayed
identities for \(\pi\), sine, and \(|x|\) off zero have been proved for a finite
triple, the obstruction follows.  The paper proves the internal absolute-value
option \(\mathsf A=\Abs\); the printed pair \((W_{\pi},W_{\sin})\) instantiates the
other two rows only under \((H_{\mathcal W})\).  It contains no unconditional
endpoint-root instance for \(U_{126}\).  This is an immediate extraction from
\Cref{thm:separation,prop:three-witness-stability}: the endpoint obstruction in
\Cref{thm:separation} shows why no pure EML term defined at \(0\) can be
substituted for the positive root term, so the Richardson input to
\Cref{thm:obstruction-spine,thm:main-hierarchy} is the punctured triple
interface, while the endpoint-root string \(U_{126}\) remains outside that input.

\begin{theorem}[Maximal verified-triple closure of the Richardson clause]
\label{thm:unconditional-submission-richardson-closure}
\label{thm:verified-triple-only-main-richardson-clause}
In the submission manuscript with no appeal to \((H_{\mathcal W})\), the
Richardson clause used by
\Cref{thm:obstruction-spine,thm:proc-ams-core-obstruction-spine,thm:main-hierarchy}
is exactly the following maximal verified-triple theorem and nothing stronger;
this is the form to which every submission-facing Richardson citation is to be
read:
for every finite pure EML triple \((\mathsf P,\mathsf S,\mathsf A)\) satisfying
\[
  \mathsf P=\pi,
  \qquad \mathsf S(x)=\sin x\quad(x\in\RR),
  \qquad \mathsf A(x)=|x|\quad(x\ne0),
\]
the source-indexed EML fragment obtained from Richardson's punctured language by
\(\tau_{\mathsf P,\mathsf S,\mathsf A}\) has no zero-identity algorithm, no
computable semantic normalizer into finite decidable objects that reflects and
preserves equality on the transported source domains, and no equality-reflecting
semantic compiler into finite decidable protocol objects.
Conversely, if a source-locked finite pure EML triple gives the same
value-preserving Richardson translation and hence the same three obstruction
consequences through the structural recursion \(\tau\), then the three displayed
rows must hold.  Thus the verified-triple condition is not merely a sufficient
hypothesis but the necessary and sufficient interface for this paper's
Richardson theorem chain.
The manuscript proves the third row unconditionally for
\(\mathsf A=\Abs=G(M(-,-))\), but it does not prove the first two rows for the
particular printed names \((U_{47},U_{111})\).  Hence the exact printed pair
\((U_{47},U_{111})\) is not an unconditional instance of the displayed theorem;
it becomes an instance precisely after those two printed real-domain identities
are assumed as \((H_{\mathcal W})\) or supplied as proved hypotheses.  More generally,
\Cref{thm:richardson-witness-row-rigidity} classifies every source-locked finite
triple by its row defects: a defective \(\pi\), sine, or punctured absolute-value
row fails already on a primitive Richardson test expression, while a defect-free
triple is exactly a verified triple covered by \Cref{thm:separation}.  The cited
\(\mathcal B_{36}\) representative existence statement and the finite syntax
catalogue do not alter this closure; by the interface rigidity of
\Cref{thm:submission-interface-rigidity}, deleting \((H_{\mathcal W})\) removes
exactly the named printed-pair branch and no replacement-triple consequence of
the main hierarchy.
\end{theorem}

\begin{proof}
The displayed verified-triple statement is exactly \Cref{thm:separation}.  Its
proof uses Richardson's punctured undecidability theorem through
\Cref{thm:richardson-punctured-source-problem}, the effective translation and
domain preservation of
\Cref{lem:richardson-translation,thm:punctured-domain-richardson-reduction}, and
the three displayed semantic identities.  The normalizer and compiler exclusions
are the finite-decidable-equality consequences proved in the same theorem and in
\Cref{prop:L1-compiles-L0}.  Thus every downstream citation from
\Cref{thm:obstruction-spine,thm:proc-ams-core-obstruction-spine,thm:main-hierarchy}
to the Richardson obstruction factors through this verified-triple interface.

The internal proof of the third row for \(\Abs\) is
\Cref{thm:eml-punctured-richardson-core}: for nonzero real \(u\),
\(M(u,u)=u^2>0\) and the positive-domain square-root witness gives
\(G(M(u,u))=\sqrt{u^2}=|u|\).  The first two rows for the printed names are
different.  By \Cref{lem:eml-pi-sine-witnesses} the catalogue proves only that
\(U_{47}\) and \(U_{111}\) are finite pure EML strings; by
\Cref{prop:printed-pi-sine-instantiation-criterion,cor:no-unconditional-printed-richardson-triple}
the identities \(U_{47}=\pi\) and \(U_{111}(x)=\sin x\) on \(\RR\) are exactly the
extra hypotheses needed to substitute those names for \((\mathsf P,\mathsf S)\).
Therefore omitting \((H_{\mathcal W})\) removes only the printed-name
specialization, not the verified-triple theorem itself.  The stronger row-defect
classification cited in the statement is \Cref{thm:richardson-witness-row-rigidity}:
if a source-locked finite triple is not a verified triple, one of its three
primitive rows fails and the structural Richardson reduction is unavailable
already at the corresponding primitive test expression.

Finally,
\Cref{thm:B36-import-exact-scope,thm:B36-self-contained-replacement-core,thm:submission-interface-rigidity}
show that the \(\mathcal B_{36}\) import is representative existence for the
external calculator basis and is not a term-by-term proof of the Richardson
witness rows.  The interface-rigidity theorem also identifies the exact effect
of removing \((H_{\mathcal W})\): it deletes the named printed-pair branch while
leaving the replacement-triple Richardson obstruction used by
\Cref{thm:main-hierarchy}.  Consequently neither that citation nor the syntax
catalogue can create an unconditional printed \((U_{47},U_{111})\) Richardson
instance, and the submission-facing Richardson clause is precisely the
verified-triple closure stated above.
\end{proof}

\begin{corollary}[Printed-string-free Richardson spine]
\label{cor:printed-string-free-richardson-spine}
If every occurrence of the optional hypothesis \((H_{\mathcal W})\) and every
unproved printed identification of \((U_{47},U_{111})\) with \((\pi,\sin)\) is
removed from the submission, the Richardson part of
\Cref{thm:obstruction-spine,thm:proc-ams-core-obstruction-spine,thm:main-hierarchy}
remains the following unconditional theorem scheme: for each finite pure EML
triple \((\mathsf P,\mathsf S,\mathsf A)\) satisfying
\[
  \mathsf P=\pi,
  \qquad \mathsf S(x)=\sin x\quad(x\in\RR),
  \qquad \mathsf A(x)=|x|\quad(x\ne0),
\]
the transported punctured Richardson fragment has no zero-identity algorithm, no
computable semantic normalizer into finite decidable objects reflecting and
preserving equality on source domains, and no equality-reflecting compiler into
finite decidable protocol objects.  The optional printed pair
\((U_{47},U_{111})\) adds only the specialization in which the first two displayed
rows are supplied by \((H_{\mathcal W})\); it is not an additional premise of the
core hierarchy theorem.
\end{corollary}

\begin{proof}
After deleting \((H_{\mathcal W})\), the proof of
\Cref{thm:unconditional-submission-richardson-closure} still gives exactly the
displayed verified-triple obstruction, because that proof uses only
\Cref{thm:richardson-punctured-source-problem,lem:richardson-translation,thm:punctured-domain-richardson-reduction}
and the three semantic rows for the chosen triple.  The internally proved row
\(\Abs=|-|\) on \(\RR\setminus\{0\}\) remains available by
\Cref{thm:eml-punctured-richardson-core}; the first two rows may be supplied by
any finite verified witnesses and do not depend on the printed names.
\Cref{thm:richardson-witness-row-rigidity} proves the converse boundary: if the
printed pair is used as those first two witnesses, then the identities
\(U_{47}=\pi\) and \(U_{111}(x)=\sin x\) on \(\RR\) are precisely the missing row
data.  Hence removing \((H_{\mathcal W})\) deletes only that named
specialization, and leaves no residual unconditional statement about the exact
printed roots \(U_{47}\) and \(U_{111}\).  The downstream obstruction spine and
main hierarchy cite the verified-triple Richardson interface through
\Cref{thm:unconditional-submission-richardson-closure}, so their Richardson input
survives unchanged as the displayed theorem scheme.
\end{proof}

\begin{remark}[Dependency cut]
The source-locking and \(\mathcal B_{36}\)-noninstantiation clauses are separate
from the Richardson obstruction: \Cref{thm:unconditional-submission-richardson-closure} uses
a punctured triple interface and does not import the endpoint-root string.
\end{remark}

\begin{corollary}[Triple-relative expression/protocol separation]
\label{cor:strict-L0-L1}
For every finite punctured EML Richardson triple
\((\mathsf P,\mathsf S,\mathsf A)\) satisfying the three identities of
\Cref{thm:separation}, the imported \(L0\) expression-universality result for
\((\eml,1)\) does not provide a computable semantic normalizer that reflects
and preserves source-domain equality for the translated Richardson fragment.
\end{corollary}

\begin{proof}
A computable semantic normalizer for the translated punctured fragment with the
stated equality-reflection and preservation property would contradict the
normalizer obstruction in \Cref{thm:separation}.  The only
absolute-value identity needed is the punctured row \(|x|\) off zero; no endpoint
square-root identity is used.
\end{proof}

\begin{proposition}[Protocol compilation boundary]
\label{prop:L1-compiles-L0}
Let \(\mathcal E\) be an effective expression grammar with semantic equivalence
relation \(\equiv\), interpreted for a source-indexed partial-real fragment as
equality on the specified common source domains.  Suppose an \(L1\) protocol has
computable normal forms with decidable equality.  If there is a computable
semantic compiler \(C\colon\mathcal E\to N\) into those normal forms such that
\[
  E\equiv E' \quad\Longleftrightarrow\quad C(E)=C(E')
\]
for all source expressions in the fragment under consideration, then semantic
equivalence in \(\mathcal E\) is decidable.  Consequently no Richardson
identity-to-zero source fragment, transported with its domains and retaining this
equality-reflection clause, admits such a compiler into the Zeckendorf \(L1\)
protocol.
\end{proposition}

\begin{proof}
Given two source expressions \(E,E'\), compute \(C(E)
\) and \(C(E')\).  Since equality of target normal forms is decidable in an
\(L1\) protocol, decide whether \(C(E)=C(E')\).  The displayed reflection and
preservation condition makes this decision equivalent to deciding
\(E\equiv E'\), including equality on the specified common domains in the
partial-real case.  If the source fragment is Richardson-undecidable in this
identity-to-zero sense, this would contradict \Cref{lem:richardson-fragment}
after the effective transported translation of \Cref{lem:richardson-translation}.  The Zeckendorf target has
decidable normal-form equality by \Cref{prop:fold-confluence} and is an \(L1\)
protocol by \Cref{thm:zeckendorf-protocol-univ}, so it is covered by the general
argument.
\end{proof}

\begin{corollary}[Compilation scope]
\label{cor:compile-scope}
The Zeckendorf \(L1\) protocol is a computable protocol for natural-number
arithmetic after value transport, but it is not a computable semantic compiler
or equality-reflecting normalizer for an \(L0\) expression grammar whose
source-indexed partial-real equivalence is Richardson-undecidable.
\end{corollary}

\begin{proof}
The positive protocol scope is \Cref{thm:zeckendorf-protocol-univ}.  The excluded
compiler scope is exactly \Cref{prop:L1-compiles-L0}: a compiler reflecting
source-domain semantic equivalence into decidable Zeckendorf normal-form equality
would decide the Richardson-undecidable source identity-to-zero relation.
\end{proof}

\begin{remark}[Transported Richardson boundary]
The Richardson obstruction concerns transported real semantics on the verified
Richardson fragment only.  It is not a decidability statement for arbitrary raw
EML terms on complex principal-branch domains.
Every equality used in \Cref{lem:richardson-translation} is an equality of
partial real functions on the domains of the source Richardson expressions.
The proof never constructs or compares arbitrary EML terms outside that
translated fragment.
\end{remark}

\begin{proposition}[External universality versus verified obstruction]
\label{prop:external-versus-locked-package}
\label{thm:no-hidden-L0-semantic-normalizer}
The Odrzywo{\l}ek \(L0\) import and the Richardson obstruction have different
logical roles: the former supplies expression representatives for a calculator
basis, while the latter requires the explicit three-witness interface of
\Cref{thm:no-core-only-richardson-translation}.
\end{proposition}

\begin{proof}
The \(L0\) import is \Cref{thm:eml-universality}; the obstruction is
\Cref{thm:separation}.  The proof of \Cref{thm:separation} cites only the core
EML witnesses, the three-witness interface, and Richardson's theorem.
\end{proof}


% End original section: sec05_richardson_interfaces.tex
% Original section: sec06_discussion.tex
\section{Final remarks}\label{sec:discussion}

The hierarchy separates representation, protocol, and certificate evidence.
Expression universality supplies representatives in the cited elementary
library.  Protocol universality adds computable Zeckendorf normal forms and
value transport through \(\NN\).  The Richardson interface prevents the
verified punctured fragment from admitting a computable semantic normalizer,
while the one-free-monogenic-orbit theorem identifies the exact cyclic-ray
boundary for a faithful ambient-product readout.  These are independent
obstructions to two different collapses of the hierarchy.

The \(L2\) conclusions now have a sharp intrinsic form.  In the terminology of
\Cref{def:fixed-finite-state-transfer,def:cert-univ}, fixed finite-state transfer
is degreewise: the matrix may depend on \(q\), but not on the resolution \(m\).
The unmodified Zeckendorf fold satisfies this condition, and
\Cref{thm:primitive-intrinsic-moment-transfer} strengthens it to
\[
  S_q^{\mathrm{int}}(m)
  =c_q\lambda_q^m+O(\theta_q^m),
  \qquad 0\le\theta_q<\lambda_q ,
\]
for every fixed \(q\ge1\).  Hence the fixed-degree pressure exists and equals
\(\log\lambda_q\).

There is nevertheless no rational-function transfer matrix in the moment marker
uniform over the entire moment tower.
By \Cref{thm:intrinsic-all-degree-transfer-obstruction}, the bivariate series
\[
  \sum_{m\ge1}\sum_{q\ge0}S_q^{\mathrm{int}}(m)z^mw^q
\]
is nonrational, and no single rational matrix law in the moment marker \(w\)
produces all \(G_m(w)=\sum_qS_q^{\mathrm{int}}(m)w^q\).  This obstruction is
intrinsic to the original \(\Fold_m\) fibres.  It is forced by the exact height
law: the diagonal sequence \(S_m^{\mathrm{int}}(m)\) has no fixed transfer
formula because it grows like
\(\exp((\frac12\log\varphi+o(1))m^2)\), beyond the coefficient budget of every
rational bivariate series.

The computable-cover theorem records a separate boundary.  Finite-fibre
certificates do not force even degreewise transfer, since a one-spike cover can
make its second moment superexponential while preserving the normal forms and
all \(L1\) readouts.  Thus the intrinsic theorem separates fixed-degree transfer
from a rational-function all-degree law in the moment marker, whereas the cover
theorem separates fixed-resolution certificates from fixed-degree transfer.  The
exact intrinsic height and pressure results are properties of the original fold
fibres.

The appendix consists of consequences of the finite-fibre objects and is not an
input to any hierarchy or separation theorem above.

% End original section: sec06_discussion.tex
% Original section: sec07_derived_consequences.tex
\section{Appendix: finite-fiber support}
\label{sec:consequences}

This appendix records three elementary support facts that are downstream
from the main constructions.  The first is a finite-mixture
decomposition for the finite fiber measure
\[
  \nu_m=\sum_{x\in X_m}\delta_{d_m(x)} .
\]
The second is a fold-numeral restatement of the natural numbers
object after \(\Val\) has identified
\(X_{\infty}^{\mathrm{fin}}\) with \(\NN\).  Between them we record the
fixed-resolution star-kernel compression associated with a finite fibre-size
multiset.  These facts are included only
as reader support; none of the proofs of the \(L0/L1/L2\) hierarchy or
of its separations uses the results in this section.

\begin{proposition}[Finite mixture decomposition]
\label{thm:mixture}
Fix \(m\ge1\) and \(q\ge1\).  Let
\[
  R_m^{(q)}=\{(\omega^{(1)},\ldots,\omega^{(q)})\in\Omega_m^q:
  \Fold_m(\omega^{(1)})=\cdots=\Fold_m(\omega^{(q)})\},
\]
and let \(U_{m,q}\) be the uniform probability measure on
\(R_m^{(q)}\).  For \(x\in X_m\), put
\[
  F_x=\Fold_m^{-1}(x),\qquad \nu_x=\mathrm{Unif}(F_x).
\]
Then
\begin{equation}\label{eq:mixture}
  U_{m,q}
  =\sum_{x\in X_m}\frac{d_m(x)^q}{S_q(m)}\,\nu_x^{\otimes q}.
\end{equation}
Equivalently, if \(X\) denotes the common fold image under
\(U_{m,q}\), then
\[
  \Pr[X=x]=\frac{d_m(x)^q}{S_q(m)},
\]
and, conditional on \(X=x\), the coordinates
\(\omega^{(1)},\ldots,\omega^{(q)}\) are independent and uniform on
\(F_x\).
\end{proposition}

\begin{proof}
The decomposition
\[
  R_m^{(q)}=\bigsqcup_{x\in X_m}F_x^q
\]
is disjoint, and \(|F_x^q|=d_m(x)^q\).  The right hand side of
\eqref{eq:mixture} assigns to each point of \(F_x^q\) the mass
\[
  \frac{d_m(x)^q}{S_q(m)}\,d_m(x)^{-q}=\frac1{S_q(m)}.
\]
Since \(|R_m^{(q)}|=S_q(m)\), this is the uniform law.
\end{proof}

The measures \(\nu_x^{\otimes q}\) have disjoint supports.  Hence the
convex hull
\[
  \mathcal C_{m,q}
  =\operatorname{conv}\{\nu_x^{\otimes q}:x\in X_m\}
\]
is a finite simplex whose extreme points are these product laws.  This
is the elementary probabilistic form of the finite-fiber certificate:
the collision moments weight the finite fibers, but they do not by
themselves assert any finite-state transfer matrix.

\begin{proposition}[Finite Hankel star kernel]
\label{thm:star-kernel}
For a finite fibre-size multiset \(D\), define
\(M_q(D)=\sum_{d\in D}d^q\) and \(K_D(i,j)=M_{i+j}(D)\) for \(i,j\ge0\).  If the
distinct values of \(D\) are \(a_1,\ldots,a_s\) with multiplicities
\(n_1,\ldots,n_s\), then every finite principal compression of \(K_D\) is the
Gram matrix of the vectors
\[
  i\longmapsto (\sqrt{n_1}a_1^i,\ldots,\sqrt{n_s}a_s^i).
\]
Equivalently,
\[
  K_D(i,j)=\sum_{\alpha=1}^s n_\alpha a_\alpha^i a_\alpha^j .
\]
Thus the Hankel kernel is a finite diagonal nonnegative compression of the
Vandermonde star feature map attached to the fixed fibre-size atoms, in standard
nonnegative-matrix terminology \cite{BermanPlemmons1994}.
\end{proposition}

\begin{proof}
The displayed identity follows immediately from
\(M_{i+j}(D)=\sum_\alpha n_\alpha a_\alpha^{i+j}
=\sum_\alpha n_\alpha a_\alpha^i a_\alpha^j\).  Therefore, for any finite index
set \(I\subset\NN\) and coefficients \((c_i)_{i\in I}\),
\[
  \sum_{i,j\in I}c_i c_jK_D(i,j)
  =\sum_{\alpha=1}^s n_\alpha\left(\sum_{i\in I}c_i a_\alpha^i\right)^2\ge0.
\]
This is exactly the Gram matrix of the stated feature vectors, with the positive
diagonal weights \(n_\alpha\) on the finite atom set.
\end{proof}

\begin{proposition}[Fold numerals as a natural numbers object]
\label{prop:nno}
Define
\[
  0_X=\Val^{-1}(0)\in X_{\infty}^{\mathrm{fin}},\qquad
  \mathsf S_X(c)=\Val^{-1}(\Val(c)+1).
\]
Then \((X_{\infty}^{\mathrm{fin}},0_X,\mathsf S_X)\) is
isomorphic in \(\mathbf{Set}\) to
\((\NN,0,\mathsf S)\), hence is a natural numbers object
\cite{Lawvere1964}.  Under this identification, the numeral
\[
  c_n=\Val^{-1}(n)
\]
acts on every endomorphism \(f\colon A\to A\) by \(n\)-fold
iteration \(f\mapsto f^{\circ n}\).  The level-zero readout
identifies addition with composition of iterates; the level-one
readout identifies multiplication with composition of the operators
\(f\mapsto f^{\circ n}\).
\end{proposition}

\begin{proof}
The valuation is a bijection by \Cref{prop:fold-confluence}, and it
carries \(0_X\) to \(0\) and \(\mathsf S_X\) to the usual successor
on \(\NN\).  This proves the natural-numbers-object statement in
\(\mathbf{Set}\).  The iteration identities are exactly
\Cref{thm:two-level} transported through this isomorphism.
\end{proof}

% End original section: sec07_derived_consequences.tex
